% Options for packages loaded elsewhere
\PassOptionsToPackage{unicode}{hyperref}
\PassOptionsToPackage{hyphens}{url}
\PassOptionsToPackage{dvipsnames,svgnames,x11names}{xcolor}
%
\documentclass[
  10pt,
]{article}
\usepackage{amsmath,amssymb}
\usepackage{iftex}
\ifPDFTeX
  \usepackage[T1]{fontenc}
  \usepackage[utf8]{inputenc}
  \usepackage{textcomp} % provide euro and other symbols
\else % if luatex or xetex
  \usepackage{unicode-math} % this also loads fontspec
  \defaultfontfeatures{Scale=MatchLowercase}
  \defaultfontfeatures[\rmfamily]{Ligatures=TeX,Scale=1}
\fi
\usepackage{lmodern}
\ifPDFTeX\else
  % xetex/luatex font selection
\fi
% Use upquote if available, for straight quotes in verbatim environments
\IfFileExists{upquote.sty}{\usepackage{upquote}}{}
\IfFileExists{microtype.sty}{% use microtype if available
  \usepackage[]{microtype}
  \UseMicrotypeSet[protrusion]{basicmath} % disable protrusion for tt fonts
}{}
\makeatletter
\@ifundefined{KOMAClassName}{% if non-KOMA class
  \IfFileExists{parskip.sty}{%
    \usepackage{parskip}
  }{% else
    \setlength{\parindent}{0pt}
    \setlength{\parskip}{6pt plus 2pt minus 1pt}}
}{% if KOMA class
  \KOMAoptions{parskip=half}}
\makeatother
\usepackage{xcolor}
\usepackage{longtable,booktabs,array}
\usepackage{calc} % for calculating minipage widths
% Correct order of tables after \paragraph or \subparagraph
\usepackage{etoolbox}
\makeatletter
\patchcmd\longtable{\par}{\if@noskipsec\mbox{}\fi\par}{}{}
\makeatother
% Allow footnotes in longtable head/foot
\IfFileExists{footnotehyper.sty}{\usepackage{footnotehyper}}{\usepackage{footnote}}
\makesavenoteenv{longtable}
\setlength{\emergencystretch}{3em} % prevent overfull lines
\providecommand{\tightlist}{%
  \setlength{\itemsep}{0pt}\setlength{\parskip}{0pt}}
\setcounter{secnumdepth}{5}

\usepackage{amsmath,amssymb}
\usepackage{graphicx}
\usepackage{longtable,booktabs,array}
\usepackage{xcolor}
\usepackage[margin=2.0cm]{geometry}
\usepackage{microtype}
\setlength{\emergencystretch}{3em}
\setcounter{MaxMatrixCols}{20}
\pdfinfoomitdate=1
\pdftrailerid{}
\pdfsuppressptexinfo=-1
\renewcommand{\arraystretch}{1.12}
\usepackage{fvextra}
\RecustomVerbatimEnvironment{verbatim}{Verbatim}{breaklines,breakanywhere,fontsize=\small}
\fvset{breaklines,breakanywhere}
\allowdisplaybreaks
\sloppy
\ifLuaTeX
  \usepackage{selnolig}  % disable illegal ligatures
\fi
\IfFileExists{bookmark.sty}{\usepackage{bookmark}}{\usepackage{hyperref}}
\IfFileExists{xurl.sty}{\usepackage{xurl}}{} % add URL line breaks if available
\urlstyle{same}
\hypersetup{
  pdftitle={The stabiliser of the bulk},
  pdfauthor={The stabiliser: found and named by Claude (Anthropic), 2026-10-24; This account and its checks: Claude (Anthropic), 2026-10-25; Underlying theory, metric and Lovelock weights: Patrick L. Nash; Repository once-ere/Dirac\_claude, folder universe\_model\_3},
  colorlinks=true,
  linkcolor={blue},
  filecolor={Maroon},
  citecolor={Blue},
  urlcolor={blue},
  pdfcreator={LaTeX via pandoc}}

\title{The stabiliser of the bulk}
\usepackage{etoolbox}
\makeatletter
\providecommand{\subtitle}[1]{% add subtitle to \maketitle
  \apptocmd{\@title}{\par {\large #1 \par}}{}{}
}
\makeatother
\subtitle{What the stabiliser of UM-3 is, why the field equations of the
author's metric require it, what it does and what it does not do, and
why UM-3 needs it: a complete account for a newcomer, with its
provenance}
\author{The stabiliser: found and named by Claude (Anthropic),
2026-10-24 \and This account and its checks: Claude (Anthropic),
2026-10-25 \and Underlying theory, metric and Lovelock weights: Patrick
L. Nash \and Repository once-ere/Dirac\_claude, folder
universe\_model\_3}
\date{}

\begin{document}
\maketitle

{
\hypersetup{linkcolor=}
\setcounter{tocdepth}{3}
\tableofcontents
}
\hypertarget{the-fold-definition-and-examples}{%
\section{The fold: definition and
examples}\label{the-fold-definition-and-examples}}

\textbf{Definition.} The \emph{fold} is the seven-dimensional surface of
the author's eight-dimensional spacetime on which the cosine of the
angle \(z = 6Hx_0\) of the hidden direction vanishes:

\[
\mathcal{F} = \big\{X = (x_0, x_1, \dots, x_7):\ \cos(6Hx_0) = 0\big\},\qquad
\text{in one period } 0 < z < \pi\text{ of the metric: } z = \frac{\pi}{2},\quad x_0 = \frac{\pi}{12H} \tag{FD-1}
\] \textbf{Symbols in (FD-1).}

\begin{itemize}
\tightlist
\item
  \(\mathcal{F}\): The fold: the seven-dimensional surface
  \(\cos(6Hx_0) = 0\) of the eight-dimensional spacetime (FD-1), in one
  period \(0 < z < \pi\) the surface \(z = \pi/2\), \(x_0 = \pi/(12H)\),
  where \(\sigma = 0\); \(\vert_{\mathcal{F}}\) means restricted to the
  fold. Status: DERIVED.
\item
  \(X = \{x_0, \dots, x_7\}\): The author's coordinates of the
  eight-dimensional spacetime, \(X = \{x_0, x_1, \dots, x_7\}\) with the
  hidden direction first (Dirac\_Einstein-Lovelock\_4+4\_Nash.nb,
  section DEFCHART); \(X^M\) is the coordinate with the index \(M\) and
  \(\mathrm{d}^8X = \mathrm{d}x_0\,\mathrm{d}x_1\cdots\mathrm{d}x_7\)
  the coordinate volume element. Unit: length (each coordinate). Status:
  AUTHOR.
\item
  \(x_0\): The hidden space-like coordinate (listed first); the metric
  depends on it through \(z = 6Hx_0\), periodically with the period
  \(\pi/(6H)\). Unit: length. Status: AUTHOR.
\item
  \(x_1\): The first coordinate of ordinary 3-space. Unit: length.
  Status: AUTHOR.
\item
  \(x_7\): The third of the three extra time-like coordinates. Unit:
  length. Status: AUTHOR.
\item
  \(H\): The author's constant (an inverse length) that fixes the period
  \(\pi/(6H)\) of \(x_0\), the Lovelock couplings
  \(\alpha_{(p)} = w_{(p)}H^{2-2p}\) and the cosmological constant
  \(\Lambda = -7H^2\) (the author's original value \(-10H^2\) was
  corrected on 2026-10-25); in UM-3 a free parameter. Unit: GeV
  (1/length). Status: AUTHOR.
\item
  \(z\): The angle variable of the hidden direction, \(z = 6Hx_0\); one
  period of \(x_0\) is \(0 < z < \pi\), the fold is \(z = \pi/2\). Unit:
  dimensionless. Status: AUTHOR.
\item
  \(\pi\): The number pi, the ratio of the circumference of a circle to
  its diameter. Unit: dimensionless. Value: 3.141592654. Status:
  STANDARD.
\end{itemize}

\(X = \{x_0, \dots, x_7\}\) are the author's coordinates: \(x_0\) the
hidden space-like direction, \(x_1, x_2, x_3\) ordinary space, \(x_4\)
the time, \(x_5, x_6, x_7\) the three extra times; \(H\) is the author's
constant (an inverse length). The author's metric depends on \(x_0\)
only through \(\cot^2z\) and \((\sin^2z)^{1/6}\), so it repeats with the
period \(\pi/(6H)\) of \(x_0\), and every period \(0 < z < \pi\)
contains exactly one surface of \(\mathcal{F}\), at \(z = \pi/2\):
\(\mathcal{F}\) is the union of these copies, one per period (example
2), and ``the fold'' without a qualifier is the copy in the period that
contains our world. A point of the fold is labelled by the seven
coordinates \(x_1, \dots, x_7\), and the metric restricted to the fold
(\(\mathrm{d}x_0 = 0\), \(z = \pi/2\)) is

\[
g_{MN}\,\mathrm{d}X^M\mathrm{d}X^N\big\vert_{\mathcal{F}} = \mathrm{e}^{2a_4}\big(\mathrm{d}x_1^2 + \mathrm{d}x_2^2 + \mathrm{d}x_3^2\big)
 - \mathrm{d}x_4^2 - \mathrm{e}^{-2a_4}b_4^2\big(\mathrm{d}x_5^2 + \mathrm{d}x_6^2 + \mathrm{d}x_7^2\big) \tag{FD-2}
\] \textbf{Symbols in (FD-2).}

\begin{itemize}
\tightlist
\item
  \(g_{MN}\): The metric tensor of the eight-dimensional spacetime
  (signature \(+,+,+,+,-,-,-,-\) in the order \(x_0, \dots, x_7\)):
  \(g_{MN}\,\mathrm{d}X^M\mathrm{d}X^N\) is the squared interval,
  \(g^{MN}\) its inverse, \(\lvert g\rvert\) the absolute value of its
  determinant (\(\sqrt{\lvert g\rvert}\,\mathrm{d}^8X\) is the proper
  volume element). Unit: dimensionless. Status: AUTHOR.
\item
  \(M\): A coordinate index of eight dimensions (also with a numbered
  copy \(M_1, M_2, \dots\)): it runs over \(x_0, x_1, \dots, x_7\); a
  repeated index is summed. Status: STANDARD.
\item
  \(N\): A coordinate index of eight dimensions (also
  \(N_1, N_2, \dots\)), running over \(x_0, \dots, x_7\); a repeated
  index is summed. Status: STANDARD.
\item
  \(\mathrm{d}\): The differential: \(\mathrm{d}X\) is an infinitesimal
  change of \(X\), \(\mathrm{d}X/\mathrm{d}Y\) the derivative of \(X\)
  with respect to \(Y\), \(\mathrm{d}^nX\) the volume element of \(n\)
  variables. Status: STANDARD.
\item
  \(X = \{x_0, \dots, x_7\}\): The author's coordinates of the
  eight-dimensional spacetime, \(X = \{x_0, x_1, \dots, x_7\}\) with the
  hidden direction first (Dirac\_Einstein-Lovelock\_4+4\_Nash.nb,
  section DEFCHART); \(X^M\) is the coordinate with the index \(M\) and
  \(\mathrm{d}^8X = \mathrm{d}x_0\,\mathrm{d}x_1\cdots\mathrm{d}x_7\)
  the coordinate volume element. Unit: length (each coordinate). Status:
  AUTHOR.
\item
  \(\mathcal{F}\): The fold: the seven-dimensional surface
  \(\cos(6Hx_0) = 0\) of the eight-dimensional spacetime (FD-1), in one
  period \(0 < z < \pi\) the surface \(z = \pi/2\), \(x_0 = \pi/(12H)\),
  where \(\sigma = 0\); \(\vert_{\mathcal{F}}\) means restricted to the
  fold. Status: DERIVED.
\item
  \(\mathrm{e}\): Euler's number, the base of the natural logarithm;
  \(\mathrm{e}^{X}\) is the exponential of \(X\). Unit: dimensionless.
  Value: 2.718281828. Status: STANDARD.
\item
  \(a_4\): The author's metric function of the time \(x_4\) that scales
  ordinary 3-space: on the fold \(\mathrm{e}^{a_4}\) is the scale factor
  of cosmology and \(a_4'\) the Hubble rate. Unit: dimensionless.
  Status: AUTHOR.
\item
  \(b_4\): The author's metric function of the time \(x_4\) in the
  extra-time components, \(b_4 > 0\), never identically 1 (the author's
  order of 2026-10-24); the scale factor (the factor of the lengths) of
  the extra times is \(\mathrm{e}^{-a_4}b_4\) on the fold and
  \(\mathrm{e}^{-a_4}b_4\,\mathrm{e}^{-H\lvert\sigma\rvert}\) at the
  distance \(\sigma\) from it (the square root of the warp factor times
  \(\mathrm{e}^{-a_4}b_4\)). Unit: dimensionless. Status: AUTHOR.
\item
  \(x_1\): The first coordinate of ordinary 3-space. Unit: length.
  Status: AUTHOR.
\item
  \(x_2\): The second coordinate of ordinary 3-space. Unit: length.
  Status: AUTHOR.
\item
  \(x_3\): The third coordinate of ordinary 3-space. Unit: length.
  Status: AUTHOR.
\item
  \(x_4\): The time; a prime on a function of \(x_4\) alone is
  \(\mathrm{d}/\mathrm{d}x_4\). Unit: length (= time, \(c = 1\)).
  Status: AUTHOR.
\item
  \(x_5\): The first of the three extra time-like coordinates. Unit:
  length. Status: AUTHOR.
\item
  \(x_6\): The second of the three extra time-like coordinates. Unit:
  length. Status: AUTHOR.
\item
  \(x_7\): The third of the three extra time-like coordinates. Unit:
  length. Status: AUTHOR.
\end{itemize}

with three space-like and four time-like directions; \(a_4\) and \(b_4\)
are the author's two functions of the time \(x_4\). Notation: \(\ln\) is
the natural logarithm, \(\mathrm{e}\) Euler's number,
\(\exp(y) = \mathrm{e}^y\), \(\sin\), \(\cos\) and \(\cot = \cos/\sin\)
the trigonometric functions, \(\tanh\) the hyperbolic tangent,
\(\lvert y\rvert\) the absolute value; \(f(0^+)\) and \(f(0^-)\) are the
limits of \(f(\sigma)\) as \(\sigma\) approaches 0 from above and from
below.

\textbf{Why this surface is called a fold.} Measured by the proper
distance along the hidden direction,

\[
\sigma = \mathrm{sgn}\Big(\frac{\pi}{2} - z\Big)\frac{\ln\sin z}{6H},\qquad \mathrm{d}\sigma = \lvert\cot z\rvert\,\mathrm{d}x_0,\qquad
\sin z = \mathrm{e}^{-6H\lvert\sigma\rvert} \tag{FD-3}
\] \textbf{Symbols in (FD-3).}

\begin{itemize}
\tightlist
\item
  \(\sigma\): The proper coordinate of the hidden direction:
  \(\mathrm{d}\sigma = \lvert\cot z\rvert\,\mathrm{d}x_0\),
  \(\sigma = \mathrm{sgn}(\pi/2 - z)\ln(\sin z)/(6H)\) on one period
  \(0 < z < \pi\); \(\sigma = 0\) is the fold \(z = \pi/2\),
  \(\sigma \to \mp\infty\) the tips \(z \to 0, \pi\); one period of
  \(x_0\) is the whole line \(-\infty < \sigma < \infty\). Unit: length.
  Status: DERIVED.
\item
  \(\mathrm{sgn}\): The sign function: \(\mathrm{sgn}(y) = +1\) for
  \(y > 0\), \(-1\) for \(y < 0\). Status: STANDARD.
\item
  \(\pi\): The number pi, the ratio of the circumference of a circle to
  its diameter. Unit: dimensionless. Value: 3.141592654. Status:
  STANDARD.
\item
  \(z\): The angle variable of the hidden direction, \(z = 6Hx_0\); one
  period of \(x_0\) is \(0 < z < \pi\), the fold is \(z = \pi/2\). Unit:
  dimensionless. Status: AUTHOR.
\item
  \(H\): The author's constant (an inverse length) that fixes the period
  \(\pi/(6H)\) of \(x_0\), the Lovelock couplings
  \(\alpha_{(p)} = w_{(p)}H^{2-2p}\) and the cosmological constant
  \(\Lambda = -7H^2\) (the author's original value \(-10H^2\) was
  corrected on 2026-10-25); in UM-3 a free parameter. Unit: GeV
  (1/length). Status: AUTHOR.
\item
  \(\mathrm{d}\): The differential: \(\mathrm{d}X\) is an infinitesimal
  change of \(X\), \(\mathrm{d}X/\mathrm{d}Y\) the derivative of \(X\)
  with respect to \(Y\), \(\mathrm{d}^nX\) the volume element of \(n\)
  variables. Status: STANDARD.
\item
  \(x_0\): The hidden space-like coordinate (listed first); the metric
  depends on it through \(z = 6Hx_0\), periodically with the period
  \(\pi/(6H)\). Unit: length. Status: AUTHOR.
\item
  \(\mathrm{e}\): Euler's number, the base of the natural logarithm;
  \(\mathrm{e}^{X}\) is the exponential of \(X\). Unit: dimensionless.
  Value: 2.718281828. Status: STANDARD.
\end{itemize}

(the fold is \(\sigma = 0\); the half-period \(0 < z < \pi/2\) is
\(\sigma < 0\), the half-period \(\pi/2 < z < \pi\) is \(\sigma > 0\);
the two ends of the period, the tips \(z \to 0\) and \(z \to \pi\), lie
at \(\sigma \to -\infty\) and \(\sigma \to +\infty\)), the author's
metric is

\[
g_{MN}\,\mathrm{d}X^M\mathrm{d}X^N = \mathrm{d}\sigma^2 + \mathrm{e}^{-2H\lvert\sigma\rvert}\Big[\mathrm{e}^{2a_4}\big(\mathrm{d}x_1^2
 + \mathrm{d}x_2^2 + \mathrm{d}x_3^2\big) - \mathrm{e}^{-2a_4}b_4^2\big(\mathrm{d}x_5^2 + \mathrm{d}x_6^2 + \mathrm{d}x_7^2\big)\Big]
 - \mathrm{d}x_4^2 \tag{FD-4}
\] \textbf{Symbols in (FD-4).}

\begin{itemize}
\tightlist
\item
  \(g_{MN}\): The metric tensor of the eight-dimensional spacetime
  (signature \(+,+,+,+,-,-,-,-\) in the order \(x_0, \dots, x_7\)):
  \(g_{MN}\,\mathrm{d}X^M\mathrm{d}X^N\) is the squared interval,
  \(g^{MN}\) its inverse, \(\lvert g\rvert\) the absolute value of its
  determinant (\(\sqrt{\lvert g\rvert}\,\mathrm{d}^8X\) is the proper
  volume element). Unit: dimensionless. Status: AUTHOR.
\item
  \(M\): A coordinate index of eight dimensions (also with a numbered
  copy \(M_1, M_2, \dots\)): it runs over \(x_0, x_1, \dots, x_7\); a
  repeated index is summed. Status: STANDARD.
\item
  \(N\): A coordinate index of eight dimensions (also
  \(N_1, N_2, \dots\)), running over \(x_0, \dots, x_7\); a repeated
  index is summed. Status: STANDARD.
\item
  \(\mathrm{d}\): The differential: \(\mathrm{d}X\) is an infinitesimal
  change of \(X\), \(\mathrm{d}X/\mathrm{d}Y\) the derivative of \(X\)
  with respect to \(Y\), \(\mathrm{d}^nX\) the volume element of \(n\)
  variables. Status: STANDARD.
\item
  \(X = \{x_0, \dots, x_7\}\): The author's coordinates of the
  eight-dimensional spacetime, \(X = \{x_0, x_1, \dots, x_7\}\) with the
  hidden direction first (Dirac\_Einstein-Lovelock\_4+4\_Nash.nb,
  section DEFCHART); \(X^M\) is the coordinate with the index \(M\) and
  \(\mathrm{d}^8X = \mathrm{d}x_0\,\mathrm{d}x_1\cdots\mathrm{d}x_7\)
  the coordinate volume element. Unit: length (each coordinate). Status:
  AUTHOR.
\item
  \(\sigma\): The proper coordinate of the hidden direction:
  \(\mathrm{d}\sigma = \lvert\cot z\rvert\,\mathrm{d}x_0\),
  \(\sigma = \mathrm{sgn}(\pi/2 - z)\ln(\sin z)/(6H)\) on one period
  \(0 < z < \pi\); \(\sigma = 0\) is the fold \(z = \pi/2\),
  \(\sigma \to \mp\infty\) the tips \(z \to 0, \pi\); one period of
  \(x_0\) is the whole line \(-\infty < \sigma < \infty\). Unit: length.
  Status: DERIVED.
\item
  \(\mathrm{e}\): Euler's number, the base of the natural logarithm;
  \(\mathrm{e}^{X}\) is the exponential of \(X\). Unit: dimensionless.
  Value: 2.718281828. Status: STANDARD.
\item
  \(H\): The author's constant (an inverse length) that fixes the period
  \(\pi/(6H)\) of \(x_0\), the Lovelock couplings
  \(\alpha_{(p)} = w_{(p)}H^{2-2p}\) and the cosmological constant
  \(\Lambda = -7H^2\) (the author's original value \(-10H^2\) was
  corrected on 2026-10-25); in UM-3 a free parameter. Unit: GeV
  (1/length). Status: AUTHOR.
\item
  \(a_4\): The author's metric function of the time \(x_4\) that scales
  ordinary 3-space: on the fold \(\mathrm{e}^{a_4}\) is the scale factor
  of cosmology and \(a_4'\) the Hubble rate. Unit: dimensionless.
  Status: AUTHOR.
\item
  \(b_4\): The author's metric function of the time \(x_4\) in the
  extra-time components, \(b_4 > 0\), never identically 1 (the author's
  order of 2026-10-24); the scale factor (the factor of the lengths) of
  the extra times is \(\mathrm{e}^{-a_4}b_4\) on the fold and
  \(\mathrm{e}^{-a_4}b_4\,\mathrm{e}^{-H\lvert\sigma\rvert}\) at the
  distance \(\sigma\) from it (the square root of the warp factor times
  \(\mathrm{e}^{-a_4}b_4\)). Unit: dimensionless. Status: AUTHOR.
\item
  \(x_1\): The first coordinate of ordinary 3-space. Unit: length.
  Status: AUTHOR.
\item
  \(x_2\): The second coordinate of ordinary 3-space. Unit: length.
  Status: AUTHOR.
\item
  \(x_3\): The third coordinate of ordinary 3-space. Unit: length.
  Status: AUTHOR.
\item
  \(x_4\): The time; a prime on a function of \(x_4\) alone is
  \(\mathrm{d}/\mathrm{d}x_4\). Unit: length (= time, \(c = 1\)).
  Status: AUTHOR.
\item
  \(x_5\): The first of the three extra time-like coordinates. Unit:
  length. Status: AUTHOR.
\item
  \(x_6\): The second of the three extra time-like coordinates. Unit:
  length. Status: AUTHOR.
\item
  \(x_7\): The third of the three extra time-like coordinates. Unit:
  length. Status: AUTHOR.
\end{itemize}

It depends on \(\sigma\) only through \(\lvert\sigma\rvert\). The
function \(\mathcal{W}(\sigma) = \mathrm{e}^{-2H\lvert\sigma\rvert}\),
which multiplies the squared steps of ordinary space and of the extra
times, is the warp factor of the metric (its square root
\(\mathrm{e}^{-H\lvert\sigma\rvert}\) multiplies the lengths); the time
\(x_4\) is not warped. Four properties make the surface \(\sigma = 0\) a
fold, and together they single it out within the period:

\begin{enumerate}
\def\labelenumi{\arabic{enumi}.}
\tightlist
\item
  \emph{It is a mirror.} The reflection \(\sigma \to -\sigma\) (that is
  \(z \to \pi - z\), or \(x_0 \to \pi/(6H) - x_0\)) maps the metric to
  itself, and the fold is the only surface of the period that the
  reflection leaves in place: the two halves of the period are mirror
  images of each other, joined along the fold, as the two halves of a
  folded sheet of paper are joined along the crease.
\item
  \emph{The warp factor is largest on it.} The warp factor
  \((\sin^2z)^{1/6} = \mathrm{e}^{-2H\lvert\sigma\rvert}\) equals 1 on
  the fold and decreases on both sides.
\item
  \emph{The metric has a kink on it.} The metric is continuous across
  the fold, but its first derivative along \(\sigma\) jumps. The
  extrinsic curvature of the surfaces \(\sigma = \text{const}\) changes
  sign across the fold,
\end{enumerate}

\[
K^a{}_b = \frac12\,g^{ac}\,\partial_\sigma g_{cb}:\qquad
K^{1}{}_{1} = K^{2}{}_{2} = K^{3}{}_{3} = K^{5}{}_{5} = K^{6}{}_{6} = K^{7}{}_{7} = -H\,\mathrm{sgn}(\sigma),\qquad K^{4}{}_{4} = 0 \tag{FD-5}
\] \textbf{Symbols in (FD-5).}

\begin{itemize}
\tightlist
\item
  \(K^a{}_b\): The extrinsic curvature of a surface
  \(\sigma = \mathrm{const}\) of the hidden direction,
  \(K^a{}_b = \frac12 g^{ac}\partial_\sigma g_{cb}\) for \(a, b\) along
  the surface (FD-5): \(-H\,\mathrm{sgn}(\sigma)\) along ordinary space
  and the extra times, 0 along the time. Unit: GeV (1/length). Status:
  DERIVED.
\item
  \(a\): A frame (orthonormal) index running over \(0, 1, \dots, 7\)
  (frame direction \((0)\) along \(\sigma\), \((1), (2), (3)\) along
  ordinary space, \((4)\) along the time, \((5), (6), (7)\) along the
  extra times); raised and lowered with the frame metric
  \(\eta_{ab} = \mathrm{diag}(1, 1, 1, 1, -1, -1, -1, -1)\); on the fold
  (\(a, b \in \{1, \dots, 7\}\)) the frame of the fold; a repeated index
  is summed. Status: STANDARD.
\item
  \(b\): A frame index (as \(a\)). Status: STANDARD.
\item
  \(c\): A frame index (as \(a\)). Status: STANDARD.
\item
  \(g_{MN}\): The metric tensor of the eight-dimensional spacetime
  (signature \(+,+,+,+,-,-,-,-\) in the order \(x_0, \dots, x_7\)):
  \(g_{MN}\,\mathrm{d}X^M\mathrm{d}X^N\) is the squared interval,
  \(g^{MN}\) its inverse, \(\lvert g\rvert\) the absolute value of its
  determinant (\(\sqrt{\lvert g\rvert}\,\mathrm{d}^8X\) is the proper
  volume element). Unit: dimensionless. Status: AUTHOR.
\item
  \(\partial\): A partial derivative: \(\partial_M\) with respect to the
  coordinate \(X^M\), \(\partial_\mu\) with respect to \(x^\mu\),
  \(\partial_\sigma\) with respect to \(\sigma\),
  \(\partial X/\partial Y\) the derivative of \(X\) with respect to
  \(Y\) at fixed other variables. Status: STANDARD.
\item
  \(\sigma\): The label \(\sigma\) in \(\partial_\sigma\): the
  derivative with respect to the proper hidden coordinate \(\sigma\).
  Status: STANDARD.
\item
  \(\sigma\): The proper coordinate of the hidden direction:
  \(\mathrm{d}\sigma = \lvert\cot z\rvert\,\mathrm{d}x_0\),
  \(\sigma = \mathrm{sgn}(\pi/2 - z)\ln(\sin z)/(6H)\) on one period
  \(0 < z < \pi\); \(\sigma = 0\) is the fold \(z = \pi/2\),
  \(\sigma \to \mp\infty\) the tips \(z \to 0, \pi\); one period of
  \(x_0\) is the whole line \(-\infty < \sigma < \infty\). Unit: length.
  Status: DERIVED.
\item
  \(H\): The author's constant (an inverse length) that fixes the period
  \(\pi/(6H)\) of \(x_0\), the Lovelock couplings
  \(\alpha_{(p)} = w_{(p)}H^{2-2p}\) and the cosmological constant
  \(\Lambda = -7H^2\) (the author's original value \(-10H^2\) was
  corrected on 2026-10-25); in UM-3 a free parameter. Unit: GeV
  (1/length). Status: AUTHOR.
\item
  \(\mathrm{sgn}\): The sign function: \(\mathrm{sgn}(y) = +1\) for
  \(y > 0\), \(-1\) for \(y < 0\). Status: STANDARD.
\end{itemize}

(\(a, b\) frame indices along the fold); it jumps by \(-2H\) along
ordinary space and along the extra times. The Einstein-Lovelock
equations of UM-3,
\(\sum_p\alpha_{(p)}\mathcal{E}_{(p)}{}^a{}_b + \Lambda\delta^a{}_b = \kappa\mathcal{T}^a{}_b\)
(\(\mathcal{E}_{(p)}\) the Lovelock tensor of order \(p = 1, 2, 3\),
\(\alpha_{(p)} = w_{(p)}H^{2-2p}\) its coupling with the weights
\(w = (1, 1/20, 1/360)\), \(\Lambda = -7H^2\) the cosmological constant,
\(\kappa = 8\pi G_8\) the gravitational coupling of eight dimensions,
\(\mathcal{T}^a{}_b\) the energy-momentum tensor, \(\delta^a{}_b\) the
Kronecker delta), are linear in the second \(\sigma\)-derivatives of the
metric, so the jump puts a Dirac delta function
\(\delta_{\mathrm{D}}(\sigma)\) into them: the metric solves them only
if the fold carries a thin source \(\mathcal{S}^a{}_b\), an
energy-momentum tensor concentrated on \(\sigma = 0\), given by the
integral of the equations across the fold (the junction conditions),

\[
\kappa\,\mathcal{S}^a{}_b = \int_{0^-}^{0^+}\mathrm{d}\sigma\,\sum_{p=1}^{3}\alpha_{(p)}\,\mathcal{E}_{(p)}{}^a{}_b \tag{FD-11}
\] \textbf{Symbols in (FD-11).}

\begin{itemize}
\tightlist
\item
  \(\kappa\): The eight-dimensional gravitational coupling,
  \(\kappa = 8\pi G_8\). Unit: GeV\(^{-6}\). Status: AUTHOR.
\item
  \(\mathcal{S}^a{}_b\): The thin source of the fold: the
  energy-momentum tensor concentrated on the fold, integrated across it,
  \(\kappa\mathcal{S}^a{}_b = \int_{0^-}^{0^+}\mathrm{d}\sigma\sum_p\alpha_{(p)}\mathcal{E}_{(p)}{}^a{}_b\)
  (FD-11, the junction conditions); \(\mathcal{S}^4{}_4\) is minus its
  energy density, \(\mathcal{S}^1{}_1\) and \(\mathcal{S}^5{}_5\) its
  pressures along ordinary space and along the extra times (frame
  components along the fold, eight-dimensional normalisation: an energy
  per unit volume of the six other directions of the fold). Unit:
  GeV\(^7\). Value:
  \(\kappa\mathcal{S}^a{}_b = -\frac{32}{5}H\delta^a{}_b\) at the static
  block. Status: DERIVED.
\item
  \(a\): A frame (orthonormal) index running over \(0, 1, \dots, 7\)
  (frame direction \((0)\) along \(\sigma\), \((1), (2), (3)\) along
  ordinary space, \((4)\) along the time, \((5), (6), (7)\) along the
  extra times); raised and lowered with the frame metric
  \(\eta_{ab} = \mathrm{diag}(1, 1, 1, 1, -1, -1, -1, -1)\); on the fold
  (\(a, b \in \{1, \dots, 7\}\)) the frame of the fold; a repeated index
  is summed. Status: STANDARD.
\item
  \(b\): A frame index (as \(a\)). Status: STANDARD.
\item
  \(\mathrm{d}\): The differential: \(\mathrm{d}X\) is an infinitesimal
  change of \(X\), \(\mathrm{d}X/\mathrm{d}Y\) the derivative of \(X\)
  with respect to \(Y\), \(\mathrm{d}^nX\) the volume element of \(n\)
  variables. Status: STANDARD.
\item
  \(\sigma\): The proper coordinate of the hidden direction:
  \(\mathrm{d}\sigma = \lvert\cot z\rvert\,\mathrm{d}x_0\),
  \(\sigma = \mathrm{sgn}(\pi/2 - z)\ln(\sin z)/(6H)\) on one period
  \(0 < z < \pi\); \(\sigma = 0\) is the fold \(z = \pi/2\),
  \(\sigma \to \mp\infty\) the tips \(z \to 0, \pi\); one period of
  \(x_0\) is the whole line \(-\infty < \sigma < \infty\). Unit: length.
  Status: DERIVED.
\item
  \(p\): The Lovelock order \(p = 1, 2, 3\) (Einstein, Gauss-Bonnet,
  cubic), as an index and as a number; in the Kohn-Sham equations a
  second label of a Kohn-Sham orbital (as \(o\), with \(\delta_{op}\)
  the Kronecker delta of the two labels). Status: STANDARD.
\item
  \(\alpha_{(p)}\): The Lovelock coupling of order \(p\),
  \(\alpha_{(p)} = w_{(p)}H^{2-2p}\) (\(\alpha_{(1)} = 1\)). Unit:
  GeV\(^{2-2p}\). Status: AUTHOR.
\item
  \(\mathcal{E}_{(p)}{}^M{}_N\): The Lovelock tensor of order \(p\): the
  variation of \(\sqrt{\lvert g\rvert}\,\mathcal{L}_{(p)}\) with respect
  to the metric (\(\mathcal{E}_{(1)}\) the Einstein tensor,
  \(\mathcal{E}_{(2)}\) the Gauss-Bonnet tensor); with frame indices its
  frame components. Unit: 1/length\(^{2p}\). Status: STANDARD.
\end{itemize}

with \(\mathcal{S}^4{}_4 = -(\text{its energy density})\),
\(\mathcal{S}^1{}_1 = \mathcal{S}^2{}_2 = \mathcal{S}^3{}_3\) its
pressure along ordinary space and
\(\mathcal{S}^5{}_5 = \mathcal{S}^6{}_6 = \mathcal{S}^7{}_7\) its
pressure along the extra times; \(\mathcal{S}\) is \(\mathcal{T}\)
integrated across the fold, an energy per unit volume of the six other
directions of the fold (unit GeV\(^7\)). On the static block (\(a_4\)
and \(b_4\) constant) and on an expanding fold whose extra times are
frozen (\(b_4\) proportional to \(\mathrm{e}^{a_4}\);
\(\mathsf{A} = a_4'/H\), \(\mathsf{A}_1 = a_4''/H^2\))

\[
\begin{aligned}
&\text{static block:}\quad \kappa\mathcal{S}^a{}_b = -\frac{32}{5}H\,\delta^a{}_b\quad\text{(a pure tension)},\\
&\text{frozen extra times:}\quad \kappa\mathcal{S}^1{}_1 = -\frac{16}{15}\big(6 + 3\mathsf{A}^2 + 2\mathsf{A}_1\big)H,\quad
 \kappa\mathcal{S}^4{}_4 = -\frac{16}{5}\big(2 + \mathsf{A}^2\big)H,\\
&\qquad\qquad \kappa\mathcal{S}^5{}_5 = -\frac45\big(8 + 8\mathsf{A}^2 + \mathsf{A}^4 + 4\mathsf{A}_1 + \mathsf{A}^2\mathsf{A}_1\big)H
\end{aligned} \tag{FD-12}
\] \textbf{Symbols in (FD-12).}

\begin{itemize}
\tightlist
\item
  \(\kappa\): The eight-dimensional gravitational coupling,
  \(\kappa = 8\pi G_8\). Unit: GeV\(^{-6}\). Status: AUTHOR.
\item
  \(\mathcal{S}^a{}_b\): The thin source of the fold: the
  energy-momentum tensor concentrated on the fold, integrated across it,
  \(\kappa\mathcal{S}^a{}_b = \int_{0^-}^{0^+}\mathrm{d}\sigma\sum_p\alpha_{(p)}\mathcal{E}_{(p)}{}^a{}_b\)
  (FD-11, the junction conditions); \(\mathcal{S}^4{}_4\) is minus its
  energy density, \(\mathcal{S}^1{}_1\) and \(\mathcal{S}^5{}_5\) its
  pressures along ordinary space and along the extra times (frame
  components along the fold, eight-dimensional normalisation: an energy
  per unit volume of the six other directions of the fold). Unit:
  GeV\(^7\). Value:
  \(\kappa\mathcal{S}^a{}_b = -\frac{32}{5}H\delta^a{}_b\) at the static
  block. Status: DERIVED.
\item
  \(a\): A frame (orthonormal) index running over \(0, 1, \dots, 7\)
  (frame direction \((0)\) along \(\sigma\), \((1), (2), (3)\) along
  ordinary space, \((4)\) along the time, \((5), (6), (7)\) along the
  extra times); raised and lowered with the frame metric
  \(\eta_{ab} = \mathrm{diag}(1, 1, 1, 1, -1, -1, -1, -1)\); on the fold
  (\(a, b \in \{1, \dots, 7\}\)) the frame of the fold; a repeated index
  is summed. Status: STANDARD.
\item
  \(b\): A frame index (as \(a\)). Status: STANDARD.
\item
  \(H\): The author's constant (an inverse length) that fixes the period
  \(\pi/(6H)\) of \(x_0\), the Lovelock couplings
  \(\alpha_{(p)} = w_{(p)}H^{2-2p}\) and the cosmological constant
  \(\Lambda = -7H^2\) (the author's original value \(-10H^2\) was
  corrected on 2026-10-25); in UM-3 a free parameter. Unit: GeV
  (1/length). Status: AUTHOR.
\item
  \(\delta\): The Kronecker delta: with two indices \(\delta^M_N = 1\)
  if \(M = N\) and 0 otherwise (likewise \(\delta_{X\mathrm{L}}\) is 1
  if the label \(X\) of a half of dirac16complex is L and 0 otherwise);
  with upper and lower index strings the generalised Kronecker delta
  \(\delta^{M_1\cdots M_n}_{N_1\cdots N_n}\), the determinant of the
  \(n \times n\) matrix of the Kronecker deltas \(\delta^{M_a}_{N_b}\)
  (\(+1\) or \(-1\) if the upper indices are an even or odd permutation
  of the lower ones, 0 otherwise). Status: STANDARD.
\item
  \(\mathsf{A}\): The expansion rate of ordinary space in units of
  \(H\), \(\mathsf{A} = a_4'/H\). Unit: dimensionless. Value: 0 at the
  static block; \(H_4^{(0)}/H\) today (\(H_4^{(0)}\) the Hubble
  constant); 1 at the author's vacuum, which is not a solution with the
  corrected weights.
\item
  \(\mathsf{A}_1\): The time derivative of \(\mathsf{A}\) in units of
  \(H\), \(\mathsf{A}_1 = \mathsf{A}'/H = a_4''/H^2\). Unit:
  dimensionless.
\end{itemize}

the tension plus matter that obeys the Friedmann equations of general
relativity. The fold is therefore not only a place but a physical
object: a thin layer that carries energy and momentum (a brane). 4.
\emph{The author's coordinate \(x_0\) degenerates on it.} The metric
component \(g_{x_0x_0} = \cot^2(6Hx_0)\) vanishes on the fold and
nowhere else in the period; near the fold the proper distance grows with
the square of the coordinate distance,

\[
\lvert\sigma\rvert = 3H\Big(x_0 - \frac{\pi}{12H}\Big)^2\Big[1 + 6H^2\Big(x_0 - \frac{\pi}{12H}\Big)^2 + \dots\Big] \tag{FD-6}
\] \textbf{Symbols in (FD-6).}

\begin{itemize}
\tightlist
\item
  \(\sigma\): The proper coordinate of the hidden direction:
  \(\mathrm{d}\sigma = \lvert\cot z\rvert\,\mathrm{d}x_0\),
  \(\sigma = \mathrm{sgn}(\pi/2 - z)\ln(\sin z)/(6H)\) on one period
  \(0 < z < \pi\); \(\sigma = 0\) is the fold \(z = \pi/2\),
  \(\sigma \to \mp\infty\) the tips \(z \to 0, \pi\); one period of
  \(x_0\) is the whole line \(-\infty < \sigma < \infty\). Unit: length.
  Status: DERIVED.
\item
  \(H\): The author's constant (an inverse length) that fixes the period
  \(\pi/(6H)\) of \(x_0\), the Lovelock couplings
  \(\alpha_{(p)} = w_{(p)}H^{2-2p}\) and the cosmological constant
  \(\Lambda = -7H^2\) (the author's original value \(-10H^2\) was
  corrected on 2026-10-25); in UM-3 a free parameter. Unit: GeV
  (1/length). Status: AUTHOR.
\item
  \(x_0\): The hidden space-like coordinate (listed first); the metric
  depends on it through \(z = 6Hx_0\), periodically with the period
  \(\pi/(6H)\). Unit: length. Status: AUTHOR.
\item
  \(\pi\): The number pi, the ratio of the circumference of a circle to
  its diameter. Unit: dimensionless. Value: 3.141592654. Status:
  STANDARD.
\end{itemize}

so the two coordinate values on either side of \(x_0 = \pi/(12H)\) at
the same coordinate distance lie at the same proper distance from the
fold: the coordinate line of \(x_0\) is folded onto the line of proper
distances (a fold map, example 5).

\textbf{What lives on the fold.} In UM-3 the Standard Model (the Higgs
doublet with the author's ad hoc vacuum, a time-dependent choice of its
vacuum value; the gauge fields of
\(SU(3)_C \times SU(2)_L \times U(1)_Y\); three generations of quarks
and leptons and three right-handed neutrinos) and the fold fluid (its
tension and the pressure along the extra times that keeps their volume
fixed) live on the fold: they are fields at \(\sigma = 0\), functions of
\(x_1, \dots, x_4\) (no field depends on the extra times). Integrated
over the extra times, the fold is the four-dimensional universe that we
observe, with the metric
\(\mathrm{diag}(\mathrm{e}^{2a_4}, \mathrm{e}^{2a_4}, \mathrm{e}^{2a_4}, -1)\)
in \(x_1, \dots, x_4\). dirac16complex (the author's spinor field of
eight dimensions) and gravity live in the bulk \(\sigma \ne 0\) as well
as on the fold; the bound states \(\chi\) of dirac16complex are the part
of it that the fold holds, and the portals (couplings between
dirac16complex and Standard-Model fields) act at \(\sigma = 0\). In UM-3
files ``the fold'' means this surface; the examples below use the word
for other surfaces with the same properties. The word is not used in its
other senses: a set of equal energy levels is called degenerate, and a
growth by the factor \(\mathrm{e}\) is called an e-folding.

\textbf{Examples of folds that are meaningful for UM-3.}

\begin{enumerate}
\def\labelenumi{\arabic{enumi}.}
\tightlist
\item
  \emph{Our fold.} The fold \(z = \pi/2\) of the period that contains
  our world (FD-1): the seven-dimensional world-volume of the Standard
  Model, and after the integration over the extra times the
  four-dimensional universe that we observe.
\item
  \emph{The copies of the fold.} Because the metric repeats with the
  period \(\pi/(6H)\), every surface
  \(x_0 = \pi/(12H) + (\text{an integer multiple of } \pi/(6H))\) is a
  copy of the fold, one per period. Neighbouring copies are separated by
  a tip (\(z\) a multiple of \(\pi\)), where the warp factor
  \((\sin^2z)^{1/6}\) vanishes and \(g_{x_0x_0} = \cot^2z\) diverges,
  and which lies at an infinite proper distance from both: in proper
  distance every copy has its own complete line
  \(-\infty < \sigma < \infty\). A compact direction of length \(L\)
  would give every field a tower of particles with the masses
  \(2\pi n/L\), \(n = 1, 2, \dots\) (a Kaluza-Klein tower); because the
  copies are infinitely far apart, the period \(\pi/(6H)\) of \(x_0\)
  produces no such tower with the spacing \(12H\), and the \(x_0\) mass
  of the bound states is their binding to the fold.
\item
  \emph{A folded sheet of paper.} The crease of a folded sheet is a fold
  in this sense: the two halves are mirror images (property 1) joined
  along the crease, the sheet is continuous there and its direction
  changes abruptly (property 3). In UM-3 the direction across the crease
  is the hidden direction \(\sigma\), the seven directions along the
  crease are \(x_1, \dots, x_7\), and the crease is the fold.
\item
  \emph{A folded function.} The warp factor
  \(\mathrm{e}^{-2H\lvert\sigma\rvert}\) is the simplest folded
  function: continuous, mirror symmetric, largest at \(\sigma = 0\),
  with a jump of its derivative,
\end{enumerate}

\[
\frac{\mathrm{d}}{\mathrm{d}\sigma}\mathrm{e}^{-2H\lvert\sigma\rvert} = -2H\,\mathrm{sgn}(\sigma)\,\mathrm{e}^{-2H\lvert\sigma\rvert},\qquad
\frac{\mathrm{d}^2}{\mathrm{d}\sigma^2}\mathrm{e}^{-2H\lvert\sigma\rvert} = 4H^2\,\mathrm{e}^{-2H\lvert\sigma\rvert} - 4H\,\delta_{\mathrm{D}}(\sigma) \tag{FD-7}
\] \textbf{Symbols in (FD-7).}

\begin{itemize}
\tightlist
\item
  \(\mathrm{d}\): The differential: \(\mathrm{d}X\) is an infinitesimal
  change of \(X\), \(\mathrm{d}X/\mathrm{d}Y\) the derivative of \(X\)
  with respect to \(Y\), \(\mathrm{d}^nX\) the volume element of \(n\)
  variables. Status: STANDARD.
\item
  \(\sigma\): The proper coordinate of the hidden direction:
  \(\mathrm{d}\sigma = \lvert\cot z\rvert\,\mathrm{d}x_0\),
  \(\sigma = \mathrm{sgn}(\pi/2 - z)\ln(\sin z)/(6H)\) on one period
  \(0 < z < \pi\); \(\sigma = 0\) is the fold \(z = \pi/2\),
  \(\sigma \to \mp\infty\) the tips \(z \to 0, \pi\); one period of
  \(x_0\) is the whole line \(-\infty < \sigma < \infty\). Unit: length.
  Status: DERIVED.
\item
  \(\mathrm{e}\): Euler's number, the base of the natural logarithm;
  \(\mathrm{e}^{X}\) is the exponential of \(X\). Unit: dimensionless.
  Value: 2.718281828. Status: STANDARD.
\item
  \(H\): The author's constant (an inverse length) that fixes the period
  \(\pi/(6H)\) of \(x_0\), the Lovelock couplings
  \(\alpha_{(p)} = w_{(p)}H^{2-2p}\) and the cosmological constant
  \(\Lambda = -7H^2\) (the author's original value \(-10H^2\) was
  corrected on 2026-10-25); in UM-3 a free parameter. Unit: GeV
  (1/length). Status: AUTHOR.
\item
  \(\mathrm{sgn}\): The sign function: \(\mathrm{sgn}(y) = +1\) for
  \(y > 0\), \(-1\) for \(y < 0\). Status: STANDARD.
\item
  \(\delta_{\mathrm{D}}\): The Dirac delta function:
  \(\int\mathrm{d}\sigma\,\delta_{\mathrm{D}}(\sigma)F(\sigma) = F(0)\).
  Unit: 1/length (for \(\delta_{\mathrm{D}}(\sigma)\)). Status:
  STANDARD.
\end{itemize}

whose delta function is the origin of the thin source of property 3.
Every metric function of FD-4 depends on \(\sigma\) through
\(\lvert\sigma\rvert\) and is folded at \(\sigma = 0\) in this sense. 5.
\emph{A fold map.} In the mathematics of smooth maps the fold (Whitney's
fold singularity) is the map \(u \mapsto u^2\) of a line onto a
half-line: it folds the line at \(u = 0\), two points \(\pm u\) land on
the same point, and the derivative of the map vanishes on the fold. Near
the fold the map from the author's coordinate \(x_0\) to the distance
\(\lvert\sigma\rvert\) from the fold is of this kind (FD-6), and
\(g_{x_0x_0} = \cot^2z\), the square of the derivative
\(\mathrm{d}\sigma/\mathrm{d}x_0\), vanishes there. 6. \emph{The brane
of Randall and Sundrum.} The five-dimensional metric

\[
\mathrm{d}y_{\mathrm{RS}}^2 + \mathrm{e}^{-2k_{\mathrm{RS}}\lvert y_{\mathrm{RS}}\rvert}\big(\mathrm{d}x_1^2 + \mathrm{d}x_2^2 + \mathrm{d}x_3^2
 - \mathrm{d}x_4^2\big),\qquad \lambda_{\mathrm{RS}} = \frac{6k_{\mathrm{RS}}}{\kappa_5},\qquad \Lambda_{\mathrm{RS}} = -\frac{6k_{\mathrm{RS}}^2}{\kappa_5} \tag{FD-8}
\] \textbf{Symbols in (FD-8).}

\begin{itemize}
\tightlist
\item
  \(\mathrm{d}\): The differential: \(\mathrm{d}X\) is an infinitesimal
  change of \(X\), \(\mathrm{d}X/\mathrm{d}Y\) the derivative of \(X\)
  with respect to \(Y\), \(\mathrm{d}^nX\) the volume element of \(n\)
  variables. Status: STANDARD.
\item
  \(y_{\mathrm{RS}}\): The proper coordinate across the brane of the
  five-dimensional model of Randall and Sundrum (FD-8, an example, not
  UM-3). Unit: length. Status: STANDARD.
\item
  \(\mathrm{e}\): Euler's number, the base of the natural logarithm;
  \(\mathrm{e}^{X}\) is the exponential of \(X\). Unit: dimensionless.
  Value: 2.718281828. Status: STANDARD.
\item
  \(k_{\mathrm{RS}}\): The inverse curvature length of the bulk of the
  model of Randall and Sundrum: the warp factor of the brane is
  \(\mathrm{e}^{-2k_{\mathrm{RS}}\lvert y_{\mathrm{RS}}\rvert}\) (FD-8,
  an example, not UM-3). Unit: GeV (1/length). Status: STANDARD.
\item
  \(x_1\): The first coordinate of ordinary 3-space. Unit: length.
  Status: AUTHOR.
\item
  \(x_2\): The second coordinate of ordinary 3-space. Unit: length.
  Status: AUTHOR.
\item
  \(x_3\): The third coordinate of ordinary 3-space. Unit: length.
  Status: AUTHOR.
\item
  \(x_4\): The time; a prime on a function of \(x_4\) alone is
  \(\mathrm{d}/\mathrm{d}x_4\). Unit: length (= time, \(c = 1\)).
  Status: AUTHOR.
\item
  \(\lambda_{\mathrm{RS}}\): The tension (energy per unit volume of
  ordinary space) of the brane of the model of Randall and Sundrum,
  \(\lambda_{\mathrm{RS}} = 6k_{\mathrm{RS}}/\kappa_5\) (FD-8, an
  example, not UM-3). Unit: GeV\(^4\). Status: STANDARD.
\item
  \(\kappa_5\): The gravitational coupling of five dimensions in the
  model of Randall and Sundrum, \(\kappa_5 = 8\pi G_5\) (FD-8, an
  example, not UM-3). Unit: GeV\(^{-3}\). Status: STANDARD.
\item
  \(\Lambda_{\mathrm{RS}}\): The vacuum energy density of the bulk of
  the model of Randall and Sundrum,
  \(\Lambda_{\mathrm{RS}} = -6k_{\mathrm{RS}}^2/\kappa_5\) (FD-8, an
  example, not UM-3). Unit: GeV\(^5\). Status: STANDARD.
\end{itemize}

has a fold at \(y_{\mathrm{RS}} = 0\) with the properties 1 to 3 (its
coordinate \(y_{\mathrm{RS}}\) is already the proper distance, so
property 4 does not arise); Einstein's equations require there a brane
of the tension \(\lambda_{\mathrm{RS}}\) (Israel's junction condition)
and in the bulk the vacuum energy \(\Lambda_{\mathrm{RS}}\), and the
bulk is otherwise empty. The fold of UM-3 is of the same kind: its
static block is a vacuum of the bulk with a fold of pure tension
(FD-12). It differs in that the time \(x_4\) is not warped, three extra
times are warped together with ordinary space, the Lovelock terms add to
the junction conditions, and the space is not cut at the brane: Randall
and Sundrum identify the two sides of their brane (one physical side),
while the fold of UM-3 has two distinct sides (example 7). 7. \emph{A
particle bound to a thin layer.} A particle held by an attractive thin
layer has a wave function that falls off exponentially with the distance
on both sides, and the layer makes its derivative jump (for the
first-order Dirac equation, the wave function itself). In UM-3 the
particles \(\chi\) of dark matter are the states of dirac16complex bound
to the fold; at rest and without their self-interaction (the field is
\(\Psi = \mathrm{e}^{3H\lvert\sigma\rvert}\chi(\sigma)\mathrm{e}^{-\mathrm{i}\omega_{\mathrm{b}}x_4}\))

\[
\begin{aligned}
&\chi(\sigma) = \mathrm{e}^{-q_{\mathrm{b}}\lvert\sigma\rvert}\chi(0^+)\ \ (\sigma > 0),\qquad
\chi(\sigma) = \mathrm{e}^{-q_{\mathrm{b}}\lvert\sigma\rvert}\chi(0^-)\ \ (\sigma < 0),\\
&\chi(0^+) = \exp\Big(\gamma^{(0)}\big(\mu_{\mathrm{L}}P_{\mathrm{L}} + \mu_{\mathrm{R}}P_{\mathrm{R}}\big)\Big)\chi(0^-),\qquad
\omega_{\mathrm{b}} = \sqrt{m_\Psi^2 - q_{\mathrm{b}}^2}
\end{aligned} \tag{FD-9}
\] \textbf{Symbols in (FD-9).}

\begin{itemize}
\tightlist
\item
  \(\chi(\sigma)\): The profile of a mode of dirac16complex along the
  hidden direction,
  \(\Psi = \mathrm{e}^{3H\lvert\sigma\rvert}\chi(\sigma)\mathrm{e}^{\mathrm{i}(\vec k\cdot\vec x - \omega x_4)}\)
  (a 16-component spinor function of \(\sigma\)); as a four-dimensional
  field, \(\chi\) denotes the particle of a bound state. Unit:
  GeV\(^{2}\) (with the \(V_*\) normalisation). Status: DERIVED.
\item
  \(\sigma\): The proper coordinate of the hidden direction:
  \(\mathrm{d}\sigma = \lvert\cot z\rvert\,\mathrm{d}x_0\),
  \(\sigma = \mathrm{sgn}(\pi/2 - z)\ln(\sin z)/(6H)\) on one period
  \(0 < z < \pi\); \(\sigma = 0\) is the fold \(z = \pi/2\),
  \(\sigma \to \mp\infty\) the tips \(z \to 0, \pi\); one period of
  \(x_0\) is the whole line \(-\infty < \sigma < \infty\). Unit: length.
  Status: DERIVED.
\item
  \(\mathrm{e}\): Euler's number, the base of the natural logarithm;
  \(\mathrm{e}^{X}\) is the exponential of \(X\). Unit: dimensionless.
  Value: 2.718281828. Status: STANDARD.
\item
  \(q_{\mathrm{b}}\): The inverse width of a state of dirac16complex
  bound to the fold: its profile falls as
  \(\mathrm{e}^{-q_{\mathrm{b}}\lvert\sigma\rvert}\). Unit: GeV. Status:
  DERIVED.
\item
  \(\gamma^M,\ \gamma^{(a)}\): The Dirac matrices of dirac16complex
  along the coordinates, \(\gamma^M = \mathbf{e}_a{}^MT^a\),
  \(\gamma_M = g_{MN}\gamma^N\); \(\gamma^{(a)} = T^a\) along the frame
  direction \((a)\) (\(\gamma^{(0)}\) along \(\sigma\), \(\gamma^{(4)}\)
  along the time). Unit: dimensionless. Status: AUTHOR.
\item
  \(\mu_X,\ \mu_{\mathrm{L}},\ \mu_{\mathrm{R}}\): The dimensionless
  strengths of the fold for the left-handed and right-handed halves of
  dirac16complex, set by the mean fields of the Standard Model at the
  fold: the jump of the profile across the fold is the exponential of
  \(\gamma^{(0)}(\mu_{\mathrm{L}}P_{\mathrm{L}} + \mu_{\mathrm{R}}P_{\mathrm{R}})\);
  negative values attract. Unit: dimensionless. Status: DERIVED.
\item
  \(P_{\mathrm{L}}\): The author's left-handed chiral projector,
  \(P_{\mathrm{L}} = (\mathbf{1} - \Gamma_8)/2\). Unit: dimensionless.
  Status: AUTHOR.
\item
  \(P_{\mathrm{R}}\): The author's right-handed chiral projector,
  \(P_{\mathrm{R}} = (\mathbf{1} + \Gamma_8)/2\). Unit: dimensionless.
  Status: AUTHOR.
\item
  \(\Gamma_8\): The chirality matrix of eight dimensions,
  \(\Gamma_8 = T_0T_1\cdots T_7\) (the author's T16\^{}A{[}8{]})
  \(= \mathrm{diag}(-\mathbf{1}_8, \mathbf{1}_8)\);
  \(\Gamma_8^2 = \mathbf{1}\); it anticommutes with every \(T_a\). Unit:
  dimensionless. Status: AUTHOR.
\item
  \(\mathbf{1},\ \mathbf{1}_8\): The identity matrix of the size fixed
  by the context (\(16 \times 16\) on dirac16complex, \(4 \times 4\) on
  a four-dimensional spinor, \(3 \times 3\) in generation or flavour
  space); \(\mathbf{1}_8\) is the \(8 \times 8\) identity matrix. Unit:
  dimensionless. Status: STANDARD.
\item
  \(\omega_{\mathrm{b}}\): The energy of a state of dirac16complex bound
  to the fold at zero momentum: the four-dimensional mass of the
  particle \(\chi\). Unit: GeV. Status: DERIVED.
\item
  \(m_\Psi\): The mass parameter of dirac16complex in the bulk. Unit:
  GeV. Status: RECORD.
\end{itemize}

where \(\mu_{\mathrm{L}}\), \(\mu_{\mathrm{R}}\) are the strengths of
the fold for the left- and right-handed halves of dirac16complex (set by
the Standard-Model fields on the fold) and
\(q_{\mathrm{b}} = -m_\Psi(\mu_{\mathrm{L}} + \mu_{\mathrm{R}})\tanh\varsigma/(2\varsigma)\),
\(\varsigma = \sqrt{\mu_{\mathrm{L}}\mu_{\mathrm{R}}}\)
(\(q_{\mathrm{b}} = -m_\Psi\mu_{\mathrm{L}}/2\) for
\(\mu_{\mathrm{R}} = 0\)). The profile is not a function of
\(\lvert\sigma\rvert\): it jumps at the fold,
\(\chi(0^+) \ne \chi(0^-)\). 8. \emph{The fold as a mirror for
dirac16complex.} The reflection across the fold acts on the profiles
\(\chi(\sigma)\) in two ways,

\[
(\mathsf{R}_1\chi)(\sigma) = T_0T_1T_2T_3\,\chi(-\sigma),\qquad (\mathsf{R}_2\chi)(\sigma) = T_0T_5T_6T_7\,\chi(-\sigma) \tag{FD-10}
\] \textbf{Symbols in (FD-10).}

\begin{itemize}
\tightlist
\item
  \(\mathsf{R}_1\): The reflection across the fold acting on a profile
  of dirac16complex together with the reflection of ordinary space,
  \((\mathsf{R}_1\chi)(\sigma) = T_0T_1T_2T_3\,\chi(-\sigma)\) (FD-10):
  a symmetry of dirac16complex, of gravity and of the mean field of the
  gas of the particles \(\chi\); the bound states of Krein charge \(-1\)
  are even, those of Krein charge \(+1\) odd. Unit: dimensionless.
  Status: DERIVED.
\item
  \(\mathsf{R}_2\): The reflection across the fold acting on a profile
  of dirac16complex together with the reflection of the extra times,
  \((\mathsf{R}_2\chi)(\sigma) = T_0T_5T_6T_7\,\chi(-\sigma)\) (FD-10):
  it maps the bound states of one Krein charge onto those of the other.
  Unit: dimensionless. Status: DERIVED.
\item
  \(\chi(\sigma)\): The profile of a mode of dirac16complex along the
  hidden direction,
  \(\Psi = \mathrm{e}^{3H\lvert\sigma\rvert}\chi(\sigma)\mathrm{e}^{\mathrm{i}(\vec k\cdot\vec x - \omega x_4)}\)
  (a 16-component spinor function of \(\sigma\)); as a four-dimensional
  field, \(\chi\) denotes the particle of a bound state. Unit:
  GeV\(^{2}\) (with the \(V_*\) normalisation). Status: DERIVED.
\item
  \(\sigma\): The proper coordinate of the hidden direction:
  \(\mathrm{d}\sigma = \lvert\cot z\rvert\,\mathrm{d}x_0\),
  \(\sigma = \mathrm{sgn}(\pi/2 - z)\ln(\sin z)/(6H)\) on one period
  \(0 < z < \pi\); \(\sigma = 0\) is the fold \(z = \pi/2\),
  \(\sigma \to \mp\infty\) the tips \(z \to 0, \pi\); one period of
  \(x_0\) is the whole line \(-\infty < \sigma < \infty\). Unit: length.
  Status: DERIVED.
\item
  \(T_a,\ T^a\): The author's eight real \(16 \times 16\) Dirac matrices
  T16\^{}A{[}a{]} of 4+4.nb (\(a = 0, \dots, 7\); \(T_0, \dots, T_7\)),
  with \(\{T_a, T_b\} = 2\eta_{ab}\mathbf{1}\); \(T^a = \eta^{ab}T_b\).
  Unit: dimensionless. Status: AUTHOR.
\end{itemize}

(\(T_a\) the author's Dirac matrices; for states at rest
\(\mathsf{R}_1\) is the reflection across the fold combined with the
reflection of ordinary space, \(\mathsf{R}_2\) the reflection across the
fold combined with the reflection of the extra times). The Krein charge
of a state is its share of the conserved charge of dirac16complex, which
is indefinite: a state can carry the charge \(-1\) or \(+1\).
\(\mathsf{R}_1\) is a symmetry of dirac16complex, of gravity and of the
mean field of a gas of the particles \(\chi\) (the average field that
the particles of the gas create and feel): the four bound states of
Krein charge \(-1\) (the particles \(\chi\)) are even under it and the
four of Krein charge \(+1\) are odd, so the mean field of the gas cannot
mix the two signs of the indefinite charge. \(\mathsf{R}_2\) maps the
bound states of one Krein charge onto those of the other and makes the
two sets degenerate without the self-interaction. 9. \emph{The planes of
Horava and Witten.} In M-theory (the eleven-dimensional theory whose
limits include the superstring theories) on an interval (a circle whose
points are identified with their mirror images), the two ten-dimensional
planes left in place by the reflection are mirror surfaces (property 1);
each carries one \(E_8\) gauge group while gravity fills the bulk. As in
UM-3, the gauge fields live on the mirror surface and gravity in the
bulk; unlike UM-3, the two sides of each plane are identified (one
physical side).

\textbf{What is not a fold.}

\begin{itemize}
\tightlist
\item
  \emph{The tips} \(z = 0\) and \(z = \pi\) (mod \(\pi\)): the mirror
  exchanges them within a period, but there the warp factor vanishes,
  \(g_{x_0x_0}\) diverges, and they lie at an infinite proper distance;
  they carry no thin source.
\item
  \emph{A smooth maximum of a warp factor}, such as
  \(1/\cosh^2(H\sigma)\): mirror symmetric and largest at
  \(\sigma = 0\), but without a kink and without a thin source.
\end{itemize}

\hypertarget{how-to-read-this-document}{%
\section{How to read this document}\label{how-to-read-this-document}}

\hypertarget{who-this-document-is-for}{%
\subsection{Who this document is for}\label{who-this-document-is-for}}

It is written for a reader who is clever but new to the subject: someone
who knows what a coordinate, a derivative, an energy and a pressure are,
but not general relativity, warped geometries or the universe model
UM-3. Every term is explained where it is first used, every equation
that the document uses is written out in it, and every symbol of every
display equation is listed and defined right after the display. Nothing
has to be looked up anywhere else.

\hypertarget{the-status-labels}{%
\subsection{The status labels}\label{the-status-labels}}

Every statement carries one of the labels \textbf{AUTHOR} (from the
author, Patrick L. Nash: his notebooks
\texttt{Dirac\_Einstein-Lovelock\_4+4\_Nash.nb} and \texttt{4+4.nb} and
his orders of 2026-10-24 and 2026-10-25), \textbf{RECORD} (from the
earlier record of this repository, such as the Lagrangian of
dirac16complex), \textbf{STANDARD} (textbook physics or a published
measurement, named in the provenance), \textbf{ASSUMED} (a choice made
in UM-3 without derivation), \textbf{DERIVED} (computed and verified in
UM-3; the check that verifies it is named, e.g.~ST1), \textbf{ESTIMATE}
(an order-of-magnitude evaluation) and \textbf{OPEN} (not settled). A
result that holds only for the author's original Lovelock weights, which
UM-3 no longer uses, is called a record of the original weights.

\hypertarget{every-equation-is-followed-by-its-symbols}{%
\subsection{Every equation is followed by its
symbols}\label{every-equation-is-followed-by-its-symbols}}

A display equation carries a tag. SB-n is an equation of this document;
FD-n belongs to the definition of the fold at the beginning; FE-n (a
field equation) and IN-n (a term of the action) are equations of UM-3
that are written identically, with the same tag, in every UM-3 document
that uses them. After every display comes the list ``Symbols in
\ldots{}'' of every symbol in it, with its meaning, its unit, its value
where one exists, and its status. The units are natural units: the
reduced Planck constant and the speed of light are 1, so that a time is
a length, and energies, masses and inverse lengths are all measured in
GeV.

\hypertarget{the-checks}{%
\subsection{The checks}\label{the-checks}}

Most statements labelled DERIVED in this document are verified by the
nine checks ST1 to ST9 of a program that computes the field equations
exactly with computer algebra (ST1 to ST8 with the author's original
weights, ST9 with the corrected ones). The others are verified by checks
of the other programs of UM-3, named where they are used: the
gravitational waves and the corrected weights (W1 to W8), the curvature
and the bulk equations of the author's metric (G2 to G4), the junction
conditions and the fold sources (J4 to J7), and the proper distance of
the definition of the fold. The provenance at the end says what each
check does and how to run it.

\hypertarget{the-stabiliser-in-brief}{%
\section{The stabiliser in brief}\label{the-stabiliser-in-brief}}

UM-3 is a model universe built on a metric (a rule for measuring
distances and durations, explained below) that the author wrote down for
an eight-dimensional spacetime. The metric has a fold, the surface
defined at the beginning of this document, and around it the bulk:
everything that is not on the fold. Gravity in UM-3 obeys the
Einstein-Lovelock field equations, which say how the curvature of
spacetime is fixed by the energy that it contains. These equations
contain three weights and a cosmological constant.

\textbf{The finding (with the author's original weights).} When the
universe is old and expands slowly, as ours does, the bulk of the
author's metric is very nearly static: nothing in it changes with time.
With the weights and the cosmological constant of the author's notebook,
a static bulk of this metric is \textbf{not} a solution of the field
equations with empty space. The equations demanded a definite energy
content: a negative energy density, and a negative pressure five times
as large, the same in every direction of the bulk except the time
(SB-9). UM-3 called whatever supplies it \textbf{the stabiliser of the
bulk}, because without it the bulk could not stay static. It had no free
parameter: its energy density and its pressure were fixed by the
author's constant \(H\) and the gravitational coupling \(\kappa\).

\textbf{Who found it.} The requirement and the stabiliser are not part
of the author's specification. They were found, introduced and named by
Claude, an AI model made by Anthropic, working in Claude Code sessions
configured for the model claude-opus-5-5, on 2026-10-24 while it derived
the complete field equations of UM-3. The idea of something that holds a
static block of the bulk is older in this repository: the two earlier
universe models assumed a ``block stabiliser'' of their own (the
provenance gives the lineage); the stabiliser of UM-3, for the author's
new metric, was a different object with a definite, negative energy
density.

\textbf{The correction.} On 2026-10-25 the author ordered: ``Gravitons
must travel at c.~If these incorrect claims come from my choice of the
Lovelock couplings and my choice of the cosmological constant
\(\Lambda\) THEN my choices are wrong, and you MUST correct them.'' With
the original weights, gravitational waves in the bulk moved
\(\sqrt5 = 2.24\) times faster than light. There is exactly one choice
of the two higher weights for which gravitons travel at the speed of
light, and with it exactly one cosmological constant makes the static
bulk empty (FE-68 to FE-70). UM-3 now uses these values. \textbf{The
stabiliser was a consequence of the author's original weights and
cosmological constant; with the corrected ones the static bulk is a
vacuum, and the stabiliser vanishes} (IN-7, ST9).

\textbf{The five points of the author's order of 2026-10-25, in brief.}
1. \emph{The stabiliser couples only to gravity.} That is how UM-3
defined it, and on the static bulk it was consistent: the stabiliser was
conserved by itself (ST4). In an expanding universe the bulk source that
the original equations required released energy that flowed along the
hidden direction into the fold (ST6); reading this as an exchange of
energy between the stabiliser and the matter of the fold is an
interpretation (ASSUMED). With the corrected weights the question
disappears: there is no stabiliser, and the bulk of an expanding fold
with frozen extra times is a vacuum without any flow (FE-33). 2.
\emph{It is absent in the primordial era, whose bulk is the author's
vacuum (FE-31).} With the original weights this was true: the author's
vacuum solved the bulk equations with no source (ST2), and a medium that
couples only to gravity with a pressure five times its energy density
could not appear later out of nothing (ST6). With the corrected weights
the author's vacuum is not a solution any more, the primordial era is an
inflation of the fold, and no era of UM-3 has a stabiliser. 3. \emph{Its
existence and value are DERIVED from the requirement of a static bulk.}
True: the static bulk fixed every entry (ST1), and a bulk without any
source could only be the author's vacuum or its time reverse at every
instant (ST3, ST8). That the source was a separate medium was ASSUMED.
With the corrected weights the same derivation gives the value zero
(ST9). 4. \emph{The static bulk} is explained, with five examples, in
the part ``The static bulk''. 5. \emph{Warp and warp factor} are
defined, with four examples, in the part ``Warp and warp factor''.

\textbf{Its importance.} With the original weights the stabiliser was
what made a slowly expanding late universe with constant couplings
possible in the author's metric; it was also the most exotic ingredient
of the model, because it violated the null energy condition. Its removal
by the corrected weights changes UM-3 in four ways: gravitational waves
travel at the speed of light at every wavelength; the fold obeys the
Friedmann equations of general relativity exactly; the bulk needs no
exotic medium; and a new requirement appears on the fold, a pressure
along the extra times that holds their volume fixed, a role that the
``block stabiliser'' of the first universe model of this repository had
played for its own metric.

\hypertarget{the-language}{%
\section{The language}\label{the-language}}

\hypertarget{coordinates-distances-and-the-metric}{%
\subsection{Coordinates, distances and the
metric}\label{coordinates-distances-and-the-metric}}

A point of space and time is labelled by numbers, its coordinates.
Ordinary space needs three (\(x_1, x_2, x_3\)) and time one (\(x_4\)).
The spacetime of UM-3 has eight: \(x_0\), a hidden direction that we do
not see, \(x_1, x_2, x_3\) ordinary space, \(x_4\) the time, and
\(x_5, x_6, x_7\) three extra directions that behave like time (the
extra times). The metric is the rule that turns a small step of the
coordinates into the squared length (or the squared duration) of the
step. A familiar example is the surface of a globe of radius
\(R_{\mathrm{g}}\), on which a point is labelled by its angle
\(\theta_{\mathrm{g}}\) from the north pole and its longitude
\(\phi_{\mathrm{g}}\):

\[
\mathrm{d}\ell_{\mathrm{g}}^2 = R_{\mathrm{g}}^2\,\mathrm{d}\theta_{\mathrm{g}}^2 + R_{\mathrm{g}}^2\sin^2\theta_{\mathrm{g}}\;\mathrm{d}\phi_{\mathrm{g}}^2 \tag{SB-1}
\] \textbf{Symbols in (SB-1).}

\begin{itemize}
\tightlist
\item
  \(\mathrm{d}\): The differential: \(\mathrm{d}X\) is an infinitesimal
  change of \(X\), \(\mathrm{d}X/\mathrm{d}Y\) the derivative of \(X\)
  with respect to \(Y\), \(\mathrm{d}^nX\) the volume element of \(n\)
  variables. Status: STANDARD.
\item
  \(\ell_{\mathrm{g}}\): The length measured along the surface of a
  globe; \(\mathrm{d}\ell_{\mathrm{g}}\) is the length of a small step
  on it (an example of a metric, not UM-3). Unit: length. Status:
  STANDARD.
\item
  \(R_{\mathrm{g}}\): The radius of a globe (an example, not UM-3).
  Unit: length. Status: STANDARD.
\item
  \(\theta_{\mathrm{g}}\): The angle of a point of a globe from its
  north pole: 0 at the north pole, \(\pi/2\) on the equator, \(\pi\) at
  the south pole (an example, not UM-3). Unit: radian. Status: STANDARD.
\item
  \(\phi_{\mathrm{g}}\): The longitude of a point of a globe (an
  example, not UM-3). Unit: radian. Status: STANDARD.
\end{itemize}

A step of the longitude by \(\mathrm{d}\phi_{\mathrm{g}}\) is long on
the equator (\(\sin\theta_{\mathrm{g}} = 1\)) and short near the poles:
the factor \(R_{\mathrm{g}}^2\sin^2\theta_{\mathrm{g}}\) in front of
\(\mathrm{d}\phi_{\mathrm{g}}^2\) says how large one unit of longitude
is at each latitude (\(\sin\) is the sine function; \(\mathrm{d}\) marks
a small step). Status: STANDARD.

The metric of the eight-dimensional spacetime is written
\(g_{MN}\,\mathrm{d}X^M\mathrm{d}X^N\): the indices \(M\) and \(N\) run
over the eight coordinates \(X = \{x_0, \dots, x_7\}\), a repeated index
is summed, and the numbers \(g_{MN}\) are the components of the metric.
In spacetime a step can be space-like (a positive squared interval: a
distance) or time-like (a negative one: minus the square of a duration).
UM-3 counts \(x_0, x_1, x_2, x_3\) as space-like and
\(x_4, x_5, x_6, x_7\) as time-like (the signature \(+,+,+,+,-,-,-,-\)).
The author's metric is written in FD-4 of the definition of the fold
with the proper distance \(\sigma\) from the fold:
\(\mathrm{d}\sigma^2\) for the hidden direction, a factor that depends
on \(\sigma\) in front of ordinary space and of the extra times, and
\(-\mathrm{d}x_4^2\) for the time. The two functions \(a_4(x_4)\) and
\(b_4(x_4)\) of the time are the author's: at the depth
\(\lvert\sigma\rvert\), a coordinate step of ordinary space has the
length \(\mathrm{e}^{a_4}\mathrm{e}^{-H\lvert\sigma\rvert}\) times the
step, and a coordinate step of an extra time the duration
\(\mathrm{e}^{-a_4}b_4\,\mathrm{e}^{-H\lvert\sigma\rvert}\) times the
step; on the fold, \(\mathrm{e}^{a_4}\) is the size of ordinary space
and \(\mathrm{e}^{-a_4}b_4\) that of the extra times (\(\mathrm{e}\) is
Euler's number, \(\lvert\sigma\rvert\) the absolute value). Status:
AUTHOR (the metric), DERIVED (its form in \(\sigma\), FD-3 and FD-4).

\hypertarget{warp-and-warp-factor}{%
\subsection{Warp and warp factor}\label{warp-and-warp-factor}}

\textbf{Definition.} A metric is \emph{warped} along a direction when
the size of other directions depends on the position along it. In a
metric of the form
\(\mathrm{d}\sigma^2 + \mathcal{W}(\sigma)\,\mathrm{d}y^2 + \dots\) the
function \(\mathcal{W}(\sigma)\) that multiplies the squared steps
\(\mathrm{d}y^2\) of the other directions is the \textbf{warp factor}.
The length of a fixed coordinate step \(\mathrm{d}y\) is
\(\sqrt{\mathcal{W}}\,\mathrm{d}y\), so lengths along those directions
scale with \(\sqrt{\mathcal{W}}\). The \textbf{warp} is the dependence
itself, measured by the rate at which the natural logarithm
\(\ln\mathcal{W}\) changes along \(\sigma\). In UM-3:

\[
\mathcal{W}(\sigma) = \mathrm{e}^{-2H\lvert\sigma\rvert} = (\sin^2z)^{1/6},\qquad
\frac{\mathrm{d}\ln\mathcal{W}}{\mathrm{d}\sigma} = -2H\,\mathrm{sgn}(\sigma) \tag{SB-2}
\] \textbf{Symbols in (SB-2).}

\begin{itemize}
\tightlist
\item
  \(\mathcal{W}\): The warp factor: in a warped metric, the function of
  the position along one direction that multiplies the squared steps of
  other directions (its square root multiplies their lengths); in UM-3
  the function of \(\sigma\) that multiplies ordinary space and the
  extra times,
  \(\mathcal{W} = \mathrm{e}^{-2H\lvert\sigma\rvert} = (\sin^2z)^{1/6}\),
  equal to 1 on the fold. Unit: dimensionless. Status: DERIVED.
\item
  \(\mathrm{e}\): Euler's number, the base of the natural logarithm;
  \(\mathrm{e}^{X}\) is the exponential of \(X\). Unit: dimensionless.
  Value: 2.718281828. Status: STANDARD.
\item
  \(H\): The author's constant (an inverse length) that fixes the period
  \(\pi/(6H)\) of \(x_0\), the Lovelock couplings
  \(\alpha_{(p)} = w_{(p)}H^{2-2p}\) and the cosmological constant
  \(\Lambda = -7H^2\) (the author's original value \(-10H^2\) was
  corrected on 2026-10-25); in UM-3 a free parameter. Unit: GeV
  (1/length). Status: AUTHOR.
\item
  \(\sigma\): The proper coordinate of the hidden direction:
  \(\mathrm{d}\sigma = \lvert\cot z\rvert\,\mathrm{d}x_0\),
  \(\sigma = \mathrm{sgn}(\pi/2 - z)\ln(\sin z)/(6H)\) on one period
  \(0 < z < \pi\); \(\sigma = 0\) is the fold \(z = \pi/2\),
  \(\sigma \to \mp\infty\) the tips \(z \to 0, \pi\); one period of
  \(x_0\) is the whole line \(-\infty < \sigma < \infty\). Unit: length.
  Status: DERIVED.
\item
  \(z\): The angle variable of the hidden direction, \(z = 6Hx_0\); one
  period of \(x_0\) is \(0 < z < \pi\), the fold is \(z = \pi/2\). Unit:
  dimensionless. Status: AUTHOR.
\item
  \(\mathrm{d}\): The differential: \(\mathrm{d}X\) is an infinitesimal
  change of \(X\), \(\mathrm{d}X/\mathrm{d}Y\) the derivative of \(X\)
  with respect to \(Y\), \(\mathrm{d}^nX\) the volume element of \(n\)
  variables. Status: STANDARD.
\item
  \(\mathrm{sgn}\): The sign function: \(\mathrm{sgn}(y) = +1\) for
  \(y > 0\), \(-1\) for \(y < 0\). Status: STANDARD.
\end{itemize}

The warp factor multiplies ordinary space (together with
\(\mathrm{e}^{2a_4}\)) and the extra times (together with
\(\mathrm{e}^{-2a_4}b_4^2\)), but not the time \(x_4\) and not
\(\sigma\) itself. The warp factor is largest, 1, on the fold and falls
on both sides. A coordinate step of ordinary space at the depth
\(\lvert\sigma\rvert\) in the bulk is
\(\mathrm{e}^{-H\lvert\sigma\rvert}\) times as long as the same
coordinate step on the fold: the deeper in the bulk, the smaller
ordinary space and the extra times become. The warp factor is the same
on both sides of the fold; the warp, its logarithmic rate of change, has
the constant size \(2H\) and changes its sign on the fold
(\(\mathrm{sgn}\) is the sign function): this jump is the kink of the
fold (property 3 of the definition). Status: AUTHOR (the metric),
DERIVED (the form in \(\sigma\)).

\textbf{Examples of warps.}

\begin{itemize}
\tightlist
\item
  \emph{A globe} (SB-1): the longitude is warped along the angle from
  the pole, with the warp factor
  \(R_{\mathrm{g}}^2\sin^2\theta_{\mathrm{g}}\) in the proper coordinate
  \(R_{\mathrm{g}}\theta_{\mathrm{g}}\) (the distance from the pole
  along a meridian).
\item
  \emph{The fold of UM-3} (SB-2): ordinary space and the extra times are
  warped along the hidden direction.
\item
  \emph{The brane of Randall and Sundrum} (FD-8, example 6 of the
  definition of the fold): the four directions of our world are warped
  by the warp factor
  \(\mathrm{e}^{-2k_{\mathrm{RS}}\lvert y_{\mathrm{RS}}\rvert}\) along a
  fifth direction. Unlike in UM-3, the time is warped as well.
\item
  \emph{Time near the Earth.} Near a mass the time is warped. In the
  weak gravity of the Earth the factor in front of \(-\mathrm{d}x_4^2\)
  is \(1 + 2\Phi_{\mathrm{N}}\), where \(\Phi_{\mathrm{N}}\) is the
  Newtonian gravitational potential (negative, and larger in size closer
  to the Earth). A clock lower down therefore runs slower. Pound and
  Rebka measured this in 1960 over the 22.5 m height of a tower, and the
  clocks of the GPS satellites are corrected for it. In UM-3 the time is
  not warped along the hidden direction: a clock deep in the bulk runs
  at the same rate as a clock on the fold.
\end{itemize}

\hypertarget{the-bulk}{%
\subsection{The bulk}\label{the-bulk}}

\textbf{Definition.} The bulk is all of the eight-dimensional spacetime
except the fold: the points with \(\sigma \neq 0\). It extends on both
sides of the fold to the tips, which lie at an infinite proper distance.
Gravity and dirac16complex live in the bulk as well as on the fold;
dirac16complex is the author's field of matter of spin one half in eight
dimensions, a spinor with sixteen complex components. The Standard Model
of particle physics lives only on the fold. Status: ASSUMED (where the
fields live), AUTHOR (the metric and dirac16complex).

\hypertarget{energy-pressure-and-the-energy-momentum-tensor}{%
\subsection{Energy, pressure and the energy-momentum
tensor}\label{energy-pressure-and-the-energy-momentum-tensor}}

Matter and fields carry energy and momentum. Their densities and flows
are collected in the energy-momentum tensor, an eight-by-eight array
whose entries are measured along the eight directions of an orthonormal
frame: the frame direction \((0)\) along \(\sigma\), \((1), (2), (3)\)
along ordinary space, \((4)\) along the time and \((5), (6), (7)\) along
the extra times. For a medium at rest the entry along the time is minus
its energy density, the entries along the other directions are its
pressures, and an entry that mixes the time with another direction is a
flow of energy along that direction. A medium at rest with the same
pressure \(p_{\mathrm{med}}\) in every direction but the time has

\[
\mathcal{T}_{\mathrm{med}}{}^a{}_b = \mathrm{diag}\big(p_{\mathrm{med}}, p_{\mathrm{med}}, p_{\mathrm{med}}, p_{\mathrm{med}},
 -\rho_{\mathrm{med}}, p_{\mathrm{med}}, p_{\mathrm{med}}, p_{\mathrm{med}}\big) \tag{SB-3}
\] \textbf{Symbols in (SB-3).}

\begin{itemize}
\tightlist
\item
  \(\mathcal{T}_{\mathrm{med}}{}^a{}_b\): The energy-momentum tensor of
  one medium (frame components): for a medium at rest the entry along
  the time is minus its energy density, the entries along the other
  directions are its pressures, and the entries that mix the time with
  another direction are its flows of energy along that direction. Unit:
  GeV\(^8\) (energy per volume in eight dimensions). Status: STANDARD.
\item
  \(a\): A frame (orthonormal) index running over \(0, 1, \dots, 7\)
  (frame direction \((0)\) along \(\sigma\), \((1), (2), (3)\) along
  ordinary space, \((4)\) along the time, \((5), (6), (7)\) along the
  extra times); raised and lowered with the frame metric
  \(\eta_{ab} = \mathrm{diag}(1, 1, 1, 1, -1, -1, -1, -1)\); on the fold
  (\(a, b \in \{1, \dots, 7\}\)) the frame of the fold; a repeated index
  is summed. Status: STANDARD.
\item
  \(b\): A frame index (as \(a\)). Status: STANDARD.
\item
  \(\mathrm{diag}\): A diagonal matrix with the listed entries on its
  diagonal (in the order of the coordinates \(x_0, x_1, \dots, x_7\) for
  a tensor of eight dimensions, \(x_1, \dots, x_7\) for a tensor of the
  fold, \(x_1, \dots, x_4\) for a four-dimensional tensor). Status:
  STANDARD.
\item
  \(p_{\mathrm{med}}\): The pressure of a medium at rest, the same along
  every direction but the time (for a medium of the bulk: along
  \(\sigma\), ordinary space and the extra times). Unit: GeV\(^8\).
  Status: STANDARD.
\item
  \(\rho_{\mathrm{med}}\): The energy density (energy per volume) of a
  medium at rest. Unit: GeV\(^8\). Status: STANDARD.
\end{itemize}

Three examples in the four dimensions of our world: dust (pressureless
matter, \(p_{\mathrm{med}} = 0\)), light
(\(p_{\mathrm{med}} = \rho_{\mathrm{med}}/3\)) and the energy of the
vacuum (\(p_{\mathrm{med}} = -\rho_{\mathrm{med}}\)). A positive vacuum
energy is a tension: it pulls instead of pushing, like a stretched skin.
Yet its gravity is repulsive, because in general relativity pressure
gravitates too: the source of the attraction is \(\rho + 3p\) in four
dimensions, and a tension with \(p = -\rho\) makes it negative. This is
why a positive cosmological constant makes a universe expand faster. Two
energy conditions classify media: the weak energy condition (the energy
density is never negative) and the null energy condition (the energy
density plus the pressure along any space-like direction is never
negative; for the frame of UM-3, which has four time-like directions,
the condition is taken for light moving along the time \(x_4\) and one
space-like direction). Ordinary matter, light and a positive vacuum
energy satisfy both. A negative vacuum energy, such as the negative
cosmological constant of the bulk of UM-3, violates the weak condition
but not the null one, since its energy density plus its pressure is
zero. A medium that violates the null condition is called exotic.
Status: STANDARD.

\hypertarget{the-field-equations-of-gravity}{%
\subsection{The field equations of
gravity}\label{the-field-equations-of-gravity}}

Einstein's insight was that gravity is the curvature of spacetime, and
that the curvature is fixed by the energy-momentum tensor. One side of
his field equations is a combination of derivatives of the metric (the
Einstein tensor); the other is the energy-momentum tensor times the
gravitational coupling. UM-3 uses the generalisation found by Lovelock,
which in eight dimensions adds two more combinations of the curvature,
each with a weight. The action (the quantity whose variation gives the
equations) and its curvature terms are

\[
\begin{aligned}
S_{\mathrm{EL}} &= \frac{1}{2\kappa}\int\mathrm{d}^8X\,\sqrt{\lvert g\rvert}\,\Big(\sum_{p=1}^{3}\alpha_{(p)}\,\mathcal{L}_{(p)}
 - 2\Lambda\Big),\qquad \kappa = 8\pi G_8,\\
\alpha_{(p)} &= w_{(p)}\,H^{2-2p},\qquad w_{(1)} = 1,\quad w_{(2)} = \frac{1}{20},\quad w_{(3)} = \frac{1}{360},
 \qquad \Lambda = -7H^2
\end{aligned} \tag{IN-2}
\] \textbf{Symbols in (IN-2).}

\begin{itemize}
\tightlist
\item
  \(S_{\mathrm{EL}}\): The Einstein-Lovelock action of gravity in eight
  dimensions. Unit: dimensionless. Status: AUTHOR.
\item
  \(\kappa\): The eight-dimensional gravitational coupling,
  \(\kappa = 8\pi G_8\). Unit: GeV\(^{-6}\). Status: AUTHOR.
\item
  \(\mathrm{d}\): The differential: \(\mathrm{d}X\) is an infinitesimal
  change of \(X\), \(\mathrm{d}X/\mathrm{d}Y\) the derivative of \(X\)
  with respect to \(Y\), \(\mathrm{d}^nX\) the volume element of \(n\)
  variables. Status: STANDARD.
\item
  \(X = \{x_0, \dots, x_7\}\): The author's coordinates of the
  eight-dimensional spacetime, \(X = \{x_0, x_1, \dots, x_7\}\) with the
  hidden direction first (Dirac\_Einstein-Lovelock\_4+4\_Nash.nb,
  section DEFCHART); \(X^M\) is the coordinate with the index \(M\) and
  \(\mathrm{d}^8X = \mathrm{d}x_0\,\mathrm{d}x_1\cdots\mathrm{d}x_7\)
  the coordinate volume element. Unit: length (each coordinate). Status:
  AUTHOR.
\item
  \(g_{MN}\): The metric tensor of the eight-dimensional spacetime
  (signature \(+,+,+,+,-,-,-,-\) in the order \(x_0, \dots, x_7\)):
  \(g_{MN}\,\mathrm{d}X^M\mathrm{d}X^N\) is the squared interval,
  \(g^{MN}\) its inverse, \(\lvert g\rvert\) the absolute value of its
  determinant (\(\sqrt{\lvert g\rvert}\,\mathrm{d}^8X\) is the proper
  volume element). Unit: dimensionless. Status: AUTHOR.
\item
  \(p\): The Lovelock order \(p = 1, 2, 3\) (Einstein, Gauss-Bonnet,
  cubic), as an index and as a number; in the Kohn-Sham equations a
  second label of a Kohn-Sham orbital (as \(o\), with \(\delta_{op}\)
  the Kronecker delta of the two labels). Status: STANDARD.
\item
  \(\alpha_{(p)}\): The Lovelock coupling of order \(p\),
  \(\alpha_{(p)} = w_{(p)}H^{2-2p}\) (\(\alpha_{(1)} = 1\)). Unit:
  GeV\(^{2-2p}\). Status: AUTHOR.
\item
  \(\mathcal{L}_{(p)}\): The Lovelock density of order \(p\) (a scalar
  built from \(p\) Riemann tensors): \(\mathcal{L}_{(1)} = R\)
  (Einstein-Hilbert), \(\mathcal{L}_{(2)}\) the Gauss-Bonnet density,
  \(\mathcal{L}_{(3)}\) the cubic Lovelock density. Unit:
  1/length\(^{2p}\). Status: STANDARD.
\item
  \(\Lambda\): The eight-dimensional cosmological constant,
  \(\Lambda = -7H^2\): with the corrected weights the value for which
  the static block is a vacuum (its Lovelock tensors add to \(7H^2\) in
  every direction). The author's original value was \(\Lambda = -10H^2\)
  (he writes the pure number \(\Lambda_{\mathrm{N}} = 10\) in
  \(0 = -\Lambda_{\mathrm{N}} + \sum_p H^{-2p}w_{(p)}\mathcal{E}_{(p)}\)).
  Unit: GeV\(^2\). Value: \(-7H^2\); the author's original \(-10H^2\).
  Status: AUTHOR (the term), DERIVED (the value).
\item
  \(\pi\): The number pi, the ratio of the circumference of a circle to
  its diameter. Unit: dimensionless. Value: 3.141592654. Status:
  STANDARD.
\item
  \(G_8\): The eight-dimensional Newton constant
  (\(\kappa = 8\pi G_8\)). Unit: GeV\(^{-6}\). Status: AUTHOR.
\item
  \(w_{(p)}\): The Lovelock weights, \(w_{(1)} = 1\),
  \(w_{(2)} = 1/20\), \(w_{(3)} = 1/360\): the only weights for which
  gravitational waves travel at the speed of light in the bulk and on
  the fold (both coefficients of the bulk wave,
  \(1 - 40w_{(2)} + 360w_{(3)}\) and \(1 - 24w_{(2)} + 72w_{(3)}\),
  vanish). The author's original weights \(w_{(2)} = 1/32\),
  \(w_{(3)} = 1/1152\) (Dirac\_Einstein-Lovelock\_4+4\_Nash.nb,
  In{[}134{]}) made bulk gravitons faster than light; they were
  corrected at the author's order of 2026-10-25. Unit: dimensionless.
  Value: \((1, 1/20, 1/360)\); the author's original
  \((1, 1/32, 1/1152)\). Status: AUTHOR (the form), DERIVED (the
  values).
\item
  \(H\): The author's constant (an inverse length) that fixes the period
  \(\pi/(6H)\) of \(x_0\), the Lovelock couplings
  \(\alpha_{(p)} = w_{(p)}H^{2-2p}\) and the cosmological constant
  \(\Lambda = -7H^2\) (the author's original value \(-10H^2\) was
  corrected on 2026-10-25); in UM-3 a free parameter. Unit: GeV
  (1/length). Status: AUTHOR.
\end{itemize}

\[
\begin{aligned}
\mathcal{L}_{(p)} &= \frac{1}{2^{p}}\,\delta^{M_1N_1\cdots M_pN_p}_{P_1Q_1\cdots P_pQ_p}\,
 R^{P_1Q_1}{}_{M_1N_1}\cdots R^{P_pQ_p}{}_{M_pN_p},\\
\mathcal{L}_{(1)} &= R,\qquad \mathcal{L}_{(2)} = R^2 - 4R_{MN}R^{MN} + R_{MNPQ}R^{MNPQ}
\end{aligned} \tag{IN-3}
\] \textbf{Symbols in (IN-3).}

\begin{itemize}
\tightlist
\item
  \(\mathcal{L}_{(p)}\): The Lovelock density of order \(p\) (a scalar
  built from \(p\) Riemann tensors): \(\mathcal{L}_{(1)} = R\)
  (Einstein-Hilbert), \(\mathcal{L}_{(2)}\) the Gauss-Bonnet density,
  \(\mathcal{L}_{(3)}\) the cubic Lovelock density. Unit:
  1/length\(^{2p}\). Status: STANDARD.
\item
  \(p\): The Lovelock order \(p = 1, 2, 3\) (Einstein, Gauss-Bonnet,
  cubic), as an index and as a number; in the Kohn-Sham equations a
  second label of a Kohn-Sham orbital (as \(o\), with \(\delta_{op}\)
  the Kronecker delta of the two labels). Status: STANDARD.
\item
  \(\delta\): The Kronecker delta: with two indices \(\delta^M_N = 1\)
  if \(M = N\) and 0 otherwise (likewise \(\delta_{X\mathrm{L}}\) is 1
  if the label \(X\) of a half of dirac16complex is L and 0 otherwise);
  with upper and lower index strings the generalised Kronecker delta
  \(\delta^{M_1\cdots M_n}_{N_1\cdots N_n}\), the determinant of the
  \(n \times n\) matrix of the Kronecker deltas \(\delta^{M_a}_{N_b}\)
  (\(+1\) or \(-1\) if the upper indices are an even or odd permutation
  of the lower ones, 0 otherwise). Status: STANDARD.
\item
  \(M\): A coordinate index of eight dimensions (also with a numbered
  copy \(M_1, M_2, \dots\)): it runs over \(x_0, x_1, \dots, x_7\); a
  repeated index is summed. Status: STANDARD.
\item
  \(N\): A coordinate index of eight dimensions (also
  \(N_1, N_2, \dots\)), running over \(x_0, \dots, x_7\); a repeated
  index is summed. Status: STANDARD.
\item
  \(P\): A coordinate index of eight dimensions (also
  \(P_1, P_2, \dots\)), running over \(x_0, \dots, x_7\); a repeated
  index is summed. Status: STANDARD.
\item
  \(Q\): A coordinate index of eight dimensions (also
  \(Q_1, Q_2, \dots\)), running over \(x_0, \dots, x_7\); a repeated
  index is summed. Status: STANDARD.
\item
  \(R^{MN}{}_{PQ},\ R_{MN},\ R\): The curvature of the metric
  \(g_{MN}\): with four indices the Riemann tensor (convention
  \([\nabla_P, \nabla_Q]V^M = R^M{}_{NPQ}V^N\)), with two indices the
  Ricci tensor \(R_{MN} = R^P{}_{MPN}\), without indices the scalar
  curvature \(R = g^{MN}R_{MN}\); with frame indices
  \(R^{(a)(b)}{}_{(c)(d)}\) its frame components,
  \(R^{(a)(b)}{}_{(a)(b)}\) (no sum) the sectional curvature of the
  frame plane \((a, b)\). Unit: 1/length\(^2\). Status: STANDARD.
\end{itemize}

The curvature is measured by the Riemann tensor \(R\), which says how
much a direction turns when it is carried around a small loop; the
generalised Kronecker delta \(\delta\) is the totally antisymmetric
combination of ordinary Kronecker deltas, \(+1\) or \(-1\) when its
lower indices are an even or odd rearrangement of its upper ones and 0
otherwise. The weights in IN-2 are the corrected ones; the author's
notebook had \(w_{(2)} = 1/32\), \(w_{(3)} = 1/1152\) and
\(\Lambda = -10H^2\) (he writes the pure number 10 in his vacuum
equations). Varying the action with respect to the inverse metric
\(g^{MN}\) and dividing by \(\sqrt{\lvert g\rvert}\) gives the Lovelock
tensors

\[
\begin{aligned}
\mathcal{E}_{(p)}{}^M{}_N &= -\frac{1}{2^{p+1}}\,\delta^{MM_1N_1\cdots M_pN_p}_{NP_1Q_1\cdots P_pQ_p}\,
 R^{P_1Q_1}{}_{M_1N_1}\cdots R^{P_pQ_p}{}_{M_pN_p},\\
\mathcal{E}_{(1)}{}^M{}_N &= R^M{}_N - \frac12\,\delta^M_N\,R,\\
\mathcal{E}_{(2)}{}^M{}_N &= 2\big(R\,R^M{}_N - 2R^M{}_PR^P{}_N - 2R^{MP}{}_{NQ}R^Q{}_P + R^{MPQS}R_{NPQS}\big)
 - \frac12\,\delta^M_N\,\mathcal{L}_{(2)}
\end{aligned} \tag{FE-2}
\] \textbf{Symbols in (FE-2).}

\begin{itemize}
\tightlist
\item
  \(\mathcal{E}_{(p)}{}^M{}_N\): The Lovelock tensor of order \(p\): the
  variation of \(\sqrt{\lvert g\rvert}\,\mathcal{L}_{(p)}\) with respect
  to the metric (\(\mathcal{E}_{(1)}\) the Einstein tensor,
  \(\mathcal{E}_{(2)}\) the Gauss-Bonnet tensor); with frame indices its
  frame components. Unit: 1/length\(^{2p}\). Status: STANDARD.
\item
  \(p\): The Lovelock order \(p = 1, 2, 3\) (Einstein, Gauss-Bonnet,
  cubic), as an index and as a number; in the Kohn-Sham equations a
  second label of a Kohn-Sham orbital (as \(o\), with \(\delta_{op}\)
  the Kronecker delta of the two labels). Status: STANDARD.
\item
  \(M\): A coordinate index of eight dimensions (also with a numbered
  copy \(M_1, M_2, \dots\)): it runs over \(x_0, x_1, \dots, x_7\); a
  repeated index is summed. Status: STANDARD.
\item
  \(N\): A coordinate index of eight dimensions (also
  \(N_1, N_2, \dots\)), running over \(x_0, \dots, x_7\); a repeated
  index is summed. Status: STANDARD.
\item
  \(\delta\): The Kronecker delta: with two indices \(\delta^M_N = 1\)
  if \(M = N\) and 0 otherwise (likewise \(\delta_{X\mathrm{L}}\) is 1
  if the label \(X\) of a half of dirac16complex is L and 0 otherwise);
  with upper and lower index strings the generalised Kronecker delta
  \(\delta^{M_1\cdots M_n}_{N_1\cdots N_n}\), the determinant of the
  \(n \times n\) matrix of the Kronecker deltas \(\delta^{M_a}_{N_b}\)
  (\(+1\) or \(-1\) if the upper indices are an even or odd permutation
  of the lower ones, 0 otherwise). Status: STANDARD.
\item
  \(P\): A coordinate index of eight dimensions (also
  \(P_1, P_2, \dots\)), running over \(x_0, \dots, x_7\); a repeated
  index is summed. Status: STANDARD.
\item
  \(Q\): A coordinate index of eight dimensions (also
  \(Q_1, Q_2, \dots\)), running over \(x_0, \dots, x_7\); a repeated
  index is summed. Status: STANDARD.
\item
  \(R^{MN}{}_{PQ},\ R_{MN},\ R\): The curvature of the metric
  \(g_{MN}\): with four indices the Riemann tensor (convention
  \([\nabla_P, \nabla_Q]V^M = R^M{}_{NPQ}V^N\)), with two indices the
  Ricci tensor \(R_{MN} = R^P{}_{MPN}\), without indices the scalar
  curvature \(R = g^{MN}R_{MN}\); with frame indices
  \(R^{(a)(b)}{}_{(c)(d)}\) its frame components,
  \(R^{(a)(b)}{}_{(a)(b)}\) (no sum) the sectional curvature of the
  frame plane \((a, b)\). Unit: 1/length\(^2\). Status: STANDARD.
\item
  \(S\): A coordinate index of eight dimensions, running over
  \(x_0, \dots, x_7\); a repeated index is summed. Status: STANDARD.
\item
  \(\mathcal{L}_{(p)}\): The Lovelock density of order \(p\) (a scalar
  built from \(p\) Riemann tensors): \(\mathcal{L}_{(1)} = R\)
  (Einstein-Hilbert), \(\mathcal{L}_{(2)}\) the Gauss-Bonnet density,
  \(\mathcal{L}_{(3)}\) the cubic Lovelock density. Unit:
  1/length\(^{2p}\). Status: STANDARD.
\end{itemize}

(\(R^M{}_N\) the Ricci tensor and \(R\) the scalar curvature,
contractions of the Riemann tensor), and the field equations of the bulk
are

\[
\sum_{p=1}^{3}\alpha_{(p)}\,\mathcal{E}_{(p)}{}^M{}_N + \Lambda\,\delta^M_N
 = \kappa\,\mathcal{T}_\Psi{}^M{}_N\qquad (\sigma \ne 0) \tag{FE-1}
\] \textbf{Symbols in (FE-1).}

\begin{itemize}
\tightlist
\item
  \(p\): The Lovelock order \(p = 1, 2, 3\) (Einstein, Gauss-Bonnet,
  cubic), as an index and as a number; in the Kohn-Sham equations a
  second label of a Kohn-Sham orbital (as \(o\), with \(\delta_{op}\)
  the Kronecker delta of the two labels). Status: STANDARD.
\item
  \(\alpha_{(p)}\): The Lovelock coupling of order \(p\),
  \(\alpha_{(p)} = w_{(p)}H^{2-2p}\) (\(\alpha_{(1)} = 1\)). Unit:
  GeV\(^{2-2p}\). Status: AUTHOR.
\item
  \(\mathcal{E}_{(p)}{}^M{}_N\): The Lovelock tensor of order \(p\): the
  variation of \(\sqrt{\lvert g\rvert}\,\mathcal{L}_{(p)}\) with respect
  to the metric (\(\mathcal{E}_{(1)}\) the Einstein tensor,
  \(\mathcal{E}_{(2)}\) the Gauss-Bonnet tensor); with frame indices its
  frame components. Unit: 1/length\(^{2p}\). Status: STANDARD.
\item
  \(M\): A coordinate index of eight dimensions (also with a numbered
  copy \(M_1, M_2, \dots\)): it runs over \(x_0, x_1, \dots, x_7\); a
  repeated index is summed. Status: STANDARD.
\item
  \(N\): A coordinate index of eight dimensions (also
  \(N_1, N_2, \dots\)), running over \(x_0, \dots, x_7\); a repeated
  index is summed. Status: STANDARD.
\item
  \(\Lambda\): The eight-dimensional cosmological constant,
  \(\Lambda = -7H^2\): with the corrected weights the value for which
  the static block is a vacuum (its Lovelock tensors add to \(7H^2\) in
  every direction). The author's original value was \(\Lambda = -10H^2\)
  (he writes the pure number \(\Lambda_{\mathrm{N}} = 10\) in
  \(0 = -\Lambda_{\mathrm{N}} + \sum_p H^{-2p}w_{(p)}\mathcal{E}_{(p)}\)).
  Unit: GeV\(^2\). Value: \(-7H^2\); the author's original \(-10H^2\).
  Status: AUTHOR (the term), DERIVED (the value).
\item
  \(\delta\): The Kronecker delta: with two indices \(\delta^M_N = 1\)
  if \(M = N\) and 0 otherwise (likewise \(\delta_{X\mathrm{L}}\) is 1
  if the label \(X\) of a half of dirac16complex is L and 0 otherwise);
  with upper and lower index strings the generalised Kronecker delta
  \(\delta^{M_1\cdots M_n}_{N_1\cdots N_n}\), the determinant of the
  \(n \times n\) matrix of the Kronecker deltas \(\delta^{M_a}_{N_b}\)
  (\(+1\) or \(-1\) if the upper indices are an even or odd permutation
  of the lower ones, 0 otherwise). Status: STANDARD.
\item
  \(\kappa\): The eight-dimensional gravitational coupling,
  \(\kappa = 8\pi G_8\). Unit: GeV\(^{-6}\). Status: AUTHOR.
\item
  \(\mathcal{T}_\Psi\): The energy-momentum tensor of dirac16complex,
  \(\mathcal{T}_{MN} = -(2/\sqrt{\lvert g\rvert})\,\delta S_\Psi/\delta g^{MN}\).
  Unit: GeV\(^8\). Status: STANDARD.
\item
  \(\sigma\): The proper coordinate of the hidden direction:
  \(\mathrm{d}\sigma = \lvert\cot z\rvert\,\mathrm{d}x_0\),
  \(\sigma = \mathrm{sgn}(\pi/2 - z)\ln(\sin z)/(6H)\) on one period
  \(0 < z < \pi\); \(\sigma = 0\) is the fold \(z = \pi/2\),
  \(\sigma \to \mp\infty\) the tips \(z \to 0, \pi\); one period of
  \(x_0\) is the whole line \(-\infty < \sigma < \infty\). Unit: length.
  Status: DERIVED.
\end{itemize}

The left-hand side is built from the metric alone. The right-hand side
is the energy content of the bulk: dirac16complex, and, with the
original weights, the stabiliser. A metric solves the equations only if
the right-hand side has exactly the value that the left-hand side
computes; a solution with nothing on the right is a \emph{vacuum}.
Status: AUTHOR (the theory), STANDARD (Lovelock gravity), DERIVED (the
corrected weights).

The divergence of every Lovelock tensor vanishes identically (the
Bianchi identity). The right-hand side must therefore be conserved: the
energy and momentum of all sources together cannot be created or
destroyed, only moved. A single medium that interacts with nothing but
gravity is conserved by itself:

\[
\nabla_M\,\mathcal{T}_{\mathrm{med}}{}^M{}_N = 0 \tag{SB-4}
\] \textbf{Symbols in (SB-4).}

\begin{itemize}
\tightlist
\item
  \(\nabla\): The covariant derivative of a tensor with the Levi-Civita
  connection of the metric \(g_{MN}\) (Christoffel symbols
  \(\Gamma^P{}_{MN}\)). Status: STANDARD.
\item
  \(\mathcal{T}_{\mathrm{med}}{}^a{}_b\): The energy-momentum tensor of
  one medium (frame components): for a medium at rest the entry along
  the time is minus its energy density, the entries along the other
  directions are its pressures, and the entries that mix the time with
  another direction are its flows of energy along that direction. Unit:
  GeV\(^8\) (energy per volume in eight dimensions). Status: STANDARD.
\item
  \(M\): A coordinate index of eight dimensions (also with a numbered
  copy \(M_1, M_2, \dots\)): it runs over \(x_0, x_1, \dots, x_7\); a
  repeated index is summed. Status: STANDARD.
\item
  \(N\): A coordinate index of eight dimensions (also
  \(N_1, N_2, \dots\)), running over \(x_0, \dots, x_7\); a repeated
  index is summed. Status: STANDARD.
\end{itemize}

(\(\nabla\) is the covariant derivative, the derivative that takes into
account how the directions of the coordinates change from point to
point.) Status: STANDARD.

\textbf{The junction conditions.} The metric has a kink at the fold, so
its second derivatives along \(\sigma\) contain a Dirac delta function
there, and the field equations hold across the fold only if the fold
carries a thin source \(\mathcal{S}\), its energy-momentum tensor
integrated across it (unit GeV\(^7\): an energy density times a length):
the integral FD-11 of the definition of the fold. For the Einstein term
alone this is Israel's condition (1966):
\(\kappa\mathcal{S}^a{}_b = -([K^a{}_b] - \delta^a{}_b[K])\), with the
jump \([K^a{}_b]\) of the extrinsic curvature of the fold (FD-5) and its
trace \([K]\); for the Gauss-Bonnet and cubic terms the integral adds
terms of higher order in the extrinsic curvature and the curvature of
the fold (Davis 2003; Gravanis and Willison 2003). Status: STANDARD.

\hypertarget{the-static-bulk}{%
\section{The static bulk}\label{the-static-bulk}}

\hypertarget{what-static-means}{%
\subsection{What static means}\label{what-static-means}}

\textbf{Definition.} A spacetime is \emph{static} when nothing in it
changes with the time and nothing flows: its metric does not depend on
the time, and no motion that the reversal of time would change takes
place in it. In the author's metric the time dependence sits in the two
functions \(a_4(x_4)\) and \(b_4(x_4)\). The bulk is static when both
are constant, that is when the rates \(\mathsf{A} = a_4'/H\) (the rate
at which ordinary space grows) and \(\mathsf{U} = b_4'/(Hb_4)\) (the
rate at which \(b_4\) changes) both vanish:

\[
\begin{aligned}
&\mathsf{A} = \frac{a_4'}{H} = 0,\qquad \mathsf{U} = \frac{b_4'}{Hb_4} = 0:\\
&g_{MN}\,\mathrm{d}X^M\mathrm{d}X^N = \mathrm{d}\sigma^2 + \mathrm{e}^{-2H\lvert\sigma\rvert}\Big[\mathrm{e}^{2a_4}\big(\mathrm{d}x_1^2
 + \mathrm{d}x_2^2 + \mathrm{d}x_3^2\big) - \mathrm{e}^{-2a_4}b_4^2\big(\mathrm{d}x_5^2 + \mathrm{d}x_6^2 + \mathrm{d}x_7^2\big)\Big]
 - \mathrm{d}x_4^2
\end{aligned} \tag{SB-5}
\] \textbf{Symbols in (SB-5).}

\begin{itemize}
\tightlist
\item
  \(\mathsf{A}\): The expansion rate of ordinary space in units of
  \(H\), \(\mathsf{A} = a_4'/H\). Unit: dimensionless. Value: 0 at the
  static block; \(H_4^{(0)}/H\) today (\(H_4^{(0)}\) the Hubble
  constant); 1 at the author's vacuum, which is not a solution with the
  corrected weights.
\item
  \(a_4\): The author's metric function of the time \(x_4\) that scales
  ordinary 3-space: on the fold \(\mathrm{e}^{a_4}\) is the scale factor
  of cosmology and \(a_4'\) the Hubble rate. Unit: dimensionless.
  Status: AUTHOR.
\item
  \({}'\): The derivative of a function of the time \(x_4\) alone with
  respect to \(x_4\) (\(a_4' = \mathrm{d}a_4/\mathrm{d}x_4\); two primes
  the second derivative); on a function of \(\sigma\) alone (a profile)
  the derivative with respect to \(\sigma\). Status: STANDARD.
\item
  \(H\): The author's constant (an inverse length) that fixes the period
  \(\pi/(6H)\) of \(x_0\), the Lovelock couplings
  \(\alpha_{(p)} = w_{(p)}H^{2-2p}\) and the cosmological constant
  \(\Lambda = -7H^2\) (the author's original value \(-10H^2\) was
  corrected on 2026-10-25); in UM-3 a free parameter. Unit: GeV
  (1/length). Status: AUTHOR.
\item
  \(\mathsf{U}\): The rate of change of \(b_4\) in units of \(H\),
  \(\mathsf{U} = b_4'/(Hb_4)\). Unit: dimensionless. Value:
  \(\mathsf{U} = \mathsf{A}\) with frozen extra times.
\item
  \(b_4\): The author's metric function of the time \(x_4\) in the
  extra-time components, \(b_4 > 0\), never identically 1 (the author's
  order of 2026-10-24); the scale factor (the factor of the lengths) of
  the extra times is \(\mathrm{e}^{-a_4}b_4\) on the fold and
  \(\mathrm{e}^{-a_4}b_4\,\mathrm{e}^{-H\lvert\sigma\rvert}\) at the
  distance \(\sigma\) from it (the square root of the warp factor times
  \(\mathrm{e}^{-a_4}b_4\)). Unit: dimensionless. Status: AUTHOR.
\item
  \(g_{MN}\): The metric tensor of the eight-dimensional spacetime
  (signature \(+,+,+,+,-,-,-,-\) in the order \(x_0, \dots, x_7\)):
  \(g_{MN}\,\mathrm{d}X^M\mathrm{d}X^N\) is the squared interval,
  \(g^{MN}\) its inverse, \(\lvert g\rvert\) the absolute value of its
  determinant (\(\sqrt{\lvert g\rvert}\,\mathrm{d}^8X\) is the proper
  volume element). Unit: dimensionless. Status: AUTHOR.
\item
  \(M\): A coordinate index of eight dimensions (also with a numbered
  copy \(M_1, M_2, \dots\)): it runs over \(x_0, x_1, \dots, x_7\); a
  repeated index is summed. Status: STANDARD.
\item
  \(N\): A coordinate index of eight dimensions (also
  \(N_1, N_2, \dots\)), running over \(x_0, \dots, x_7\); a repeated
  index is summed. Status: STANDARD.
\item
  \(\mathrm{d}\): The differential: \(\mathrm{d}X\) is an infinitesimal
  change of \(X\), \(\mathrm{d}X/\mathrm{d}Y\) the derivative of \(X\)
  with respect to \(Y\), \(\mathrm{d}^nX\) the volume element of \(n\)
  variables. Status: STANDARD.
\item
  \(X = \{x_0, \dots, x_7\}\): The author's coordinates of the
  eight-dimensional spacetime, \(X = \{x_0, x_1, \dots, x_7\}\) with the
  hidden direction first (Dirac\_Einstein-Lovelock\_4+4\_Nash.nb,
  section DEFCHART); \(X^M\) is the coordinate with the index \(M\) and
  \(\mathrm{d}^8X = \mathrm{d}x_0\,\mathrm{d}x_1\cdots\mathrm{d}x_7\)
  the coordinate volume element. Unit: length (each coordinate). Status:
  AUTHOR.
\item
  \(\sigma\): The proper coordinate of the hidden direction:
  \(\mathrm{d}\sigma = \lvert\cot z\rvert\,\mathrm{d}x_0\),
  \(\sigma = \mathrm{sgn}(\pi/2 - z)\ln(\sin z)/(6H)\) on one period
  \(0 < z < \pi\); \(\sigma = 0\) is the fold \(z = \pi/2\),
  \(\sigma \to \mp\infty\) the tips \(z \to 0, \pi\); one period of
  \(x_0\) is the whole line \(-\infty < \sigma < \infty\). Unit: length.
  Status: DERIVED.
\item
  \(\mathrm{e}\): Euler's number, the base of the natural logarithm;
  \(\mathrm{e}^{X}\) is the exponential of \(X\). Unit: dimensionless.
  Value: 2.718281828. Status: STANDARD.
\item
  \(x_1\): The first coordinate of ordinary 3-space. Unit: length.
  Status: AUTHOR.
\item
  \(x_2\): The second coordinate of ordinary 3-space. Unit: length.
  Status: AUTHOR.
\item
  \(x_3\): The third coordinate of ordinary 3-space. Unit: length.
  Status: AUTHOR.
\item
  \(x_4\): The time; a prime on a function of \(x_4\) alone is
  \(\mathrm{d}/\mathrm{d}x_4\). Unit: length (= time, \(c = 1\)).
  Status: AUTHOR.
\item
  \(x_5\): The first of the three extra time-like coordinates. Unit:
  length. Status: AUTHOR.
\item
  \(x_6\): The second of the three extra time-like coordinates. Unit:
  length. Status: AUTHOR.
\item
  \(x_7\): The third of the three extra time-like coordinates. Unit:
  length. Status: AUTHOR.
\end{itemize}

with constant \(a_4\) and \(b_4\) (a prime is the derivative with
respect to the time \(x_4\)). UM-3 calls this the \emph{static block}.
Status: DERIVED (the form), ASSUMED (that the late bulk is the static
block to first approximation).

\hypertarget{static-is-not-the-same-as-empty}{%
\subsection{Static is not the same as
empty}\label{static-is-not-the-same-as-empty}}

A static spacetime may be empty or it may need a source, and an empty
spacetime may be static or not:

\begin{longtable}[]{@{}
  >{\raggedright\arraybackslash}p{(\columnwidth - 4\tabcolsep) * \real{0.3333}}
  >{\raggedright\arraybackslash}p{(\columnwidth - 4\tabcolsep) * \real{0.3333}}
  >{\raggedright\arraybackslash}p{(\columnwidth - 4\tabcolsep) * \real{0.3333}}@{}}
\toprule\noalign{}
\begin{minipage}[b]{\linewidth}\raggedright
spacetime
\end{minipage} & \begin{minipage}[b]{\linewidth}\raggedright
static?
\end{minipage} & \begin{minipage}[b]{\linewidth}\raggedright
empty (a vacuum)?
\end{minipage} \\
\midrule\noalign{}
\endhead
\bottomrule\noalign{}
\endlastfoot
the static block of UM-3 (SB-5), corrected weights & yes & yes (IN-7) \\
the static block of UM-3 (SB-5), the author's original weights & yes &
no: it needed the stabiliser (SB-9) \\
the author's vacuum (FE-31) & no: ordinary space grows at the rate \(H\)
& with the original weights yes; with the corrected weights no \\
the space outside a star & yes & yes \\
Einstein's static universe (SB-6) & yes & no: it needs dust and a
cosmological constant \\
\end{longtable}

\hypertarget{examples-of-static-spacetimes-and-static-media}{%
\subsection{Examples of static spacetimes and static
media}\label{examples-of-static-spacetimes-and-static-media}}

\textbf{Example 1: the static block of UM-3} (SB-5). It is the late
universe to first approximation. Today ordinary space expands at the
Hubble rate \(H_4^{(0)} = 1.44\times10^{-42}\) GeV. If the author's
constant \(H\) is much larger, the rate \(\mathsf{A} = H_4^{(0)}/H\) is
tiny and the bulk is static up to corrections of order \(\mathsf{A}\).
The extra times are frozen (\(\mathrm{e}^{-a_4}b_4\) constant,
\(\mathsf{U} = \mathsf{A}\)), so that the couplings of the Standard
Model do not change. With the corrected weights the static block is an
exact vacuum, and so is the bulk of an expanding fold with frozen extra
times for every expansion history (FE-33). Status: ASSUMED (the late
universe), DERIVED (the vacuum).

\textbf{Example 2: the space outside a star or a planet}
(Schwarzschild's solution of 1916). It is static and empty: a static
vacuum. Not every static spacetime needs a source. Status: STANDARD.

\textbf{Example 3: Einstein's static universe} (1917). Einstein wanted a
universe that does not change. With dust of the mass density
\(\rho_{\mathrm{E}}\) and nothing else his equations do not allow it,
because gravity pulls the dust together. He added a positive
cosmological constant \(\Lambda_{\mathrm{E}}\), whose gravity is
repulsive (it acts like a positive vacuum energy, a tension), and found
a static universe, a three-dimensional sphere of radius
\(R_{\mathrm{E}}\), but only for one exact balance:

\[
\Lambda_{\mathrm{E}} = 4\pi G_{\mathrm{N}}\,\rho_{\mathrm{E}} = \frac{1}{R_{\mathrm{E}}^2} \tag{SB-6}
\] \textbf{Symbols in (SB-6).}

\begin{itemize}
\tightlist
\item
  \(\Lambda_{\mathrm{E}}\): The cosmological constant of Einstein's
  static universe of 1917, in four dimensions (an example, not UM-3).
  Unit: 1/length\(^2\). Status: STANDARD.
\item
  \(\pi\): The number pi, the ratio of the circumference of a circle to
  its diameter. Unit: dimensionless. Value: 3.141592654. Status:
  STANDARD.
\item
  \(G_{\mathrm{N}}\): Newton's constant of four-dimensional gravity (the
  gravity of our world, on the fold),
  \(G_{\mathrm{N}} = 1/M_{\mathrm{Pl}}^2\) with the Planck mass
  \(M_{\mathrm{Pl}} = 1.221\times10^{19}\) GeV. Unit: GeV\(^{-2}\).
  Value: \(6.708\times10^{-39}\) GeV\(^{-2}\). Status: STANDARD.
\item
  \(\rho_{\mathrm{E}}\): The mass density of the dust that fills
  Einstein's static universe (an example, not UM-3). Unit: mass per
  volume. Status: STANDARD.
\item
  \(R_{\mathrm{E}}\): The radius of the three-dimensional sphere that is
  the space of Einstein's static universe (an example, not UM-3). Unit:
  length. Status: STANDARD.
\end{itemize}

This is the classic example of a static spacetime that requires a
definite source. It is also unstable (Eddington, 1930): a small excess
of matter makes it collapse, a small deficit makes it expand for ever.
Einstein's \(\Lambda_{\mathrm{E}}\) played for his universe the part
that the stabiliser played for the bulk of UM-3 with the original
weights. Status: STANDARD.

\textbf{Example 4: the air at rest above the ground.} The atmosphere
stays at rest only because its pressure falls with the altitude exactly
so fast that the difference of pressure carries the weight of the air in
between:

\[
\frac{\mathrm{d}P_{\mathrm{air}}}{\mathrm{d}h_{\mathrm{alt}}} = -\rho_{\mathrm{air}}\,g_{\oplus} \tag{SB-7}
\] \textbf{Symbols in (SB-7).}

\begin{itemize}
\tightlist
\item
  \(\mathrm{d}\): The differential: \(\mathrm{d}X\) is an infinitesimal
  change of \(X\), \(\mathrm{d}X/\mathrm{d}Y\) the derivative of \(X\)
  with respect to \(Y\), \(\mathrm{d}^nX\) the volume element of \(n\)
  variables. Status: STANDARD.
\item
  \(P_{\mathrm{air}}\): The pressure of the air at the altitude
  \(h_{\mathrm{alt}}\) (an example, not UM-3). Unit: Pa. Value: about
  \(1.0\times10^5\) Pa at the ground. Status: STANDARD.
\item
  \(h_{\mathrm{alt}}\): The altitude above the ground (an example, not
  UM-3). Unit: m. Status: STANDARD.
\item
  \(\rho_{\mathrm{air}}\): The mass density of the air at the altitude
  \(h_{\mathrm{alt}}\) (an example, not UM-3). Unit: kg/m\(^3\). Value:
  about 1.2 kg/m\(^3\) at the ground. Status: STANDARD.
\item
  \(g_{\oplus}\): The acceleration of free fall at the surface of the
  Earth (an example, not UM-3). Unit: m/s\(^2\). Value: 9.81 m/s\(^2\).
  Status: STANDARD.
\end{itemize}

Again, the static state needs a definite distribution of pressure; any
other distribution sets the air in motion. Status: STANDARD.

\textbf{Example 5: a drum skin at rest.} A stretched skin stays flat and
at rest because its tension is the same everywhere, so that the pulls on
every small piece of it balance. Pressing it or loosening it locally
sets it in motion. The fold of UM-3 is a skin of this kind: at the
static block its thin source is a pure tension (FE-32). Status:
STANDARD.

In all these examples except the empty space outside a star, staying
static requires a definite content: a definite balance of energy and
pressure. Whether the static block of UM-3 is of this kind or a vacuum
depends on the Lovelock weights.

\hypertarget{the-stabiliser-of-the-authors-original-weights}{%
\section{The stabiliser of the author's original
weights}\label{the-stabiliser-of-the-authors-original-weights}}

This part is the record of the stabiliser as it was found, with the
weights \(w_{(2)} = 1/32\), \(w_{(3)} = 1/1152\) and the cosmological
constant \(\Lambda = -10H^2\) of the author's notebook. The next part
shows what the corrected weights change.

\hypertarget{how-the-static-block-fixes-the-source}{%
\subsection{How the static block fixes the
source}\label{how-the-static-block-fixes-the-source}}

The static block SB-5 has a simple curvature. In each half of the bulk,
the seven directions other than the time (\(\sigma\), ordinary space and
the extra times) form a space of constant negative curvature \(-H^2\),
like a slice of anti-de Sitter space: every pair of these directions has
the frame curvature \(R^{ab}{}_{ab} = -H^2\), and no pair with the time
is curved (the time is not warped). For such a space times a flat line,
the Lovelock tensor of order \(p\) (FE-2) has only diagonal entries, and
the entry along a direction depends only on the number
\(m_{\mathrm{c}}\) of the other curved directions:

\[
\begin{aligned}
&\mathcal{E}_{(p)}{}^{(a)}{}_{(a)}\Big\vert_{\text{static block}} = -\frac12\,\big(-H^2\big)^p\,\frac{m_{\mathrm{c}}!}{(m_{\mathrm{c}} - 2p)!}\\
&\big(m_{\mathrm{c}} = 6\ \text{along the hidden direction, ordinary space and the extra times},\ m_{\mathrm{c}} = 7\ \text{along the time}\big)
\end{aligned} \tag{SB-8}
\] \textbf{Symbols in (SB-8).}

\begin{itemize}
\tightlist
\item
  \(\mathcal{E}_{(p)}{}^M{}_N\): The Lovelock tensor of order \(p\): the
  variation of \(\sqrt{\lvert g\rvert}\,\mathcal{L}_{(p)}\) with respect
  to the metric (\(\mathcal{E}_{(1)}\) the Einstein tensor,
  \(\mathcal{E}_{(2)}\) the Gauss-Bonnet tensor); with frame indices its
  frame components. Unit: 1/length\(^{2p}\). Status: STANDARD.
\item
  \(a\): A frame (orthonormal) index running over \(0, 1, \dots, 7\)
  (frame direction \((0)\) along \(\sigma\), \((1), (2), (3)\) along
  ordinary space, \((4)\) along the time, \((5), (6), (7)\) along the
  extra times); raised and lowered with the frame metric
  \(\eta_{ab} = \mathrm{diag}(1, 1, 1, 1, -1, -1, -1, -1)\); on the fold
  (\(a, b \in \{1, \dots, 7\}\)) the frame of the fold; a repeated index
  is summed. Status: STANDARD.
\item
  \(p\): The Lovelock order \(p = 1, 2, 3\) (Einstein, Gauss-Bonnet,
  cubic), as an index and as a number; in the Kohn-Sham equations a
  second label of a Kohn-Sham orbital (as \(o\), with \(\delta_{op}\)
  the Kronecker delta of the two labels). Status: STANDARD.
\item
  \(H\): The author's constant (an inverse length) that fixes the period
  \(\pi/(6H)\) of \(x_0\), the Lovelock couplings
  \(\alpha_{(p)} = w_{(p)}H^{2-2p}\) and the cosmological constant
  \(\Lambda = -7H^2\) (the author's original value \(-10H^2\) was
  corrected on 2026-10-25); in UM-3 a free parameter. Unit: GeV
  (1/length). Status: AUTHOR.
\item
  \(m_{\mathrm{c}}\): The number of the other curved directions of the
  static block seen from one direction: 6 for \(\sigma\), ordinary space
  and the extra times (the seven curved directions without the given
  one), 7 for the time (which is not curved itself). Unit:
  dimensionless. Value: 6 or 7. Status: DERIVED.
\end{itemize}

(no sum over the frame index \(a\); \(m_{\mathrm{c}}!\) is the factorial
\(1\cdot2\cdots m_{\mathrm{c}}\).) Each Riemann tensor in FE-2
contributes the curvature \(-H^2\) of a pair of directions, and the
generalised Kronecker delta counts the ordered ways of choosing the
\(2p\) other directions of the \(p\) pairs among the \(m_{\mathrm{c}}\)
curved ones, \(m_{\mathrm{c}}!/(m_{\mathrm{c}} - 2p)!\). Multiplied by
the couplings \(\alpha_{(p)} = w_{(p)}H^{2-2p}\), the entries along a
warped direction and along the time are, in units of \(H^2\):

\begin{longtable}[]{@{}
  >{\raggedright\arraybackslash}p{(\columnwidth - 6\tabcolsep) * \real{0.2500}}
  >{\raggedright\arraybackslash}p{(\columnwidth - 6\tabcolsep) * \real{0.2500}}
  >{\raggedright\arraybackslash}p{(\columnwidth - 6\tabcolsep) * \real{0.2500}}
  >{\raggedright\arraybackslash}p{(\columnwidth - 6\tabcolsep) * \real{0.2500}}@{}}
\toprule\noalign{}
\begin{minipage}[b]{\linewidth}\raggedright
order \(p\)
\end{minipage} & \begin{minipage}[b]{\linewidth}\raggedright
per unit weight (warped, time)
\end{minipage} & \begin{minipage}[b]{\linewidth}\raggedright
original weights (warped, time)
\end{minipage} & \begin{minipage}[b]{\linewidth}\raggedright
corrected weights (warped, time)
\end{minipage} \\
\midrule\noalign{}
\endhead
\bottomrule\noalign{}
\endlastfoot
1, Einstein & \((15, 21)\) & \((15, 21)\) & \((15, 21)\) \\
2, Gauss-Bonnet & \((-180, -420)\) & \((-45/8, -105/8)\) &
\((-9, -21)\) \\
3, cubic & \((360, 2520)\) & \((5/16, 35/16)\) & \((1, 7)\) \\
the cosmological constant & & \((-10, -10)\) & \((-7, -7)\) \\
total & & \((-5/16, 1/16)\) & \((0, 0)\) \\
\end{longtable}

For general weights this is FE-70:

\[
\begin{aligned}
&\sum_{p=1}^{3}\alpha_{(p)}\,\mathcal{E}_{(p)}{}^1{}_1 + \Lambda = \big(15 - 180w_{(2)} + 360w_{(3)}\big)H^2 + \Lambda,\\
&\sum_{p=1}^{3}\alpha_{(p)}\,\mathcal{E}_{(p)}{}^4{}_4 + \Lambda = \big(21 - 420w_{(2)} + 2520w_{(3)}\big)H^2 + \Lambda,\\
&w_{(2)} = \frac{1}{20},\ w_{(3)} = \frac{1}{360}:\quad \text{both } 7H^2 + \Lambda,\qquad \Lambda = -7H^2;\\
&\text{curvature } -\mathsf{y}H^2:\quad (\mathsf{y} - 1)^2(\mathsf{y} - 7)\,H^2,\quad 7(\mathsf{y} - 1)^3H^2
\end{aligned} \tag{FE-70}
\] \textbf{Symbols in (FE-70).}

\begin{itemize}
\tightlist
\item
  \(p\): The Lovelock order \(p = 1, 2, 3\) (Einstein, Gauss-Bonnet,
  cubic), as an index and as a number; in the Kohn-Sham equations a
  second label of a Kohn-Sham orbital (as \(o\), with \(\delta_{op}\)
  the Kronecker delta of the two labels). Status: STANDARD.
\item
  \(\alpha_{(p)}\): The Lovelock coupling of order \(p\),
  \(\alpha_{(p)} = w_{(p)}H^{2-2p}\) (\(\alpha_{(1)} = 1\)). Unit:
  GeV\(^{2-2p}\). Status: AUTHOR.
\item
  \(\mathcal{E}_{(p)}{}^M{}_N\): The Lovelock tensor of order \(p\): the
  variation of \(\sqrt{\lvert g\rvert}\,\mathcal{L}_{(p)}\) with respect
  to the metric (\(\mathcal{E}_{(1)}\) the Einstein tensor,
  \(\mathcal{E}_{(2)}\) the Gauss-Bonnet tensor); with frame indices its
  frame components. Unit: 1/length\(^{2p}\). Status: STANDARD.
\item
  \(\Lambda\): The eight-dimensional cosmological constant,
  \(\Lambda = -7H^2\): with the corrected weights the value for which
  the static block is a vacuum (its Lovelock tensors add to \(7H^2\) in
  every direction). The author's original value was \(\Lambda = -10H^2\)
  (he writes the pure number \(\Lambda_{\mathrm{N}} = 10\) in
  \(0 = -\Lambda_{\mathrm{N}} + \sum_p H^{-2p}w_{(p)}\mathcal{E}_{(p)}\)).
  Unit: GeV\(^2\). Value: \(-7H^2\); the author's original \(-10H^2\).
  Status: AUTHOR (the term), DERIVED (the value).
\item
  \(w_{(p)}\): The Lovelock weights, \(w_{(1)} = 1\),
  \(w_{(2)} = 1/20\), \(w_{(3)} = 1/360\): the only weights for which
  gravitational waves travel at the speed of light in the bulk and on
  the fold (both coefficients of the bulk wave,
  \(1 - 40w_{(2)} + 360w_{(3)}\) and \(1 - 24w_{(2)} + 72w_{(3)}\),
  vanish). The author's original weights \(w_{(2)} = 1/32\),
  \(w_{(3)} = 1/1152\) (Dirac\_Einstein-Lovelock\_4+4\_Nash.nb,
  In{[}134{]}) made bulk gravitons faster than light; they were
  corrected at the author's order of 2026-10-25. Unit: dimensionless.
  Value: \((1, 1/20, 1/360)\); the author's original
  \((1, 1/32, 1/1152)\). Status: AUTHOR (the form), DERIVED (the
  values).
\item
  \(H\): The author's constant (an inverse length) that fixes the period
  \(\pi/(6H)\) of \(x_0\), the Lovelock couplings
  \(\alpha_{(p)} = w_{(p)}H^{2-2p}\) and the cosmological constant
  \(\Lambda = -7H^2\) (the author's original value \(-10H^2\) was
  corrected on 2026-10-25); in UM-3 a free parameter. Unit: GeV
  (1/length). Status: AUTHOR.
\item
  \(\mathsf{y}\): The curvature of a seven-dimensional space of constant
  curvature \(-\mathsf{y}H^2\) (in units of \(-H^2\)), times the time
  line: \(\mathsf{y} = 1\) is the static block. Unit: dimensionless.
\end{itemize}

A cosmological constant adds the same number to every direction, so it
can cancel the left-hand side only if the warped and the time entries
are equal. With the original weights they are not (\(155/16\) along the
warped directions and \(161/16\) along the time before \(\Lambda\)), and
no value of \(\Lambda\) makes the static block a vacuum. Status: DERIVED
(ST1 by computer algebra on the static block, ST9;
\texttt{um3\_weights.py}, W4, for general weights).

\hypertarget{the-stabiliser}{%
\subsection{The stabiliser}\label{the-stabiliser}}

With the original weights and \(\Lambda = -10H^2\) the left-hand side of
FE-1 at the static block is therefore, on both sides of the fold and
with no entry that mixes two directions,

\[
\begin{aligned}
&w_{(2)} = \frac{1}{32},\ w_{(3)} = \frac{1}{1152},\ \Lambda = -10H^2:\\
&\sum_{p=1}^{3}\alpha_{(p)}\,\mathcal{E}_{(p)}{}^a{}_b + \Lambda\,\delta^a{}_b = H^2\,\mathrm{diag}\Big(-\frac{5}{16}, -\frac{5}{16}, -\frac{5}{16},
 -\frac{5}{16}, \frac{1}{16}, -\frac{5}{16}, -\frac{5}{16}, -\frac{5}{16}\Big) = \kappa\,\mathcal{T}_{\mathrm{stab}}{}^a{}_b,\\
&\mathcal{T}_{\mathrm{stab}}{}^a{}_b = \mathrm{diag}\big(p_{\mathrm{s}}, p_{\mathrm{s}}, p_{\mathrm{s}}, p_{\mathrm{s}},
 -\rho_{\mathrm{s}}, p_{\mathrm{s}}, p_{\mathrm{s}}, p_{\mathrm{s}}\big),\qquad
 \kappa\rho_{\mathrm{s}} = -\frac{H^2}{16},\qquad \kappa p_{\mathrm{s}} = -\frac{5H^2}{16}
\end{aligned} \tag{SB-9}
\] \textbf{Symbols in (SB-9).}

\begin{itemize}
\tightlist
\item
  \(w_{(p)}\): The Lovelock weights, \(w_{(1)} = 1\),
  \(w_{(2)} = 1/20\), \(w_{(3)} = 1/360\): the only weights for which
  gravitational waves travel at the speed of light in the bulk and on
  the fold (both coefficients of the bulk wave,
  \(1 - 40w_{(2)} + 360w_{(3)}\) and \(1 - 24w_{(2)} + 72w_{(3)}\),
  vanish). The author's original weights \(w_{(2)} = 1/32\),
  \(w_{(3)} = 1/1152\) (Dirac\_Einstein-Lovelock\_4+4\_Nash.nb,
  In{[}134{]}) made bulk gravitons faster than light; they were
  corrected at the author's order of 2026-10-25. Unit: dimensionless.
  Value: \((1, 1/20, 1/360)\); the author's original
  \((1, 1/32, 1/1152)\). Status: AUTHOR (the form), DERIVED (the
  values).
\item
  \(\Lambda\): The eight-dimensional cosmological constant,
  \(\Lambda = -7H^2\): with the corrected weights the value for which
  the static block is a vacuum (its Lovelock tensors add to \(7H^2\) in
  every direction). The author's original value was \(\Lambda = -10H^2\)
  (he writes the pure number \(\Lambda_{\mathrm{N}} = 10\) in
  \(0 = -\Lambda_{\mathrm{N}} + \sum_p H^{-2p}w_{(p)}\mathcal{E}_{(p)}\)).
  Unit: GeV\(^2\). Value: \(-7H^2\); the author's original \(-10H^2\).
  Status: AUTHOR (the term), DERIVED (the value).
\item
  \(H\): The author's constant (an inverse length) that fixes the period
  \(\pi/(6H)\) of \(x_0\), the Lovelock couplings
  \(\alpha_{(p)} = w_{(p)}H^{2-2p}\) and the cosmological constant
  \(\Lambda = -7H^2\) (the author's original value \(-10H^2\) was
  corrected on 2026-10-25); in UM-3 a free parameter. Unit: GeV
  (1/length). Status: AUTHOR.
\item
  \(p\): The Lovelock order \(p = 1, 2, 3\) (Einstein, Gauss-Bonnet,
  cubic), as an index and as a number; in the Kohn-Sham equations a
  second label of a Kohn-Sham orbital (as \(o\), with \(\delta_{op}\)
  the Kronecker delta of the two labels). Status: STANDARD.
\item
  \(\alpha_{(p)}\): The Lovelock coupling of order \(p\),
  \(\alpha_{(p)} = w_{(p)}H^{2-2p}\) (\(\alpha_{(1)} = 1\)). Unit:
  GeV\(^{2-2p}\). Status: AUTHOR.
\item
  \(\mathcal{E}_{(p)}{}^M{}_N\): The Lovelock tensor of order \(p\): the
  variation of \(\sqrt{\lvert g\rvert}\,\mathcal{L}_{(p)}\) with respect
  to the metric (\(\mathcal{E}_{(1)}\) the Einstein tensor,
  \(\mathcal{E}_{(2)}\) the Gauss-Bonnet tensor); with frame indices its
  frame components. Unit: 1/length\(^{2p}\). Status: STANDARD.
\item
  \(a\): A frame (orthonormal) index running over \(0, 1, \dots, 7\)
  (frame direction \((0)\) along \(\sigma\), \((1), (2), (3)\) along
  ordinary space, \((4)\) along the time, \((5), (6), (7)\) along the
  extra times); raised and lowered with the frame metric
  \(\eta_{ab} = \mathrm{diag}(1, 1, 1, 1, -1, -1, -1, -1)\); on the fold
  (\(a, b \in \{1, \dots, 7\}\)) the frame of the fold; a repeated index
  is summed. Status: STANDARD.
\item
  \(b\): A frame index (as \(a\)). Status: STANDARD.
\item
  \(\delta\): The Kronecker delta: with two indices \(\delta^M_N = 1\)
  if \(M = N\) and 0 otherwise (likewise \(\delta_{X\mathrm{L}}\) is 1
  if the label \(X\) of a half of dirac16complex is L and 0 otherwise);
  with upper and lower index strings the generalised Kronecker delta
  \(\delta^{M_1\cdots M_n}_{N_1\cdots N_n}\), the determinant of the
  \(n \times n\) matrix of the Kronecker deltas \(\delta^{M_a}_{N_b}\)
  (\(+1\) or \(-1\) if the upper indices are an even or odd permutation
  of the lower ones, 0 otherwise). Status: STANDARD.
\item
  \(\mathrm{diag}\): A diagonal matrix with the listed entries on its
  diagonal (in the order of the coordinates \(x_0, x_1, \dots, x_7\) for
  a tensor of eight dimensions, \(x_1, \dots, x_7\) for a tensor of the
  fold, \(x_1, \dots, x_4\) for a four-dimensional tensor). Status:
  STANDARD.
\item
  \(\kappa\): The eight-dimensional gravitational coupling,
  \(\kappa = 8\pi G_8\). Unit: GeV\(^{-6}\). Status: AUTHOR.
\item
  \(\mathcal{T}_{\mathrm{stab}}{}^a{}_b\): The energy-momentum tensor of
  the stabiliser of the bulk (frame components): the source that a
  static bulk of the author's metric required with the author's original
  Lovelock weights, found and named by Claude (Anthropic) on 2026-10-24;
  with the corrected weights and \(\Lambda = -7H^2\) it is zero (IN-7).
  Unit: GeV\(^8\). Value: 0 with the corrected weights. Status: DERIVED.
\item
  \(p_{\mathrm{s}}\): The pressure of the stabiliser of the bulk with
  the author's original weights (equal along \(\sigma\), ordinary space
  and the extra times). Unit: GeV\(^8\). Value:
  \(\kappa p_{\mathrm{s}} = -5H^2/16\) with the author's original
  weights; 0 with the corrected weights. Status: DERIVED.
\item
  \(\rho_{\mathrm{s}}\): The energy density of the stabiliser of the
  bulk with the author's original weights; with frozen extra times it
  was the energy density of the bulk at which ordinary space stands
  still. Unit: GeV\(^8\). Value: \(\kappa\rho_{\mathrm{s}} = -H^2/16\)
  with the author's original weights; 0 with the corrected weights.
  Status: DERIVED.
\end{itemize}

in the order (\(\sigma\), \(x_1\), \(x_2\), \(x_3\), \(x_4\), \(x_5\),
\(x_6\), \(x_7\)). The static bulk was therefore not a vacuum: the field
equations held only if the bulk contained a source whose energy-momentum
tensor is this left-hand side divided by \(\kappa\), a medium at rest in
the language of SB-3, without any flow, with the energy density
\(\rho_{\mathrm{s}}\) and the same pressure
\(p_{\mathrm{s}} = 5\rho_{\mathrm{s}}\) in every direction but the time.
Calling this source a medium of its own, the stabiliser, assumes that
dirac16complex contributes nothing to the static bulk (its
energy-momentum tensor \(\mathcal{T}_\Psi\) zero or negligible there;
ASSUMED). Status: DERIVED (ST1).

The kink of the static block needed, with the original weights, the
source of the fold

\[
\kappa\,\mathcal{S}^{(i)}{}_{(i)} = \kappa\,\mathcal{S}^{(r)}{}_{(r)} = -\frac{61}{8}H,\qquad
\kappa\,\mathcal{S}^{(4)}{}_{(4)} = -\frac{31}{4}H\qquad(\text{original weights, static block}) \tag{SB-10}
\] \textbf{Symbols in (SB-10).}

\begin{itemize}
\tightlist
\item
  \(\kappa\): The eight-dimensional gravitational coupling,
  \(\kappa = 8\pi G_8\). Unit: GeV\(^{-6}\). Status: AUTHOR.
\item
  \(\mathcal{S}^a{}_b\): The thin source of the fold: the
  energy-momentum tensor concentrated on the fold, integrated across it,
  \(\kappa\mathcal{S}^a{}_b = \int_{0^-}^{0^+}\mathrm{d}\sigma\sum_p\alpha_{(p)}\mathcal{E}_{(p)}{}^a{}_b\)
  (FD-11, the junction conditions); \(\mathcal{S}^4{}_4\) is minus its
  energy density, \(\mathcal{S}^1{}_1\) and \(\mathcal{S}^5{}_5\) its
  pressures along ordinary space and along the extra times (frame
  components along the fold, eight-dimensional normalisation: an energy
  per unit volume of the six other directions of the fold). Unit:
  GeV\(^7\). Value:
  \(\kappa\mathcal{S}^a{}_b = -\frac{32}{5}H\delta^a{}_b\) at the static
  block. Status: DERIVED.
\item
  \(i\): A generation index (\(i = 1, 2, 3\)), which also labels the
  neutrino mass states \(\nu_1, \nu_2, \nu_3\); in a frame component
  such as \(\mathcal{S}^{(i)}{}_{(i)}\) a direction \(x_i\) of ordinary
  space (the three components are equal); a repeated index is not summed
  unless a sum sign is written. Status: STANDARD.
\item
  \(r\): An extra-time direction of the fold: \(r \in \{5, 6, 7\}\)
  (\(x_5, x_6, x_7\)); not summed (the three components are equal).
  Status: STANDARD.
\item
  \(H\): The author's constant (an inverse length) that fixes the period
  \(\pi/(6H)\) of \(x_0\), the Lovelock couplings
  \(\alpha_{(p)} = w_{(p)}H^{2-2p}\) and the cosmological constant
  \(\Lambda = -7H^2\) (the author's original value \(-10H^2\) was
  corrected on 2026-10-25); in UM-3 a free parameter. Unit: GeV
  (1/length). Status: AUTHOR.
\end{itemize}

along ordinary space, the time and the extra times. Order by order (in
units of \(H\)): Einstein \((-10, -12, -10)\), which is Israel's
\(-([K^a{}_b] - \delta^a{}_b[K])\) with the jump \(-2H\) of the
extrinsic curvature along the six warped directions; Gauss-Bonnet
\((5/2, 5, 5/2)\); cubic \((-1/8, -3/4, -1/8)\). It splits into a
tension \(61H/8\) (pressures equal to minus the energy density) and a
pressureless part \(H/8\). With the corrected weights the Gauss-Bonnet
and cubic parts are \((4, 8, 4)\) and \((-2/5, -12/5, -2/5)\), and the
total is the pure tension \(-32H/5\) in all seven directions (FE-32
below). Status: DERIVED (ST1 and ST9, from the junction conditions
FD-11; \texttt{um3\_fold.py}, J5).

\hypertarget{why-it-is-called-the-stabiliser-and-who-found-it}{%
\subsection{Why it is called the stabiliser, and who found
it}\label{why-it-is-called-the-stabiliser-and-who-found-it}}

It was the energy content without which the bulk of the author's metric
could not stay static: it \emph{stabilised} the bulk in the sense of
keeping it static. (Whether the static bulk with the stabiliser was also
stable, that is whether small disturbances die away, is a separate
question that was not analysed.) The stabiliser was not part of the
author's specification. Claude (Anthropic's AI model, working in Claude
Code) found the requirement SB-9 on 2026-10-24 while deriving the
complete coupled field equations of UM-3, introduced the medium to meet
it, and named it the stabiliser of the bulk. The name and the idea of a
medium that holds a static block had appeared before in this repository,
in the two earlier universe models, for their own metrics and with other
properties (the provenance gives the lineage). Status: DERIVED (the
requirement), ASSUMED (the medium).

\hypertarget{what-kind-of-medium-it-was}{%
\subsection{What kind of medium it
was}\label{what-kind-of-medium-it-was}}

\textbf{Its energy density and its pressure were negative.}
\(\rho_{\mathrm{s}} = -H^2/(16\kappa)\), and the pressure five times as
large, \(p_{\mathrm{s}} = 5\rho_{\mathrm{s}}\), the same along the
hidden direction, ordinary space and the extra times. It did not flow.

\textbf{It violated both energy conditions.}

\[
\kappa\rho_{\mathrm{s}} = -\frac{H^2}{16} < 0,\qquad \kappa\big(\rho_{\mathrm{s}} + p_{\mathrm{s}}\big) = -\frac{3H^2}{8} < 0 \tag{SB-11}
\] \textbf{Symbols in (SB-11).}

\begin{itemize}
\tightlist
\item
  \(\kappa\): The eight-dimensional gravitational coupling,
  \(\kappa = 8\pi G_8\). Unit: GeV\(^{-6}\). Status: AUTHOR.
\item
  \(\rho_{\mathrm{s}}\): The energy density of the stabiliser of the
  bulk with the author's original weights; with frozen extra times it
  was the energy density of the bulk at which ordinary space stands
  still. Unit: GeV\(^8\). Value: \(\kappa\rho_{\mathrm{s}} = -H^2/16\)
  with the author's original weights; 0 with the corrected weights.
  Status: DERIVED.
\item
  \(H\): The author's constant (an inverse length) that fixes the period
  \(\pi/(6H)\) of \(x_0\), the Lovelock couplings
  \(\alpha_{(p)} = w_{(p)}H^{2-2p}\) and the cosmological constant
  \(\Lambda = -7H^2\) (the author's original value \(-10H^2\) was
  corrected on 2026-10-25); in UM-3 a free parameter. Unit: GeV
  (1/length). Status: AUTHOR.
\item
  \(p_{\mathrm{s}}\): The pressure of the stabiliser of the bulk with
  the author's original weights (equal along \(\sigma\), ordinary space
  and the extra times). Unit: GeV\(^8\). Value:
  \(\kappa p_{\mathrm{s}} = -5H^2/16\) with the author's original
  weights; 0 with the corrected weights. Status: DERIVED.
\end{itemize}

The second inequality, for light moving along the time and one
space-like direction, is the violation of the null energy condition; it
is what made the stabiliser exotic. The negative cosmological constant
of the bulk violates only the weak condition, since its energy density
plus its pressure is zero. Status: DERIVED (ST5).

\textbf{Among homogeneous scalar fields, only a ghost reproduced it.} A
scalar field that fills the bulk uniformly and changes only with the
time has a kinetic energy \(K_{\mathrm{f}}\) and a potential energy
\(V_{\mathrm{f}}\) per volume, and its energy density and pressure (the
same in every direction but the time) are

\[
\rho_{\mathrm{med}} = K_{\mathrm{f}} + V_{\mathrm{f}},\qquad p_{\mathrm{med}} = K_{\mathrm{f}} - V_{\mathrm{f}};\qquad
p_{\mathrm{med}} = 5\rho_{\mathrm{med}}:\quad V_{\mathrm{f}} = -\frac{2}{3}K_{\mathrm{f}},\quad \rho_{\mathrm{med}} = \frac{K_{\mathrm{f}}}{3} \tag{SB-12}
\] \textbf{Symbols in (SB-12).}

\begin{itemize}
\tightlist
\item
  \(\rho_{\mathrm{med}}\): The energy density (energy per volume) of a
  medium at rest. Unit: GeV\(^8\). Status: STANDARD.
\item
  \(K_{\mathrm{f}}\): The kinetic energy per volume of a field that
  fills the bulk uniformly and changes only with the time \(x_4\): half
  the square of its rate of change for an ordinary field (never
  negative), minus that for a field whose kinetic term has the opposite
  sign (a ghost). Unit: GeV\(^8\). Status: STANDARD.
\item
  \(V_{\mathrm{f}}\): The potential energy per volume of a field that
  fills the bulk uniformly and changes only with the time \(x_4\). Unit:
  GeV\(^8\). Status: STANDARD.
\item
  \(p_{\mathrm{med}}\): The pressure of a medium at rest, the same along
  every direction but the time (for a medium of the bulk: along
  \(\sigma\), ordinary space and the extra times). Unit: GeV\(^8\).
  Status: STANDARD.
\end{itemize}

(derived by varying the action of the field with respect to the metric,
ST5). An ordinary field has \(K_{\mathrm{f}} \ge 0\) and so, with a
pressure five times its energy density, a positive energy density. Only
a field whose kinetic energy is negative, a so-called ghost, with
\(K_{\mathrm{f}} = 3\rho_{\mathrm{s}}\) and
\(V_{\mathrm{f}} = -2\rho_{\mathrm{s}}\), reproduced the stabiliser; on
the static bulk it would roll at a constant speed on a flat stretch of
its potential. A scalar field that depends on \(\sigma\) instead of the
time could not do it, because its pressure along \(\sigma\) differs from
its pressure along the other directions unless it is constant, and a
constant field is a vacuum energy with \(p = -\rho\). Vector fields and
other kinds of fields were not examined. Status: DERIVED (ST5, for
homogeneous scalar fields).

\textbf{If it had been a simple fluid, its sound would have been faster
than light.} For a fluid whose pressure is a fixed multiple of its
energy density, small compressions travel at the speed of sound
\(c_{\mathrm{s}}\) with
\(c_{\mathrm{s}}^2 = \mathrm{d}p/\mathrm{d}\rho\); for the stabiliser
this is 5, so its sound would have travelled \(\sqrt5 \approx 2.24\)
times faster than light. The ghost field instead propagates at the speed
of light, but its quanta carry negative energy. A general theorem
(Dubovsky, Grégoire, Nicolis and Rattazzi, 2006) states that, under
broad conditions, a medium that violates the null energy condition has
either signals faster than light or unstable excitations. Status:
STANDARD (the theorem and the fluid formula), DERIVED (their
application).

\textbf{It was conserved by itself on the static bulk.} The covariant
divergence SB-4 of the constant tensor of SB-9 vanishes on the static
block, in every component (ST4). The reason: its pressure along the
hidden direction equals its pressure along every warped direction, and
the time, along which its entry differs, is not warped. Unlike the air
of SB-7, the stabiliser needed no change of pressure along the hidden
direction: because the time is not warped there is no gravitational pull
along \(\sigma\). Status: DERIVED (ST4).

\hypertarget{how-it-would-have-evolved-on-its-own}{%
\subsection{How it would have evolved on its
own}\label{how-it-would-have-evolved-on-its-own}}

A medium that interacts only through gravity obeys SB-4. For a medium
that fills the bulk uniformly, with the energy density
\(\rho_{\mathrm{med}}(x_4)\) and the pressure \(p_{\mathrm{med}}(x_4)\)
in every warped direction, in the author's metric with arbitrary
\(a_4(x_4)\) and \(b_4(x_4)\), the time component of SB-4 reads (ST6)

\[
\rho_{\mathrm{med}}' = -3\,\frac{b_4'}{b_4}\,\big(\rho_{\mathrm{med}} + p_{\mathrm{med}}\big);\qquad
p_{\mathrm{med}} = 5\rho_{\mathrm{med}}:\quad \rho_{\mathrm{med}}\,b_4^{18} = \text{constant} \tag{SB-13}
\] \textbf{Symbols in (SB-13).}

\begin{itemize}
\tightlist
\item
  \(\rho_{\mathrm{med}}\): The energy density (energy per volume) of a
  medium at rest. Unit: GeV\(^8\). Status: STANDARD.
\item
  \({}'\): The derivative of a function of the time \(x_4\) alone with
  respect to \(x_4\) (\(a_4' = \mathrm{d}a_4/\mathrm{d}x_4\); two primes
  the second derivative); on a function of \(\sigma\) alone (a profile)
  the derivative with respect to \(\sigma\). Status: STANDARD.
\item
  \(b_4\): The author's metric function of the time \(x_4\) in the
  extra-time components, \(b_4 > 0\), never identically 1 (the author's
  order of 2026-10-24); the scale factor (the factor of the lengths) of
  the extra times is \(\mathrm{e}^{-a_4}b_4\) on the fold and
  \(\mathrm{e}^{-a_4}b_4\,\mathrm{e}^{-H\lvert\sigma\rvert}\) at the
  distance \(\sigma\) from it (the square root of the warp factor times
  \(\mathrm{e}^{-a_4}b_4\)). Unit: dimensionless. Status: AUTHOR.
\item
  \(p_{\mathrm{med}}\): The pressure of a medium at rest, the same along
  every direction but the time (for a medium of the bulk: along
  \(\sigma\), ordinary space and the extra times). Unit: GeV\(^8\).
  Status: STANDARD.
\end{itemize}

The growth of ordinary space drops out, because ordinary space grows
with \(\mathrm{e}^{a_4}\) and the extra times with
\(\mathrm{e}^{-a_4}b_4\), so that the volume of all six warped
directions together grows with \(b_4^3\). Two consequences follow. A
stabiliser that is absent (\(\rho = 0\)) stays absent. And with frozen
extra times (\(b_4\) proportional to \(\mathrm{e}^{a_4}\)) a stabiliser
on its own would have been diluted as \(\mathrm{e}^{-18a_4}\) while the
universe expands. Status: DERIVED (ST6).

\hypertarget{what-the-original-equations-required-when-the-universe-expanded}{%
\subsection{What the original equations required when the universe
expanded}\label{what-the-original-equations-required-when-the-universe-expanded}}

With frozen extra times, the condition of the late universe
(\(\mathsf{U} = \mathsf{A}\): the size \(\mathrm{e}^{-a_4}b_4\) of the
extra times constant, so that the couplings of the Standard Model stay
constant), the original equations required in the bulk and on the fold
(with \(\mathsf{A}_1 = a_4''/H^2\) the rate of change of
\(\mathsf{A}\)):

\[
\begin{aligned}
&\mathsf{U} = \mathsf{A}:\qquad \frac{\kappa\rho_{\mathrm{bulk}}}{H^2} = \frac{3\mathsf{A}^2 - 1}{16},\qquad
 \frac{\kappa\,\mathcal{T}_{\mathrm{bulk}}{}^{(0)}{}_{(0)}}{H^2} = -\frac{3\mathsf{A}^4 + 3\mathsf{A}^2\mathsf{A}_1 + 18\mathsf{A}^2 + 3\mathsf{A}_1 + 5}{16},\\
&\frac{\kappa\,\mathcal{T}_{\mathrm{bulk}}{}^{(i)}{}_{(i)}}{H^2} = -\frac{9\mathsf{A}^2 + 2\mathsf{A}_1 + 5}{16},\qquad
 \frac{\kappa\,\mathcal{T}_{\mathrm{bulk}}{}^{(r)}{}_{(r)}}{H^2} = -\frac{3\mathsf{A}^4 + 3\mathsf{A}^2\mathsf{A}_1 + 24\mathsf{A}^2 + 3\mathsf{A}_1 + 5}{16},\\
&\frac{\kappa\,\mathcal{T}_{\mathrm{bulk}}{}^{(0)}{}_{(4)}}{H^2} = \mathrm{sgn}(\sigma)\,\frac{3\mathsf{A}\big(1 + \mathsf{A}^2\big)}{16},\qquad
 \frac{\kappa\,\mathcal{S}^{(4)}{}_{(4)}}{H} = -\frac{10\mathsf{A}^2 + 31}{4},\\
&\frac{\kappa\,\mathcal{S}^{(i)}{}_{(i)}}{H} = -\frac{57\mathsf{A}^2 + 40\mathsf{A}_1 + 183}{24},\qquad
 \frac{\kappa\,\mathcal{S}^{(r)}{}_{(r)}}{H} = -\frac{2\mathsf{A}^4 + 2\mathsf{A}^2\mathsf{A}_1 + 37\mathsf{A}^2 + 20\mathsf{A}_1 + 61}{8}
\end{aligned} \tag{SB-14}
\] \textbf{Symbols in (SB-14).}

\begin{itemize}
\tightlist
\item
  \(\mathsf{U}\): The rate of change of \(b_4\) in units of \(H\),
  \(\mathsf{U} = b_4'/(Hb_4)\). Unit: dimensionless. Value:
  \(\mathsf{U} = \mathsf{A}\) with frozen extra times.
\item
  \(\mathsf{A}\): The expansion rate of ordinary space in units of
  \(H\), \(\mathsf{A} = a_4'/H\). Unit: dimensionless. Value: 0 at the
  static block; \(H_4^{(0)}/H\) today (\(H_4^{(0)}\) the Hubble
  constant); 1 at the author's vacuum, which is not a solution with the
  corrected weights.
\item
  \(\mathsf{A}_1\): The time derivative of \(\mathsf{A}\) in units of
  \(H\), \(\mathsf{A}_1 = \mathsf{A}'/H = a_4''/H^2\). Unit:
  dimensionless.
\item
  \(\kappa\): The eight-dimensional gravitational coupling,
  \(\kappa = 8\pi G_8\). Unit: GeV\(^{-6}\). Status: AUTHOR.
\item
  \(\rho_{\mathrm{bulk}}\): The energy density of the bulk sources,
  \(\rho_{\mathrm{bulk}} = -\mathcal{T}_{\mathrm{bulk}}{}^{(4)}{}_{(4)}\).
  Unit: GeV\(^8\). Status: DERIVED.
\item
  \(H\): The author's constant (an inverse length) that fixes the period
  \(\pi/(6H)\) of \(x_0\), the Lovelock couplings
  \(\alpha_{(p)} = w_{(p)}H^{2-2p}\) and the cosmological constant
  \(\Lambda = -7H^2\) (the author's original value \(-10H^2\) was
  corrected on 2026-10-25); in UM-3 a free parameter. Unit: GeV
  (1/length). Status: AUTHOR.
\item
  \(\mathcal{T}_{\mathrm{bulk}}\): The energy-momentum tensor of all
  sources of the bulk (with the corrected weights dirac16complex alone,
  \(\mathcal{T}_\Psi\); with the author's original weights also the
  stabiliser); \(\mathcal{T}_{\mathrm{bulk}}{}^{(0)}{}_b\) is its flux
  along \(\sigma\) (frame components). Unit: GeV\(^8\). Status:
  STANDARD.
\item
  \(i\): A generation index (\(i = 1, 2, 3\)), which also labels the
  neutrino mass states \(\nu_1, \nu_2, \nu_3\); in a frame component
  such as \(\mathcal{S}^{(i)}{}_{(i)}\) a direction \(x_i\) of ordinary
  space (the three components are equal); a repeated index is not summed
  unless a sum sign is written. Status: STANDARD.
\item
  \(r\): An extra-time direction of the fold: \(r \in \{5, 6, 7\}\)
  (\(x_5, x_6, x_7\)); not summed (the three components are equal).
  Status: STANDARD.
\item
  \(\mathrm{sgn}\): The sign function: \(\mathrm{sgn}(y) = +1\) for
  \(y > 0\), \(-1\) for \(y < 0\). Status: STANDARD.
\item
  \(\sigma\): The proper coordinate of the hidden direction:
  \(\mathrm{d}\sigma = \lvert\cot z\rvert\,\mathrm{d}x_0\),
  \(\sigma = \mathrm{sgn}(\pi/2 - z)\ln(\sin z)/(6H)\) on one period
  \(0 < z < \pi\); \(\sigma = 0\) is the fold \(z = \pi/2\),
  \(\sigma \to \mp\infty\) the tips \(z \to 0, \pi\); one period of
  \(x_0\) is the whole line \(-\infty < \sigma < \infty\). Unit: length.
  Status: DERIVED.
\item
  \(\mathcal{S}^a{}_b\): The thin source of the fold: the
  energy-momentum tensor concentrated on the fold, integrated across it,
  \(\kappa\mathcal{S}^a{}_b = \int_{0^-}^{0^+}\mathrm{d}\sigma\sum_p\alpha_{(p)}\mathcal{E}_{(p)}{}^a{}_b\)
  (FD-11, the junction conditions); \(\mathcal{S}^4{}_4\) is minus its
  energy density, \(\mathcal{S}^1{}_1\) and \(\mathcal{S}^5{}_5\) its
  pressures along ordinary space and along the extra times (frame
  components along the fold, eight-dimensional normalisation: an energy
  per unit volume of the six other directions of the fold). Unit:
  GeV\(^7\). Value:
  \(\kappa\mathcal{S}^a{}_b = -\frac{32}{5}H\delta^a{}_b\) at the static
  block. Status: DERIVED.
\end{itemize}

(\(\rho_{\mathrm{bulk}} = -\mathcal{T}_{\mathrm{bulk}}{}^{(4)}{}_{(4)}\)
is the energy density of all bulk sources together.) For
\(\mathsf{A} \to 0\) they reduce to SB-9 and SB-10. The entry
\(\mathcal{T}_{\mathrm{bulk}}{}^{(0)}{}_{(4)}\) is a flow of energy
along the hidden direction, equal in size on the two sides of the fold
and opposite in sign, so it jumps across the fold by
\(\frac38\mathsf{A}(1 + \mathsf{A}^2)H^2/\kappa\); the energy current,
minus this entry, points towards the fold on both sides. Status: DERIVED
(ST6, with its \(\mathsf{A}_1\) terms; \texttt{um3\_fold.py}, J5 and
J7).

\textbf{Where the energy came from.} With frozen extra times only
ordinary space grows, and the volume of the bulk grows at the rate
\(3a_4'\). The required bulk source had, for constant \(\mathsf{A}\), an
energy density plus pressure along ordinary space that is negative:

\[
\begin{aligned}
&\kappa\big(\rho_{\mathrm{bulk}} + \mathcal{T}_{\mathrm{bulk}}{}^{(1)}{}_{(1)}\big) = -\frac{3}{8}\big(1 + \mathsf{A}^2\big)H^2\qquad
 (\mathsf{U} = \mathsf{A},\ \mathsf{A}_1 = 0),\\
&\int_{-\infty}^{\infty}\mathrm{d}\sigma\;\mathrm{e}^{-6H\lvert\sigma\rvert}\;\big(-3a_4'\big)\,\kappa\big(\rho_{\mathrm{bulk}}
 + \mathcal{T}_{\mathrm{bulk}}{}^{(1)}{}_{(1)}\big) = \frac{3}{8}\,\mathsf{A}\big(1 + \mathsf{A}^2\big)H^2
\end{aligned} \tag{SB-15}
\] \textbf{Symbols in (SB-15).}

\begin{itemize}
\tightlist
\item
  \(\kappa\): The eight-dimensional gravitational coupling,
  \(\kappa = 8\pi G_8\). Unit: GeV\(^{-6}\). Status: AUTHOR.
\item
  \(\rho_{\mathrm{bulk}}\): The energy density of the bulk sources,
  \(\rho_{\mathrm{bulk}} = -\mathcal{T}_{\mathrm{bulk}}{}^{(4)}{}_{(4)}\).
  Unit: GeV\(^8\). Status: DERIVED.
\item
  \(\mathcal{T}_{\mathrm{bulk}}\): The energy-momentum tensor of all
  sources of the bulk (with the corrected weights dirac16complex alone,
  \(\mathcal{T}_\Psi\); with the author's original weights also the
  stabiliser); \(\mathcal{T}_{\mathrm{bulk}}{}^{(0)}{}_b\) is its flux
  along \(\sigma\) (frame components). Unit: GeV\(^8\). Status:
  STANDARD.
\item
  \(\mathsf{A}\): The expansion rate of ordinary space in units of
  \(H\), \(\mathsf{A} = a_4'/H\). Unit: dimensionless. Value: 0 at the
  static block; \(H_4^{(0)}/H\) today (\(H_4^{(0)}\) the Hubble
  constant); 1 at the author's vacuum, which is not a solution with the
  corrected weights.
\item
  \(H\): The author's constant (an inverse length) that fixes the period
  \(\pi/(6H)\) of \(x_0\), the Lovelock couplings
  \(\alpha_{(p)} = w_{(p)}H^{2-2p}\) and the cosmological constant
  \(\Lambda = -7H^2\) (the author's original value \(-10H^2\) was
  corrected on 2026-10-25); in UM-3 a free parameter. Unit: GeV
  (1/length). Status: AUTHOR.
\item
  \(\mathsf{U}\): The rate of change of \(b_4\) in units of \(H\),
  \(\mathsf{U} = b_4'/(Hb_4)\). Unit: dimensionless. Value:
  \(\mathsf{U} = \mathsf{A}\) with frozen extra times.
\item
  \(\mathsf{A}_1\): The time derivative of \(\mathsf{A}\) in units of
  \(H\), \(\mathsf{A}_1 = \mathsf{A}'/H = a_4''/H^2\). Unit:
  dimensionless.
\item
  \(\mathrm{d}\): The differential: \(\mathrm{d}X\) is an infinitesimal
  change of \(X\), \(\mathrm{d}X/\mathrm{d}Y\) the derivative of \(X\)
  with respect to \(Y\), \(\mathrm{d}^nX\) the volume element of \(n\)
  variables. Status: STANDARD.
\item
  \(\sigma\): The proper coordinate of the hidden direction:
  \(\mathrm{d}\sigma = \lvert\cot z\rvert\,\mathrm{d}x_0\),
  \(\sigma = \mathrm{sgn}(\pi/2 - z)\ln(\sin z)/(6H)\) on one period
  \(0 < z < \pi\); \(\sigma = 0\) is the fold \(z = \pi/2\),
  \(\sigma \to \mp\infty\) the tips \(z \to 0, \pi\); one period of
  \(x_0\) is the whole line \(-\infty < \sigma < \infty\). Unit: length.
  Status: DERIVED.
\item
  \(\infty\): Infinity: as a limit of an integral or a sum it means that
  the range has no bound. Status: STANDARD.
\item
  \(\mathrm{e}\): Euler's number, the base of the natural logarithm;
  \(\mathrm{e}^{X}\) is the exponential of \(X\). Unit: dimensionless.
  Value: 2.718281828. Status: STANDARD.
\item
  \(a_4\): The author's metric function of the time \(x_4\) that scales
  ordinary 3-space: on the fold \(\mathrm{e}^{a_4}\) is the scale factor
  of cosmology and \(a_4'\) the Hubble rate. Unit: dimensionless.
  Status: AUTHOR.
\item
  \({}'\): The derivative of a function of the time \(x_4\) alone with
  respect to \(x_4\) (\(a_4' = \mathrm{d}a_4/\mathrm{d}x_4\); two primes
  the second derivative); on a function of \(\sigma\) alone (a profile)
  the derivative with respect to \(\sigma\). Status: STANDARD.
\end{itemize}

Left to itself, a medium of the bulk changes its energy density at the
rate \(-3a_4'(\rho + p)\) while only ordinary space grows. The required
source, with \(\rho + p < 0\), would have seen its negative energy
density rise towards zero; since the field equations kept its energy
density fixed, it had to give off the energy
\(-3a_4'(\rho_{\mathrm{bulk}} + \mathcal{T}_{\mathrm{bulk}}{}^{(1)}{}_{(1)})\)
per volume and time. Integrated over the hidden direction with the
volume factor \(\mathrm{e}^{-6H\lvert\sigma\rvert}\) of the six warped
directions, this is exactly the jump of the flow: all the energy that
the bulk gave off arrived at the fold. This is a consistency identity
(it follows from the Bianchi identity), not an independent result.
Status: DERIVED (ST6).

\textbf{Where the energy went.} The source of the fold took up exactly
the energy that arrived from the bulk, for any \(a_4(x_4)\):

\[
-\big(\kappa\,\mathcal{S}^{(4)}{}_{(4)}\big)' + 3a_4'\big(\kappa\,\mathcal{S}^{(i)}{}_{(i)} - \kappa\,\mathcal{S}^{(4)}{}_{(4)}\big)
 = \kappa\Big(\mathcal{T}_{\mathrm{bulk}}{}^{(0)}{}_{(4)}\big\vert_{\sigma=0^+} - \mathcal{T}_{\mathrm{bulk}}{}^{(0)}{}_{(4)}\big\vert_{\sigma=0^-}\Big)
 \qquad (\mathsf{U} = \mathsf{A}) \tag{SB-16}
\] \textbf{Symbols in (SB-16).}

\begin{itemize}
\tightlist
\item
  \(\kappa\): The eight-dimensional gravitational coupling,
  \(\kappa = 8\pi G_8\). Unit: GeV\(^{-6}\). Status: AUTHOR.
\item
  \(\mathcal{S}^a{}_b\): The thin source of the fold: the
  energy-momentum tensor concentrated on the fold, integrated across it,
  \(\kappa\mathcal{S}^a{}_b = \int_{0^-}^{0^+}\mathrm{d}\sigma\sum_p\alpha_{(p)}\mathcal{E}_{(p)}{}^a{}_b\)
  (FD-11, the junction conditions); \(\mathcal{S}^4{}_4\) is minus its
  energy density, \(\mathcal{S}^1{}_1\) and \(\mathcal{S}^5{}_5\) its
  pressures along ordinary space and along the extra times (frame
  components along the fold, eight-dimensional normalisation: an energy
  per unit volume of the six other directions of the fold). Unit:
  GeV\(^7\). Value:
  \(\kappa\mathcal{S}^a{}_b = -\frac{32}{5}H\delta^a{}_b\) at the static
  block. Status: DERIVED.
\item
  \(i\): A generation index (\(i = 1, 2, 3\)), which also labels the
  neutrino mass states \(\nu_1, \nu_2, \nu_3\); in a frame component
  such as \(\mathcal{S}^{(i)}{}_{(i)}\) a direction \(x_i\) of ordinary
  space (the three components are equal); a repeated index is not summed
  unless a sum sign is written. Status: STANDARD.
\item
  \(a_4\): The author's metric function of the time \(x_4\) that scales
  ordinary 3-space: on the fold \(\mathrm{e}^{a_4}\) is the scale factor
  of cosmology and \(a_4'\) the Hubble rate. Unit: dimensionless.
  Status: AUTHOR.
\item
  \({}'\): The derivative of a function of the time \(x_4\) alone with
  respect to \(x_4\) (\(a_4' = \mathrm{d}a_4/\mathrm{d}x_4\); two primes
  the second derivative); on a function of \(\sigma\) alone (a profile)
  the derivative with respect to \(\sigma\). Status: STANDARD.
\item
  \(\mathcal{T}_{\mathrm{bulk}}\): The energy-momentum tensor of all
  sources of the bulk (with the corrected weights dirac16complex alone,
  \(\mathcal{T}_\Psi\); with the author's original weights also the
  stabiliser); \(\mathcal{T}_{\mathrm{bulk}}{}^{(0)}{}_b\) is its flux
  along \(\sigma\) (frame components). Unit: GeV\(^8\). Status:
  STANDARD.
\item
  \(\sigma\): The proper coordinate of the hidden direction:
  \(\mathrm{d}\sigma = \lvert\cot z\rvert\,\mathrm{d}x_0\),
  \(\sigma = \mathrm{sgn}(\pi/2 - z)\ln(\sin z)/(6H)\) on one period
  \(0 < z < \pi\); \(\sigma = 0\) is the fold \(z = \pi/2\),
  \(\sigma \to \mp\infty\) the tips \(z \to 0, \pi\); one period of
  \(x_0\) is the whole line \(-\infty < \sigma < \infty\). Unit: length.
  Status: DERIVED.
\item
  \(\mathsf{U}\): The rate of change of \(b_4\) in units of \(H\),
  \(\mathsf{U} = b_4'/(Hb_4)\). Unit: dimensionless. Value:
  \(\mathsf{U} = \mathsf{A}\) with frozen extra times.
\item
  \(\mathsf{A}\): The expansion rate of ordinary space in units of
  \(H\), \(\mathsf{A} = a_4'/H\). Unit: dimensionless. Value: 0 at the
  static block; \(H_4^{(0)}/H\) today (\(H_4^{(0)}\) the Hubble
  constant); 1 at the author's vacuum, which is not a solution with the
  corrected weights.
\end{itemize}

(\(-\mathcal{S}^{(4)}{}_{(4)}\) is the energy density of the source of
the fold and \(\mathcal{S}^{(i)}{}_{(i)}\) its pressure along ordinary
space; \(0^+\) and \(0^-\) mark the limits from above and from below.)
This is the conservation law at the fold, FE-11, again an identity of
the junction conditions. At the static block the earlier version of UM-3
split the source of the fold into a tension
\(\lambda_{\mathrm{F}} = 61HV_*/(8\kappa)\) and a pressureless part
\(\rho_{\mathrm{d}} = HV_*/(8\kappa)\) (per unit volume of ordinary
space, \(V_*\) the volume of the extra times), and the pressureless part
needed exactly the energy
\(3a_4'\rho_{\mathrm{d}} = \frac38\mathsf{A}H^2V_*/\kappa\), the part of
the jump linear in \(\mathsf{A}\). Reading this as an exchange of energy
between the stabiliser and a fold fluid is an interpretation: what the
equations fix is only the total bulk source and its flow; a medium of
the bulk that couples only to gravity could also carry its flow into a
thin layer of its own (ASSUMED). Moreover, ordinary matter on the fold
needs metric functions that depend on \(\sigma\), which the separable
metric of the author does not have (\texttt{um3\_fold.py}, J7), so these
energy balances describe the fold sources that the separable metric
requires, not ordinary matter. Status: DERIVED (ST6), ASSUMED (the
interpretation). The conservation law at the fold, for every metric of
UM-3:

\[
\begin{aligned}
&\nabla_M\,\mathcal{T}_\Psi{}^M{}_N = 0\qquad (\sigma \ne 0),\\
&\hat\nabla_a\,\mathcal{S}^a{}_b = -\Big(\mathcal{T}_{\mathrm{bulk}}{}^{(0)}{}_b\big\vert_{\sigma=0^+}
 - \mathcal{T}_{\mathrm{bulk}}{}^{(0)}{}_b\big\vert_{\sigma=0^-}\Big)
\end{aligned} \tag{FE-11}
\] \textbf{Symbols in (FE-11).}

\begin{itemize}
\tightlist
\item
  \(\nabla\): The covariant derivative of a tensor with the Levi-Civita
  connection of the metric \(g_{MN}\) (Christoffel symbols
  \(\Gamma^P{}_{MN}\)). Status: STANDARD.
\item
  \(\mathcal{T}_\Psi\): The energy-momentum tensor of dirac16complex,
  \(\mathcal{T}_{MN} = -(2/\sqrt{\lvert g\rvert})\,\delta S_\Psi/\delta g^{MN}\).
  Unit: GeV\(^8\). Status: STANDARD.
\item
  \(M\): A coordinate index of eight dimensions (also with a numbered
  copy \(M_1, M_2, \dots\)): it runs over \(x_0, x_1, \dots, x_7\); a
  repeated index is summed. Status: STANDARD.
\item
  \(N\): A coordinate index of eight dimensions (also
  \(N_1, N_2, \dots\)), running over \(x_0, \dots, x_7\); a repeated
  index is summed. Status: STANDARD.
\item
  \(\sigma\): The proper coordinate of the hidden direction:
  \(\mathrm{d}\sigma = \lvert\cot z\rvert\,\mathrm{d}x_0\),
  \(\sigma = \mathrm{sgn}(\pi/2 - z)\ln(\sin z)/(6H)\) on one period
  \(0 < z < \pi\); \(\sigma = 0\) is the fold \(z = \pi/2\),
  \(\sigma \to \mp\infty\) the tips \(z \to 0, \pi\); one period of
  \(x_0\) is the whole line \(-\infty < \sigma < \infty\). Unit: length.
  Status: DERIVED.
\item
  \(\hat\nabla\): The covariant derivative of the fold: with the
  Levi-Civita connection of the metric of the fold (for four-dimensional
  tensors of its metric
  \(\hat h_{\mu\nu} = \mathrm{diag}(\mathrm{e}^{2a_4}, \mathrm{e}^{2a_4}, \mathrm{e}^{2a_4}, -1)\)
  in the observable coordinates \(x_1, \dots, x_4\); for
  \(\hat\nabla_a\mathcal{S}^a{}_b\) of the induced metric of the
  seven-dimensional fold). Status: STANDARD.
\item
  \(\mathcal{S}^a{}_b\): The thin source of the fold: the
  energy-momentum tensor concentrated on the fold, integrated across it,
  \(\kappa\mathcal{S}^a{}_b = \int_{0^-}^{0^+}\mathrm{d}\sigma\sum_p\alpha_{(p)}\mathcal{E}_{(p)}{}^a{}_b\)
  (FD-11, the junction conditions); \(\mathcal{S}^4{}_4\) is minus its
  energy density, \(\mathcal{S}^1{}_1\) and \(\mathcal{S}^5{}_5\) its
  pressures along ordinary space and along the extra times (frame
  components along the fold, eight-dimensional normalisation: an energy
  per unit volume of the six other directions of the fold). Unit:
  GeV\(^7\). Value:
  \(\kappa\mathcal{S}^a{}_b = -\frac{32}{5}H\delta^a{}_b\) at the static
  block. Status: DERIVED.
\item
  \(a\): A frame (orthonormal) index running over \(0, 1, \dots, 7\)
  (frame direction \((0)\) along \(\sigma\), \((1), (2), (3)\) along
  ordinary space, \((4)\) along the time, \((5), (6), (7)\) along the
  extra times); raised and lowered with the frame metric
  \(\eta_{ab} = \mathrm{diag}(1, 1, 1, 1, -1, -1, -1, -1)\); on the fold
  (\(a, b \in \{1, \dots, 7\}\)) the frame of the fold; a repeated index
  is summed. Status: STANDARD.
\item
  \(b\): A frame index (as \(a\)). Status: STANDARD.
\item
  \(\mathcal{T}_{\mathrm{bulk}}\): The energy-momentum tensor of all
  sources of the bulk (with the corrected weights dirac16complex alone,
  \(\mathcal{T}_\Psi\); with the author's original weights also the
  stabiliser); \(\mathcal{T}_{\mathrm{bulk}}{}^{(0)}{}_b\) is its flux
  along \(\sigma\) (frame components). Unit: GeV\(^8\). Status:
  STANDARD.
\end{itemize}

\hypertarget{the-constraints-of-the-original-equations}{%
\subsection{The constraints of the original
equations}\label{the-constraints-of-the-original-equations}}

Two components of the bulk equations contain no second time derivative
of the metric, for any \(a_4(x_4)\) and \(b_4(x_4)\): the energy
constraint (the time component, like the Friedmann equation of
cosmology) and the flow of energy along the hidden direction. With the
original weights they were

\[
\begin{aligned}
&\kappa\,\mathcal{T}_{\mathrm{bulk}}{}^{(4)}{}_{(4)} = \frac{H^2}{16}\,\mathsf{P},\qquad
 \kappa\,\mathcal{T}_{\mathrm{bulk}}{}^{(0)}{}_{(4)} = \frac{3}{16}\,\mathsf{U}\,\mathcal{G}\,H^2\,\mathrm{sgn}(\sigma),\qquad \mathsf{V} = \mathsf{U} - \mathsf{A},\\
&\mathsf{P} = 1 - 3\mathsf{A}^2 - 3\mathsf{V}^2 - 9\mathsf{A}\mathsf{V} - 27\mathsf{A}^2\mathsf{V}^2 - 9\mathsf{A}^3\mathsf{V} - 9\mathsf{A}\mathsf{V}^3 - 5\mathsf{A}^3\mathsf{V}^3,\qquad
 \mathcal{G} = 1 + \mathsf{A}^2 + 8\mathsf{A}\mathsf{V} + \mathsf{V}^2 + 5\mathsf{A}^2\mathsf{V}^2
\end{aligned} \tag{SB-17}
\] \textbf{Symbols in (SB-17).}

\begin{itemize}
\tightlist
\item
  \(\kappa\): The eight-dimensional gravitational coupling,
  \(\kappa = 8\pi G_8\). Unit: GeV\(^{-6}\). Status: AUTHOR.
\item
  \(\mathcal{T}_{\mathrm{bulk}}\): The energy-momentum tensor of all
  sources of the bulk (with the corrected weights dirac16complex alone,
  \(\mathcal{T}_\Psi\); with the author's original weights also the
  stabiliser); \(\mathcal{T}_{\mathrm{bulk}}{}^{(0)}{}_b\) is its flux
  along \(\sigma\) (frame components). Unit: GeV\(^8\). Status:
  STANDARD.
\item
  \(H\): The author's constant (an inverse length) that fixes the period
  \(\pi/(6H)\) of \(x_0\), the Lovelock couplings
  \(\alpha_{(p)} = w_{(p)}H^{2-2p}\) and the cosmological constant
  \(\Lambda = -7H^2\) (the author's original value \(-10H^2\) was
  corrected on 2026-10-25); in UM-3 a free parameter. Unit: GeV
  (1/length). Status: AUTHOR.
\item
  \(\mathsf{P}(\mathsf{A}, \mathsf{U})\): The polynomial of the energy
  constraint of the bulk with the author's original Lovelock weights:
  the \((x_4, x_4)\) component of the bulk equations of the author's
  metric was
  \(\kappa\mathcal{T}_{\mathrm{bulk}}{}^{(4)}{}_{(4)} = (H^2/16)\,\mathsf{P}(\mathsf{A}, \mathsf{U})\),
  whatever the time derivatives of the rates (with the corrected weights
  it is \(\mathsf{A}^3(\mathsf{A} - \mathsf{U})^3H^2\)). Unit:
  dimensionless. Value: \(\mathsf{P}(0, 0) = 1\),
  \(\mathsf{P}(\pm1, 0) = 0\). Status: DERIVED.
\item
  \(\mathsf{U}\): The rate of change of \(b_4\) in units of \(H\),
  \(\mathsf{U} = b_4'/(Hb_4)\). Unit: dimensionless. Value:
  \(\mathsf{U} = \mathsf{A}\) with frozen extra times.
\item
  \(\mathcal{G}\): The factor of the \((\sigma, x_4)\) bulk equation of
  the author's metric with the author's original weights,
  \(\mathcal{G} = 5\mathsf{A}^2(\mathsf{A} - \mathsf{U})^2 - 6\mathsf{A}^2 + 6\mathsf{A}\mathsf{U} + \mathsf{U}^2 + 1\).
  Unit: dimensionless. Status: DERIVED.
\item
  \(\mathrm{sgn}\): The sign function: \(\mathrm{sgn}(y) = +1\) for
  \(y > 0\), \(-1\) for \(y < 0\). Status: STANDARD.
\item
  \(\sigma\): The proper coordinate of the hidden direction:
  \(\mathrm{d}\sigma = \lvert\cot z\rvert\,\mathrm{d}x_0\),
  \(\sigma = \mathrm{sgn}(\pi/2 - z)\ln(\sin z)/(6H)\) on one period
  \(0 < z < \pi\); \(\sigma = 0\) is the fold \(z = \pi/2\),
  \(\sigma \to \mp\infty\) the tips \(z \to 0, \pi\); one period of
  \(x_0\) is the whole line \(-\infty < \sigma < \infty\). Unit: length.
  Status: DERIVED.
\item
  \(\mathsf{V}\): The rate of change of the scale factor
  \(\mathrm{e}^{-a_4}b_4\) of the extra times on the fold in units of
  \(H\), \(\mathsf{V} = \mathsf{U} - \mathsf{A}\); zero with frozen
  extra times. Unit: dimensionless.
\item
  \(\mathsf{A}\): The expansion rate of ordinary space in units of
  \(H\), \(\mathsf{A} = a_4'/H\). Unit: dimensionless. Value: 0 at the
  static block; \(H_4^{(0)}/H\) today (\(H_4^{(0)}\) the Hubble
  constant); 1 at the author's vacuum, which is not a solution with the
  corrected weights.
\end{itemize}

written with the rate \(\mathsf{A}\) of ordinary space and the rate
\(\mathsf{V}\) of the scale factor \(\mathrm{e}^{-a_4}b_4\) of the extra
times; both polynomials are symmetric under the exchange of
\(\mathsf{A}\) and \(\mathsf{V}\). A bulk without any source needs
\(\mathsf{P} = 0\) and \(\mathsf{U}\mathcal{G} = 0\) at every instant.
For \(\mathsf{U} = 0\),
\(\mathsf{P} = (\mathsf{A}^2 - 1)^2(5\mathsf{A}^2 + 1)\) vanishes only
at \(\mathsf{A} = \pm1\); for \(\mathsf{U} \ne 0\) both polynomials
would need a common root, but their resultant has no real root except
\(\mathsf{U} = 0\). \textbf{With the original weights, a bulk without
any source was therefore at every instant the author's vacuum
\((\mathsf{A}, \mathsf{U}) = (1, 0)\) or its time reverse \((-1, 0)\)},
whatever the time dependence of the rates. A late universe that expands
slowly with constant couplings needed a bulk source. With frozen extra
times (\(\mathsf{V} = 0\)) \(\mathsf{P} = 1 - 3\mathsf{A}^2\) for any
\(a_4(x_4)\), and the energy constraint became a Friedmann equation of
the bulk,
\(3a_4'^2 = H^2 + 16\kappa\rho_{\mathrm{bulk}} = 16\kappa(\rho_{\mathrm{bulk}} - \rho_{\mathrm{s}})\):
the energy density of the stabiliser was the energy density of the bulk
at which ordinary space stands still. Status: DERIVED (ST3, ST8).

\hypertarget{the-correction-of-the-weights-removes-the-stabiliser}{%
\section{The correction of the weights removes the
stabiliser}\label{the-correction-of-the-weights-removes-the-stabiliser}}

\hypertarget{gravitons-must-travel-at-the-speed-of-light}{%
\subsection{Gravitons must travel at the speed of
light}\label{gravitons-must-travel-at-the-speed-of-light}}

A gravitational wave is a small ripple of the metric that travels. The
simplest one, of the ``+'' polarisation, multiplies the metric
components along \(x_1\) and \(x_2\) by \(\mathrm{e}^{h_{\mathrm{gw}}}\)
and \(\mathrm{e}^{-h_{\mathrm{gw}}}\) and travels along \(x_3\). Put
into the field equations of the static bulk and kept to first order in
\(h_{\mathrm{gw}}\), the Lovelock terms give it the wave equation

\[
\begin{aligned}
&\frac{c_{\mathrm{t}}}{2}\,\frac{\partial^2h_{\mathrm{gw}}}{\partial x_4^2}
 - \frac{c_{\mathrm{x}}}{2}\Big(\mathrm{e}^{2H\lvert\sigma\rvert}\,\frac{\partial^2h_{\mathrm{gw}}}{\partial x_3^2}
 + \frac{\partial^2h_{\mathrm{gw}}}{\partial\sigma^2} - 6H\,\mathrm{sgn}(\sigma)\,\frac{\partial h_{\mathrm{gw}}}{\partial\sigma}\Big) = 0,\\
&c_{\mathrm{t}} = 1 - 40w_{(2)} + 360w_{(3)},\qquad c_{\mathrm{x}} = 1 - 24w_{(2)} + 72w_{(3)}
\end{aligned} \tag{FE-68}
\] \textbf{Symbols in (FE-68).}

\begin{itemize}
\tightlist
\item
  \(c_{\mathrm{t}}\): The coefficient of the second time derivative of a
  transverse traceless gravitational wave in the bulk at the static
  block, in units of its value with Einstein gravity alone:
  \(c_{\mathrm{t}} = 1 - 40w_{(2)} + 360w_{(3)}\) (FE-68). Unit:
  dimensionless. Value: 0 with the corrected weights; \(1/16\) with the
  author's original weights. Status: DERIVED.
\item
  \(\partial\): A partial derivative: \(\partial_M\) with respect to the
  coordinate \(X^M\), \(\partial_\mu\) with respect to \(x^\mu\),
  \(\partial_\sigma\) with respect to \(\sigma\),
  \(\partial X/\partial Y\) the derivative of \(X\) with respect to
  \(Y\) at fixed other variables. Status: STANDARD.
\item
  \(h_{\mathrm{gw}}\): The amplitude of a gravitational wave of the +
  polarisation travelling along \(x_3\): the metric components along
  \(x_1\) and \(x_2\) multiplied by \(\mathrm{e}^{h_{\mathrm{gw}}}\) and
  \(\mathrm{e}^{-h_{\mathrm{gw}}}\); a function of \(\sigma\), \(x_3\)
  and \(x_4\) in the bulk, and on the fold the gravitational wave of the
  four-dimensional world. Unit: dimensionless. Status: STANDARD.
\item
  \(x_4\): The time; a prime on a function of \(x_4\) alone is
  \(\mathrm{d}/\mathrm{d}x_4\). Unit: length (= time, \(c = 1\)).
  Status: AUTHOR.
\item
  \(c_{\mathrm{x}}\): The coefficient of the second derivatives along
  ordinary space and along \(\sigma\) of the same wave, in the same
  units: \(c_{\mathrm{x}} = 1 - 24w_{(2)} + 72w_{(3)}\) (FE-68); bulk
  gravitons move with the local speed
  \(\sqrt{c_{\mathrm{x}}/c_{\mathrm{t}}}\). Unit: dimensionless. Value:
  0 with the corrected weights; \(5/16\) with the author's original
  weights. Status: DERIVED.
\item
  \(\mathrm{e}\): Euler's number, the base of the natural logarithm;
  \(\mathrm{e}^{X}\) is the exponential of \(X\). Unit: dimensionless.
  Value: 2.718281828. Status: STANDARD.
\item
  \(H\): The author's constant (an inverse length) that fixes the period
  \(\pi/(6H)\) of \(x_0\), the Lovelock couplings
  \(\alpha_{(p)} = w_{(p)}H^{2-2p}\) and the cosmological constant
  \(\Lambda = -7H^2\) (the author's original value \(-10H^2\) was
  corrected on 2026-10-25); in UM-3 a free parameter. Unit: GeV
  (1/length). Status: AUTHOR.
\item
  \(\sigma\): The proper coordinate of the hidden direction:
  \(\mathrm{d}\sigma = \lvert\cot z\rvert\,\mathrm{d}x_0\),
  \(\sigma = \mathrm{sgn}(\pi/2 - z)\ln(\sin z)/(6H)\) on one period
  \(0 < z < \pi\); \(\sigma = 0\) is the fold \(z = \pi/2\),
  \(\sigma \to \mp\infty\) the tips \(z \to 0, \pi\); one period of
  \(x_0\) is the whole line \(-\infty < \sigma < \infty\). Unit: length.
  Status: DERIVED.
\item
  \(x_3\): The third coordinate of ordinary 3-space. Unit: length.
  Status: AUTHOR.
\item
  \(\mathrm{sgn}\): The sign function: \(\mathrm{sgn}(y) = +1\) for
  \(y > 0\), \(-1\) for \(y < 0\). Status: STANDARD.
\item
  \(w_{(p)}\): The Lovelock weights, \(w_{(1)} = 1\),
  \(w_{(2)} = 1/20\), \(w_{(3)} = 1/360\): the only weights for which
  gravitational waves travel at the speed of light in the bulk and on
  the fold (both coefficients of the bulk wave,
  \(1 - 40w_{(2)} + 360w_{(3)}\) and \(1 - 24w_{(2)} + 72w_{(3)}\),
  vanish). The author's original weights \(w_{(2)} = 1/32\),
  \(w_{(3)} = 1/1152\) (Dirac\_Einstein-Lovelock\_4+4\_Nash.nb,
  In{[}134{]}) made bulk gravitons faster than light; they were
  corrected at the author's order of 2026-10-25. Unit: dimensionless.
  Value: \((1, 1/20, 1/360)\); the author's original
  \((1, 1/32, 1/1152)\). Status: AUTHOR (the form), DERIVED (the
  values).
\end{itemize}

with a time-kinetic coefficient \(c_{\mathrm{t}}\) and a gradient
coefficient \(c_{\mathrm{x}}\) (\(\partial\) marks a partial
derivative). In the bulk the wave moves with the local speed
\(\sqrt{c_{\mathrm{x}}/c_{\mathrm{t}}}\) times the speed of light.
Einstein gravity alone gives \(c_{\mathrm{t}} = c_{\mathrm{x}} = 1\),
the speed of light; the Gauss-Bonnet and cubic terms change the two
coefficients differently, because the static block is curved along
\(\sigma\), ordinary space and the extra times but not along the time.
With the author's original weights \(c_{\mathrm{t}} = 1/16\) and
\(c_{\mathrm{x}} = 5/16\): gravitational waves in the bulk moved
\(\sqrt5 = 2.24\) times faster than light. At the fold the wave has a
thin part of its own, and a wave of long wavelength collects both:

\[
\begin{aligned}
&\int_{0^-}^{0^+}\mathrm{d}\sigma\,(\dots) = \frac{16\big(w_{(2)} - 6w_{(3)}\big)}{H}\,\frac{\partial^2h_{\mathrm{gw}}}{\partial x_4^2}
 - \frac{12\big(w_{(2)} - 2w_{(3)}\big)}{H}\,\frac{\partial^2h_{\mathrm{gw}}}{\partial x_3^2}
 - \frac{c_{\mathrm{x}}}{2}\Big(\frac{\partial h_{\mathrm{gw}}}{\partial\sigma}\Big\vert_{0^+} - \frac{\partial h_{\mathrm{gw}}}{\partial\sigma}\Big\vert_{0^-}\Big)
 = \kappa\,\Pi_{\mathrm{gw}},\\
&\kappa K_{44} = \frac{c_{\mathrm{t}}}{6H} + \frac{16\big(w_{(2)} - 6w_{(3)}\big)}{H},\qquad
 \kappa K_{33} = \frac{c_{\mathrm{x}}}{4H} + \frac{12\big(w_{(2)} - 2w_{(3)}\big)}{H},\qquad
 \kappa\big(K_{44} - K_{33}\big) = -\frac{c_{\mathrm{t}}}{12H},\qquad c_{\mathrm{g}}^2 = \frac{K_{33}}{K_{44}}
\end{aligned} \tag{FE-69}
\] \textbf{Symbols in (FE-69).}

\begin{itemize}
\tightlist
\item
  \(\mathrm{d}\): The differential: \(\mathrm{d}X\) is an infinitesimal
  change of \(X\), \(\mathrm{d}X/\mathrm{d}Y\) the derivative of \(X\)
  with respect to \(Y\), \(\mathrm{d}^nX\) the volume element of \(n\)
  variables. Status: STANDARD.
\item
  \(\sigma\): The proper coordinate of the hidden direction:
  \(\mathrm{d}\sigma = \lvert\cot z\rvert\,\mathrm{d}x_0\),
  \(\sigma = \mathrm{sgn}(\pi/2 - z)\ln(\sin z)/(6H)\) on one period
  \(0 < z < \pi\); \(\sigma = 0\) is the fold \(z = \pi/2\),
  \(\sigma \to \mp\infty\) the tips \(z \to 0, \pi\); one period of
  \(x_0\) is the whole line \(-\infty < \sigma < \infty\). Unit: length.
  Status: DERIVED.
\item
  \(w_{(p)}\): The Lovelock weights, \(w_{(1)} = 1\),
  \(w_{(2)} = 1/20\), \(w_{(3)} = 1/360\): the only weights for which
  gravitational waves travel at the speed of light in the bulk and on
  the fold (both coefficients of the bulk wave,
  \(1 - 40w_{(2)} + 360w_{(3)}\) and \(1 - 24w_{(2)} + 72w_{(3)}\),
  vanish). The author's original weights \(w_{(2)} = 1/32\),
  \(w_{(3)} = 1/1152\) (Dirac\_Einstein-Lovelock\_4+4\_Nash.nb,
  In{[}134{]}) made bulk gravitons faster than light; they were
  corrected at the author's order of 2026-10-25. Unit: dimensionless.
  Value: \((1, 1/20, 1/360)\); the author's original
  \((1, 1/32, 1/1152)\). Status: AUTHOR (the form), DERIVED (the
  values).
\item
  \(H\): The author's constant (an inverse length) that fixes the period
  \(\pi/(6H)\) of \(x_0\), the Lovelock couplings
  \(\alpha_{(p)} = w_{(p)}H^{2-2p}\) and the cosmological constant
  \(\Lambda = -7H^2\) (the author's original value \(-10H^2\) was
  corrected on 2026-10-25); in UM-3 a free parameter. Unit: GeV
  (1/length). Status: AUTHOR.
\item
  \(\partial\): A partial derivative: \(\partial_M\) with respect to the
  coordinate \(X^M\), \(\partial_\mu\) with respect to \(x^\mu\),
  \(\partial_\sigma\) with respect to \(\sigma\),
  \(\partial X/\partial Y\) the derivative of \(X\) with respect to
  \(Y\) at fixed other variables. Status: STANDARD.
\item
  \(h_{\mathrm{gw}}\): The amplitude of a gravitational wave of the +
  polarisation travelling along \(x_3\): the metric components along
  \(x_1\) and \(x_2\) multiplied by \(\mathrm{e}^{h_{\mathrm{gw}}}\) and
  \(\mathrm{e}^{-h_{\mathrm{gw}}}\); a function of \(\sigma\), \(x_3\)
  and \(x_4\) in the bulk, and on the fold the gravitational wave of the
  four-dimensional world. Unit: dimensionless. Status: STANDARD.
\item
  \(x_4\): The time; a prime on a function of \(x_4\) alone is
  \(\mathrm{d}/\mathrm{d}x_4\). Unit: length (= time, \(c = 1\)).
  Status: AUTHOR.
\item
  \(x_3\): The third coordinate of ordinary 3-space. Unit: length.
  Status: AUTHOR.
\item
  \(c_{\mathrm{x}}\): The coefficient of the second derivatives along
  ordinary space and along \(\sigma\) of the same wave, in the same
  units: \(c_{\mathrm{x}} = 1 - 24w_{(2)} + 72w_{(3)}\) (FE-68); bulk
  gravitons move with the local speed
  \(\sqrt{c_{\mathrm{x}}/c_{\mathrm{t}}}\). Unit: dimensionless. Value:
  0 with the corrected weights; \(5/16\) with the author's original
  weights. Status: DERIVED.
\item
  \(\kappa\): The eight-dimensional gravitational coupling,
  \(\kappa = 8\pi G_8\). Unit: GeV\(^{-6}\). Status: AUTHOR.
\item
  \(\Pi_{\mathrm{gw}}\): The transverse traceless stress of the matter
  of the fold that drives the gravitational wave \(h_{\mathrm{gw}}\)
  (half the difference of its components along \(x_1\) and \(x_2\)),
  integrated across the fold (eight-dimensional normalisation). Unit:
  GeV\(^7\). Status: STANDARD.
\item
  \(K_{44}\): The coefficient of
  \(\partial^2h_{\mathrm{gw}}/\partial x_4^2\) in the equation of the
  long-wavelength gravitational wave (the same at every \(\sigma\)) at
  the static block, integrated across the hidden direction with the
  volume factor \(\mathrm{e}^{-6H\lvert\sigma\rvert}\): the part of the
  bulk plus the thin part at the fold. Unit: GeV\(^5\)
  (\(\kappa K_{44}\) is a length). Status: DERIVED.
\item
  \(c_{\mathrm{t}}\): The coefficient of the second time derivative of a
  transverse traceless gravitational wave in the bulk at the static
  block, in units of its value with Einstein gravity alone:
  \(c_{\mathrm{t}} = 1 - 40w_{(2)} + 360w_{(3)}\) (FE-68). Unit:
  dimensionless. Value: 0 with the corrected weights; \(1/16\) with the
  author's original weights. Status: DERIVED.
\item
  \(K_{33}\): The coefficient of
  \(-\partial^2h_{\mathrm{gw}}/\partial x_3^2\) in the same equation,
  with the same two parts; the wave travels at
  \(c_{\mathrm{g}}^2 = K_{33}/K_{44}\). Unit: GeV\(^5\). Status:
  DERIVED.
\item
  \(c_{\mathrm{g}}\): The speed of gravitational waves on the fold, in
  units of the speed of light. Unit: dimensionless. Value: 1 at every
  wavelength with the corrected weights; with the author's original
  weights the long-wavelength mode had \(c_{\mathrm{g}}^2 = 83/82\) and
  shorter waves were slower than light. Status: DERIVED.
\end{itemize}

For every pair of weights the difference of the two coefficients of the
long wave is \(-c_{\mathrm{t}}/(12\kappa H)\): it travels at the speed
of light only if \(c_{\mathrm{t}} = 0\). If \(c_{\mathrm{t}}\) were zero
but not \(c_{\mathrm{x}}\), waves in the bulk would have no time-kinetic
term at all, an infinite local speed. Gravitons travel at the speed of
light everywhere only if both coefficients vanish:
\(1 - 40w_{(2)} + 360w_{(3)} = 0\) and
\(1 - 24w_{(2)} + 72w_{(3)} = 0\), two linear equations with the single
solution \(w_{(2)} = 1/20\), \(w_{(3)} = 1/360\). With the original
weights the long wave had \(c_{\mathrm{g}}^2 = 83/82\), faster than
light (the bulk part alone \(15/2\): \(c_{\mathrm{g}} = 2.74\), a figure
that an earlier version of UM-3 quoted), and the wave bound to the fold
was slower than light at short wavelengths. With the corrected weights
the bulk part vanishes, \(\kappa K_{44} = \kappa K_{33} = 8/(15H)\), and
\(c_{\mathrm{g}} = 1\) at every wavelength. Status: DERIVED
(\texttt{um3\_weights.py}, W1 to W3, W7; \texttt{um3\_bulk.py}, B5).

\hypertarget{the-static-block-with-the-corrected-weights}{%
\subsection{The static block with the corrected
weights}\label{the-static-block-with-the-corrected-weights}}

With \(w_{(2)} = 1/20\) and \(w_{(3)} = 1/360\) the warped and the time
entries of the static block are both \(7H^2\) before the cosmological
constant (the last column of the table of SB-8; FE-70), so the
cosmological constant \(\Lambda = -7H^2\) makes the static block a
vacuum:

\[
\begin{aligned}
&\kappa\,\mathcal{T}_{\mathrm{stab}}{}^a{}_b = \Big(\sum_{p=1}^{3}\alpha_{(p)}\,\mathcal{E}_{(p)}{}^a{}_b + \Lambda\,\delta^a{}_b\Big)\Big\vert_{\text{static block}} = 0:\\
&\qquad (15 - 9 + 1 - 7)H^2 = 0\ \text{(warped directions)},\qquad (21 - 21 + 7 - 7)H^2 = 0\ \text{(the time)}
\end{aligned} \tag{IN-7}
\] \textbf{Symbols in (IN-7).}

\begin{itemize}
\tightlist
\item
  \(\kappa\): The eight-dimensional gravitational coupling,
  \(\kappa = 8\pi G_8\). Unit: GeV\(^{-6}\). Status: AUTHOR.
\item
  \(\mathcal{T}_{\mathrm{stab}}{}^a{}_b\): The energy-momentum tensor of
  the stabiliser of the bulk (frame components): the source that a
  static bulk of the author's metric required with the author's original
  Lovelock weights, found and named by Claude (Anthropic) on 2026-10-24;
  with the corrected weights and \(\Lambda = -7H^2\) it is zero (IN-7).
  Unit: GeV\(^8\). Value: 0 with the corrected weights. Status: DERIVED.
\item
  \(a\): A frame (orthonormal) index running over \(0, 1, \dots, 7\)
  (frame direction \((0)\) along \(\sigma\), \((1), (2), (3)\) along
  ordinary space, \((4)\) along the time, \((5), (6), (7)\) along the
  extra times); raised and lowered with the frame metric
  \(\eta_{ab} = \mathrm{diag}(1, 1, 1, 1, -1, -1, -1, -1)\); on the fold
  (\(a, b \in \{1, \dots, 7\}\)) the frame of the fold; a repeated index
  is summed. Status: STANDARD.
\item
  \(b\): A frame index (as \(a\)). Status: STANDARD.
\item
  \(p\): The Lovelock order \(p = 1, 2, 3\) (Einstein, Gauss-Bonnet,
  cubic), as an index and as a number; in the Kohn-Sham equations a
  second label of a Kohn-Sham orbital (as \(o\), with \(\delta_{op}\)
  the Kronecker delta of the two labels). Status: STANDARD.
\item
  \(\alpha_{(p)}\): The Lovelock coupling of order \(p\),
  \(\alpha_{(p)} = w_{(p)}H^{2-2p}\) (\(\alpha_{(1)} = 1\)). Unit:
  GeV\(^{2-2p}\). Status: AUTHOR.
\item
  \(\mathcal{E}_{(p)}{}^M{}_N\): The Lovelock tensor of order \(p\): the
  variation of \(\sqrt{\lvert g\rvert}\,\mathcal{L}_{(p)}\) with respect
  to the metric (\(\mathcal{E}_{(1)}\) the Einstein tensor,
  \(\mathcal{E}_{(2)}\) the Gauss-Bonnet tensor); with frame indices its
  frame components. Unit: 1/length\(^{2p}\). Status: STANDARD.
\item
  \(\Lambda\): The eight-dimensional cosmological constant,
  \(\Lambda = -7H^2\): with the corrected weights the value for which
  the static block is a vacuum (its Lovelock tensors add to \(7H^2\) in
  every direction). The author's original value was \(\Lambda = -10H^2\)
  (he writes the pure number \(\Lambda_{\mathrm{N}} = 10\) in
  \(0 = -\Lambda_{\mathrm{N}} + \sum_p H^{-2p}w_{(p)}\mathcal{E}_{(p)}\)).
  Unit: GeV\(^2\). Value: \(-7H^2\); the author's original \(-10H^2\).
  Status: AUTHOR (the term), DERIVED (the value).
\item
  \(\delta\): The Kronecker delta: with two indices \(\delta^M_N = 1\)
  if \(M = N\) and 0 otherwise (likewise \(\delta_{X\mathrm{L}}\) is 1
  if the label \(X\) of a half of dirac16complex is L and 0 otherwise);
  with upper and lower index strings the generalised Kronecker delta
  \(\delta^{M_1\cdots M_n}_{N_1\cdots N_n}\), the determinant of the
  \(n \times n\) matrix of the Kronecker deltas \(\delta^{M_a}_{N_b}\)
  (\(+1\) or \(-1\) if the upper indices are an even or odd permutation
  of the lower ones, 0 otherwise). Status: STANDARD.
\item
  \(H\): The author's constant (an inverse length) that fixes the period
  \(\pi/(6H)\) of \(x_0\), the Lovelock couplings
  \(\alpha_{(p)} = w_{(p)}H^{2-2p}\) and the cosmological constant
  \(\Lambda = -7H^2\) (the author's original value \(-10H^2\) was
  corrected on 2026-10-25); in UM-3 a free parameter. Unit: GeV
  (1/length). Status: AUTHOR.
\end{itemize}

The stabiliser of SB-9 is zero: order by order, Einstein \((15, 21)\),
Gauss-Bonnet \((-9, -21)\), cubic \((1, 7)\) and the cosmological
constant \((-7, -7)\) add to zero along the warped directions and along
the time. The source of the fold of the static block becomes a pure
tension, the same in all seven directions of the fold:

\[
\mathsf{A} = \mathsf{U} = \mathsf{A}_1 = \mathsf{U}_1 = 0:\qquad \kappa\,\mathcal{T}_{\mathrm{bulk}}{}^a{}_b = 0,\qquad
\kappa\,\mathcal{S}^{(i)}{}_{(i)} = \kappa\,\mathcal{S}^{(4)}{}_{(4)} = \kappa\,\mathcal{S}^{(r)}{}_{(r)} = -\frac{32}{5}H \tag{FE-32}
\] \textbf{Symbols in (FE-32).}

\begin{itemize}
\tightlist
\item
  \(\mathsf{A}\): The expansion rate of ordinary space in units of
  \(H\), \(\mathsf{A} = a_4'/H\). Unit: dimensionless. Value: 0 at the
  static block; \(H_4^{(0)}/H\) today (\(H_4^{(0)}\) the Hubble
  constant); 1 at the author's vacuum, which is not a solution with the
  corrected weights.
\item
  \(\mathsf{U}\): The rate of change of \(b_4\) in units of \(H\),
  \(\mathsf{U} = b_4'/(Hb_4)\). Unit: dimensionless. Value:
  \(\mathsf{U} = \mathsf{A}\) with frozen extra times.
\item
  \(\mathsf{A}_1\): The time derivative of \(\mathsf{A}\) in units of
  \(H\), \(\mathsf{A}_1 = \mathsf{A}'/H = a_4''/H^2\). Unit:
  dimensionless.
\item
  \(\mathsf{U}_1\): The time derivative of \(\mathsf{U}\) in units of
  \(H\), \(\mathsf{U}_1 = \mathsf{U}'/H\). Unit: dimensionless. Status:
  DERIVED.
\item
  \(\kappa\): The eight-dimensional gravitational coupling,
  \(\kappa = 8\pi G_8\). Unit: GeV\(^{-6}\). Status: AUTHOR.
\item
  \(\mathcal{T}_{\mathrm{bulk}}\): The energy-momentum tensor of all
  sources of the bulk (with the corrected weights dirac16complex alone,
  \(\mathcal{T}_\Psi\); with the author's original weights also the
  stabiliser); \(\mathcal{T}_{\mathrm{bulk}}{}^{(0)}{}_b\) is its flux
  along \(\sigma\) (frame components). Unit: GeV\(^8\). Status:
  STANDARD.
\item
  \(a\): A frame (orthonormal) index running over \(0, 1, \dots, 7\)
  (frame direction \((0)\) along \(\sigma\), \((1), (2), (3)\) along
  ordinary space, \((4)\) along the time, \((5), (6), (7)\) along the
  extra times); raised and lowered with the frame metric
  \(\eta_{ab} = \mathrm{diag}(1, 1, 1, 1, -1, -1, -1, -1)\); on the fold
  (\(a, b \in \{1, \dots, 7\}\)) the frame of the fold; a repeated index
  is summed. Status: STANDARD.
\item
  \(b\): A frame index (as \(a\)). Status: STANDARD.
\item
  \(\mathcal{S}^a{}_b\): The thin source of the fold: the
  energy-momentum tensor concentrated on the fold, integrated across it,
  \(\kappa\mathcal{S}^a{}_b = \int_{0^-}^{0^+}\mathrm{d}\sigma\sum_p\alpha_{(p)}\mathcal{E}_{(p)}{}^a{}_b\)
  (FD-11, the junction conditions); \(\mathcal{S}^4{}_4\) is minus its
  energy density, \(\mathcal{S}^1{}_1\) and \(\mathcal{S}^5{}_5\) its
  pressures along ordinary space and along the extra times (frame
  components along the fold, eight-dimensional normalisation: an energy
  per unit volume of the six other directions of the fold). Unit:
  GeV\(^7\). Value:
  \(\kappa\mathcal{S}^a{}_b = -\frac{32}{5}H\delta^a{}_b\) at the static
  block. Status: DERIVED.
\item
  \(i\): A generation index (\(i = 1, 2, 3\)), which also labels the
  neutrino mass states \(\nu_1, \nu_2, \nu_3\); in a frame component
  such as \(\mathcal{S}^{(i)}{}_{(i)}\) a direction \(x_i\) of ordinary
  space (the three components are equal); a repeated index is not summed
  unless a sum sign is written. Status: STANDARD.
\item
  \(r\): An extra-time direction of the fold: \(r \in \{5, 6, 7\}\)
  (\(x_5, x_6, x_7\)); not summed (the three components are equal).
  Status: STANDARD.
\item
  \(H\): The author's constant (an inverse length) that fixes the period
  \(\pi/(6H)\) of \(x_0\), the Lovelock couplings
  \(\alpha_{(p)} = w_{(p)}H^{2-2p}\) and the cosmological constant
  \(\Lambda = -7H^2\) (the author's original value \(-10H^2\) was
  corrected on 2026-10-25); in UM-3 a free parameter. Unit: GeV
  (1/length). Status: AUTHOR.
\end{itemize}

This is the arrangement of Randall and Sundrum (FD-8): an empty bulk
with a negative cosmological constant and a fold of pure tension, like
the drum skin of example 5. Status: DERIVED (ST9;
\texttt{um3\_weights.py}, W4; \texttt{um3\_fold.py}, J5).

\hypertarget{the-expanding-fold-with-the-corrected-weights}{%
\subsection{The expanding fold with the corrected
weights}\label{the-expanding-fold-with-the-corrected-weights}}

The corrected equations have no need of a stabiliser in any era. Their
two constraints are

\[
\begin{aligned}
&\kappa\,\mathcal{T}_{\mathrm{bulk}}{}^{(4)}{}_{(4)} = \mathsf{A}^3(\mathsf{A} - \mathsf{U})^3H^2,\qquad
 \kappa\,\mathcal{T}_{\mathrm{bulk}}{}^{(0)}{}_{(4)} = 3\mathsf{A}^2\mathsf{U}(\mathsf{A} - \mathsf{U})^2H^2\,\mathrm{sgn}(\sigma),\\
&\text{no bulk source:}\quad \mathsf{A} = 0\ \ \text{or}\ \ \mathsf{U} = \mathsf{A}\quad\text{at every instant}
\end{aligned} \tag{FE-66}
\] \textbf{Symbols in (FE-66).}

\begin{itemize}
\tightlist
\item
  \(\kappa\): The eight-dimensional gravitational coupling,
  \(\kappa = 8\pi G_8\). Unit: GeV\(^{-6}\). Status: AUTHOR.
\item
  \(\mathcal{T}_{\mathrm{bulk}}\): The energy-momentum tensor of all
  sources of the bulk (with the corrected weights dirac16complex alone,
  \(\mathcal{T}_\Psi\); with the author's original weights also the
  stabiliser); \(\mathcal{T}_{\mathrm{bulk}}{}^{(0)}{}_b\) is its flux
  along \(\sigma\) (frame components). Unit: GeV\(^8\). Status:
  STANDARD.
\item
  \(\mathsf{A}\): The expansion rate of ordinary space in units of
  \(H\), \(\mathsf{A} = a_4'/H\). Unit: dimensionless. Value: 0 at the
  static block; \(H_4^{(0)}/H\) today (\(H_4^{(0)}\) the Hubble
  constant); 1 at the author's vacuum, which is not a solution with the
  corrected weights.
\item
  \(\mathsf{U}\): The rate of change of \(b_4\) in units of \(H\),
  \(\mathsf{U} = b_4'/(Hb_4)\). Unit: dimensionless. Value:
  \(\mathsf{U} = \mathsf{A}\) with frozen extra times.
\item
  \(H\): The author's constant (an inverse length) that fixes the period
  \(\pi/(6H)\) of \(x_0\), the Lovelock couplings
  \(\alpha_{(p)} = w_{(p)}H^{2-2p}\) and the cosmological constant
  \(\Lambda = -7H^2\) (the author's original value \(-10H^2\) was
  corrected on 2026-10-25); in UM-3 a free parameter. Unit: GeV
  (1/length). Status: AUTHOR.
\item
  \(\mathrm{sgn}\): The sign function: \(\mathrm{sgn}(y) = +1\) for
  \(y > 0\), \(-1\) for \(y < 0\). Status: STANDARD.
\item
  \(\sigma\): The proper coordinate of the hidden direction:
  \(\mathrm{d}\sigma = \lvert\cot z\rvert\,\mathrm{d}x_0\),
  \(\sigma = \mathrm{sgn}(\pi/2 - z)\ln(\sin z)/(6H)\) on one period
  \(0 < z < \pi\); \(\sigma = 0\) is the fold \(z = \pi/2\),
  \(\sigma \to \mp\infty\) the tips \(z \to 0, \pi\); one period of
  \(x_0\) is the whole line \(-\infty < \sigma < \infty\). Unit: length.
  Status: DERIVED.
\end{itemize}

so a bulk without any source has, at every instant, either static
ordinary space (\(\mathsf{A} = 0\)) or frozen extra times
(\(\mathsf{U} = \mathsf{A}\)), and on both branches every component of
the bulk equations vanishes. The expanding fold of our universe is on
the second branch, where the bulk is an exact vacuum for every expansion
history and the fold obeys exactly the Friedmann equations of
four-dimensional general relativity (the equations of the expansion of a
homogeneous universe, \(3a_4'^2 = \kappa_7\rho\) for the energy density
\(\rho\) and \(-(2a_4'' + 3a_4'^2) = \kappa_7p\) for the pressure \(p\)
of its matter above the tension, both in the normalisation of the thin
source):

\[
\begin{aligned}
&\mathsf{U} = \mathsf{A},\quad \mathsf{U}_1 = \mathsf{A}_1\ \ \big(\mathrm{e}^{-a_4}b_4 = \text{constant},\ \mathcal{Z} = 1\big):\qquad
 \kappa\,\mathcal{T}_{\mathrm{bulk}}{}^a{}_b = 0\ \ \text{for any } a_4(x_4),\\
&\frac{\kappa\,\mathcal{S}^{(i)}{}_{(i)}}{H} = -\frac{16}{15}\big(6 + 3\mathsf{A}^2 + 2\mathsf{A}_1\big),\qquad
 \frac{\kappa\,\mathcal{S}^{(4)}{}_{(4)}}{H} = -\frac{16}{5}\big(2 + \mathsf{A}^2\big),\qquad
 \frac{\kappa\,\mathcal{S}^{(r)}{}_{(r)}}{H} = -\frac45\big(8 + 8\mathsf{A}^2 + \mathsf{A}^4 + 4\mathsf{A}_1 + \mathsf{A}^2\mathsf{A}_1\big),\\
&3a_4'^2 = -\kappa_7\Big(\mathcal{S}^{(4)}{}_{(4)} + \frac{32H}{5\kappa}\Big),\qquad
 -\big(2a_4'' + 3a_4'^2\big) = \kappa_7\Big(\mathcal{S}^{(i)}{}_{(i)} + \frac{32H}{5\kappa}\Big),\qquad \kappa_7 = \frac{15}{16}\,\kappa H
\end{aligned} \tag{FE-33}
\] \textbf{Symbols in (FE-33).}

\begin{itemize}
\tightlist
\item
  \(\mathsf{U}\): The rate of change of \(b_4\) in units of \(H\),
  \(\mathsf{U} = b_4'/(Hb_4)\). Unit: dimensionless. Value:
  \(\mathsf{U} = \mathsf{A}\) with frozen extra times.
\item
  \(\mathsf{A}\): The expansion rate of ordinary space in units of
  \(H\), \(\mathsf{A} = a_4'/H\). Unit: dimensionless. Value: 0 at the
  static block; \(H_4^{(0)}/H\) today (\(H_4^{(0)}\) the Hubble
  constant); 1 at the author's vacuum, which is not a solution with the
  corrected weights.
\item
  \(\mathsf{U}_1\): The time derivative of \(\mathsf{U}\) in units of
  \(H\), \(\mathsf{U}_1 = \mathsf{U}'/H\). Unit: dimensionless. Status:
  DERIVED.
\item
  \(\mathsf{A}_1\): The time derivative of \(\mathsf{A}\) in units of
  \(H\), \(\mathsf{A}_1 = \mathsf{A}'/H = a_4''/H^2\). Unit:
  dimensionless.
\item
  \(\mathrm{e}\): Euler's number, the base of the natural logarithm;
  \(\mathrm{e}^{X}\) is the exponential of \(X\). Unit: dimensionless.
  Value: 2.718281828. Status: STANDARD.
\item
  \(a_4\): The author's metric function of the time \(x_4\) that scales
  ordinary 3-space: on the fold \(\mathrm{e}^{a_4}\) is the scale factor
  of cosmology and \(a_4'\) the Hubble rate. Unit: dimensionless.
  Status: AUTHOR.
\item
  \(b_4\): The author's metric function of the time \(x_4\) in the
  extra-time components, \(b_4 > 0\), never identically 1 (the author's
  order of 2026-10-24); the scale factor (the factor of the lengths) of
  the extra times is \(\mathrm{e}^{-a_4}b_4\) on the fold and
  \(\mathrm{e}^{-a_4}b_4\,\mathrm{e}^{-H\lvert\sigma\rvert}\) at the
  distance \(\sigma\) from it (the square root of the warp factor times
  \(\mathrm{e}^{-a_4}b_4\)). Unit: dimensionless. Status: AUTHOR.
\item
  \(\mathcal{Z}\): The fold factor,
  \(\mathcal{Z}(x_4) = V_{\mathrm{t}}(\mathrm{e}^{-a_4}b_4)^3/V_*\): the
  factor by which the size of the extra times multiplies every term of
  the fold; \(\mathcal{Z} = 1\) when the extra times are frozen at their
  late value. Unit: dimensionless. Value: 1 in the late universe.
  Status: DERIVED.
\item
  \(\kappa\): The eight-dimensional gravitational coupling,
  \(\kappa = 8\pi G_8\). Unit: GeV\(^{-6}\). Status: AUTHOR.
\item
  \(\mathcal{T}_{\mathrm{bulk}}\): The energy-momentum tensor of all
  sources of the bulk (with the corrected weights dirac16complex alone,
  \(\mathcal{T}_\Psi\); with the author's original weights also the
  stabiliser); \(\mathcal{T}_{\mathrm{bulk}}{}^{(0)}{}_b\) is its flux
  along \(\sigma\) (frame components). Unit: GeV\(^8\). Status:
  STANDARD.
\item
  \(a\): A frame (orthonormal) index running over \(0, 1, \dots, 7\)
  (frame direction \((0)\) along \(\sigma\), \((1), (2), (3)\) along
  ordinary space, \((4)\) along the time, \((5), (6), (7)\) along the
  extra times); raised and lowered with the frame metric
  \(\eta_{ab} = \mathrm{diag}(1, 1, 1, 1, -1, -1, -1, -1)\); on the fold
  (\(a, b \in \{1, \dots, 7\}\)) the frame of the fold; a repeated index
  is summed. Status: STANDARD.
\item
  \(b\): A frame index (as \(a\)). Status: STANDARD.
\item
  \(x_4\): The time; a prime on a function of \(x_4\) alone is
  \(\mathrm{d}/\mathrm{d}x_4\). Unit: length (= time, \(c = 1\)).
  Status: AUTHOR.
\item
  \(\mathcal{S}^a{}_b\): The thin source of the fold: the
  energy-momentum tensor concentrated on the fold, integrated across it,
  \(\kappa\mathcal{S}^a{}_b = \int_{0^-}^{0^+}\mathrm{d}\sigma\sum_p\alpha_{(p)}\mathcal{E}_{(p)}{}^a{}_b\)
  (FD-11, the junction conditions); \(\mathcal{S}^4{}_4\) is minus its
  energy density, \(\mathcal{S}^1{}_1\) and \(\mathcal{S}^5{}_5\) its
  pressures along ordinary space and along the extra times (frame
  components along the fold, eight-dimensional normalisation: an energy
  per unit volume of the six other directions of the fold). Unit:
  GeV\(^7\). Value:
  \(\kappa\mathcal{S}^a{}_b = -\frac{32}{5}H\delta^a{}_b\) at the static
  block. Status: DERIVED.
\item
  \(i\): A generation index (\(i = 1, 2, 3\)), which also labels the
  neutrino mass states \(\nu_1, \nu_2, \nu_3\); in a frame component
  such as \(\mathcal{S}^{(i)}{}_{(i)}\) a direction \(x_i\) of ordinary
  space (the three components are equal); a repeated index is not summed
  unless a sum sign is written. Status: STANDARD.
\item
  \(r\): An extra-time direction of the fold: \(r \in \{5, 6, 7\}\)
  (\(x_5, x_6, x_7\)); not summed (the three components are equal).
  Status: STANDARD.
\item
  \(H\): The author's constant (an inverse length) that fixes the period
  \(\pi/(6H)\) of \(x_0\), the Lovelock couplings
  \(\alpha_{(p)} = w_{(p)}H^{2-2p}\) and the cosmological constant
  \(\Lambda = -7H^2\) (the author's original value \(-10H^2\) was
  corrected on 2026-10-25); in UM-3 a free parameter. Unit: GeV
  (1/length). Status: AUTHOR.
\item
  \({}'\): The derivative of a function of the time \(x_4\) alone with
  respect to \(x_4\) (\(a_4' = \mathrm{d}a_4/\mathrm{d}x_4\); two primes
  the second derivative); on a function of \(\sigma\) alone (a profile)
  the derivative with respect to \(\sigma\). Status: STANDARD.
\item
  \(\kappa_7\): The gravitational coupling of the fold: with the
  corrected Lovelock weights the thin source of the fold equals, to
  second order in the rates of expansion, the Einstein tensor of the
  seven-dimensional metric of the fold divided by
  \(\kappa_7 = 15\kappa H/16\), and exactly so in the Friedmann
  equations of a fold with frozen extra times; four-dimensional gravity
  on the fold has \(8\pi G_{\mathrm{N}} = \kappa_7/V_*\) (plus the
  induced term, if the fields of the fold induce one). Unit:
  GeV\(^{-5}\). Value: \(15\kappa H/16\). Status: DERIVED.
\end{itemize}

There is no flow of energy along the hidden direction, so the matter of
the fold is conserved by itself (FE-11 with a vanishing right-hand
side). The author's vacuum, in which the original weights needed no
source, now needs one:

\[
\begin{aligned}
&a_4 = Hx_4,\quad b_4 = \text{constant}\ (\mathsf{A} = 1,\ \mathsf{U} = \mathsf{A}_1 = \mathsf{U}_1 = 0):\qquad
 \frac{\kappa\,\mathcal{T}_{\mathrm{bulk}}{}^a{}_b}{H^2} = \mathrm{diag}\Big(-\frac75, -\frac75, -\frac75, -\frac75, 1, -\frac75, -\frac75, -\frac75\Big),\\
&\kappa\rho_{\mathrm{bulk}} = -H^2,\qquad \kappa\,\mathcal{S}^{(i)}{}_{(i)} = \kappa\,\mathcal{S}^{(r)}{}_{(r)} = -\frac{42}{5}H,\qquad
 \kappa\,\mathcal{S}^{(4)}{}_{(4)} = -\frac{28}{5}H
\end{aligned} \tag{FE-31}
\] \textbf{Symbols in (FE-31).}

\begin{itemize}
\tightlist
\item
  \(a_4\): The author's metric function of the time \(x_4\) that scales
  ordinary 3-space: on the fold \(\mathrm{e}^{a_4}\) is the scale factor
  of cosmology and \(a_4'\) the Hubble rate. Unit: dimensionless.
  Status: AUTHOR.
\item
  \(H\): The author's constant (an inverse length) that fixes the period
  \(\pi/(6H)\) of \(x_0\), the Lovelock couplings
  \(\alpha_{(p)} = w_{(p)}H^{2-2p}\) and the cosmological constant
  \(\Lambda = -7H^2\) (the author's original value \(-10H^2\) was
  corrected on 2026-10-25); in UM-3 a free parameter. Unit: GeV
  (1/length). Status: AUTHOR.
\item
  \(x_4\): The time; a prime on a function of \(x_4\) alone is
  \(\mathrm{d}/\mathrm{d}x_4\). Unit: length (= time, \(c = 1\)).
  Status: AUTHOR.
\item
  \(b_4\): The author's metric function of the time \(x_4\) in the
  extra-time components, \(b_4 > 0\), never identically 1 (the author's
  order of 2026-10-24); the scale factor (the factor of the lengths) of
  the extra times is \(\mathrm{e}^{-a_4}b_4\) on the fold and
  \(\mathrm{e}^{-a_4}b_4\,\mathrm{e}^{-H\lvert\sigma\rvert}\) at the
  distance \(\sigma\) from it (the square root of the warp factor times
  \(\mathrm{e}^{-a_4}b_4\)). Unit: dimensionless. Status: AUTHOR.
\item
  \(\mathsf{A}\): The expansion rate of ordinary space in units of
  \(H\), \(\mathsf{A} = a_4'/H\). Unit: dimensionless. Value: 0 at the
  static block; \(H_4^{(0)}/H\) today (\(H_4^{(0)}\) the Hubble
  constant); 1 at the author's vacuum, which is not a solution with the
  corrected weights.
\item
  \(\mathsf{U}\): The rate of change of \(b_4\) in units of \(H\),
  \(\mathsf{U} = b_4'/(Hb_4)\). Unit: dimensionless. Value:
  \(\mathsf{U} = \mathsf{A}\) with frozen extra times.
\item
  \(\mathsf{A}_1\): The time derivative of \(\mathsf{A}\) in units of
  \(H\), \(\mathsf{A}_1 = \mathsf{A}'/H = a_4''/H^2\). Unit:
  dimensionless.
\item
  \(\mathsf{U}_1\): The time derivative of \(\mathsf{U}\) in units of
  \(H\), \(\mathsf{U}_1 = \mathsf{U}'/H\). Unit: dimensionless. Status:
  DERIVED.
\item
  \(\kappa\): The eight-dimensional gravitational coupling,
  \(\kappa = 8\pi G_8\). Unit: GeV\(^{-6}\). Status: AUTHOR.
\item
  \(\mathcal{T}_{\mathrm{bulk}}\): The energy-momentum tensor of all
  sources of the bulk (with the corrected weights dirac16complex alone,
  \(\mathcal{T}_\Psi\); with the author's original weights also the
  stabiliser); \(\mathcal{T}_{\mathrm{bulk}}{}^{(0)}{}_b\) is its flux
  along \(\sigma\) (frame components). Unit: GeV\(^8\). Status:
  STANDARD.
\item
  \(a\): A frame (orthonormal) index running over \(0, 1, \dots, 7\)
  (frame direction \((0)\) along \(\sigma\), \((1), (2), (3)\) along
  ordinary space, \((4)\) along the time, \((5), (6), (7)\) along the
  extra times); raised and lowered with the frame metric
  \(\eta_{ab} = \mathrm{diag}(1, 1, 1, 1, -1, -1, -1, -1)\); on the fold
  (\(a, b \in \{1, \dots, 7\}\)) the frame of the fold; a repeated index
  is summed. Status: STANDARD.
\item
  \(b\): A frame index (as \(a\)). Status: STANDARD.
\item
  \(\mathrm{diag}\): A diagonal matrix with the listed entries on its
  diagonal (in the order of the coordinates \(x_0, x_1, \dots, x_7\) for
  a tensor of eight dimensions, \(x_1, \dots, x_7\) for a tensor of the
  fold, \(x_1, \dots, x_4\) for a four-dimensional tensor). Status:
  STANDARD.
\item
  \(\rho_{\mathrm{bulk}}\): The energy density of the bulk sources,
  \(\rho_{\mathrm{bulk}} = -\mathcal{T}_{\mathrm{bulk}}{}^{(4)}{}_{(4)}\).
  Unit: GeV\(^8\). Status: DERIVED.
\item
  \(\mathcal{S}^a{}_b\): The thin source of the fold: the
  energy-momentum tensor concentrated on the fold, integrated across it,
  \(\kappa\mathcal{S}^a{}_b = \int_{0^-}^{0^+}\mathrm{d}\sigma\sum_p\alpha_{(p)}\mathcal{E}_{(p)}{}^a{}_b\)
  (FD-11, the junction conditions); \(\mathcal{S}^4{}_4\) is minus its
  energy density, \(\mathcal{S}^1{}_1\) and \(\mathcal{S}^5{}_5\) its
  pressures along ordinary space and along the extra times (frame
  components along the fold, eight-dimensional normalisation: an energy
  per unit volume of the six other directions of the fold). Unit:
  GeV\(^7\). Value:
  \(\kappa\mathcal{S}^a{}_b = -\frac{32}{5}H\delta^a{}_b\) at the static
  block. Status: DERIVED.
\item
  \(i\): A generation index (\(i = 1, 2, 3\)), which also labels the
  neutrino mass states \(\nu_1, \nu_2, \nu_3\); in a frame component
  such as \(\mathcal{S}^{(i)}{}_{(i)}\) a direction \(x_i\) of ordinary
  space (the three components are equal); a repeated index is not summed
  unless a sum sign is written. Status: STANDARD.
\item
  \(r\): An extra-time direction of the fold: \(r \in \{5, 6, 7\}\)
  (\(x_5, x_6, x_7\)); not summed (the three components are equal).
  Status: STANDARD.
\end{itemize}

so the primordial era of UM-3 is no longer the author's vacuum but an
inflation of the fold on the branch of frozen extra times, driven by the
vacuum energy of a field of the fold, as in general relativity. Status:
DERIVED (ST9; \texttt{um3\_weights.py}, W5 and W6;
\texttt{um3\_geometry.py}, G4).

\hypertarget{what-takes-the-place-of-the-stabiliser}{%
\subsection{What takes the place of the
stabiliser}\label{what-takes-the-place-of-the-stabiliser}}

Nothing in the bulk: the corrected bulk is a vacuum on both branches.
But one requirement remains, on the fold. The component of the junction
along the extra times (the third line of FE-33) asks the matter of the
fold for a definite pressure along the extra times, which radiation has
and dust and vacuum energy do not. Without it the volume of the extra
times would change, and with it Newton's constant and every coupling of
the fold: since the time of nucleosynthesis Newton's constant would have
fallen by the factor 15.8, and today it would change 6800 times faster
than lunar laser ranging allows. UM-3 therefore assumes that the fold
fluid, the thin medium of the fold, supplies the missing pressure
\(q_{\mathrm{hold}}\):

\[
\begin{aligned}
&p_{\mathrm{ff}} = -\rho_{\mathrm{ff}},\qquad q_{\mathrm{ff}} = -\rho_{\mathrm{ff}} + q_{\mathrm{hold}},\qquad
 \lambda_{\mathrm{F}} = \frac{32HV_*}{5\kappa},\\
&q_{\mathrm{hold}} = -\Delta_{\mathrm{r}} - \frac45\big(\mathsf{A}^4 + \mathsf{A}^2\mathsf{A}_1\big)\frac{HV_*}{\kappa},\qquad
 \Delta_{\mathrm{r}} = q_{\mathrm{SM}} + q_\chi + q_{\mathrm{P}} - \big(\rho_{\mathrm{ff}} - \lambda_{\mathrm{F}}\big) + \frac{\rho - 3\mathcal{P}}{2}
\end{aligned} \tag{FE-65}
\] \textbf{Symbols in (FE-65).}

\begin{itemize}
\tightlist
\item
  \(p_{\mathrm{ff}}\): The pressure of the fold fluid along ordinary
  space. Unit: GeV\(^4\). Value:
  \(-\lambda_{\mathrm{F}} = -32HV_*/(5\kappa)\). Status: ASSUMED (the
  fluid), DERIVED (the value).
\item
  \(\rho_{\mathrm{ff}}\): The energy density of the fold fluid. Unit:
  GeV\(^4\). Value: \(\lambda_{\mathrm{F}} = 32HV_*/(5\kappa)\) (its
  tension). Status: ASSUMED (the fluid), DERIVED (the value).
\item
  \(q_{\mathrm{ff}}\): The pressure of the fold fluid along the extra
  times: minus its tension plus the pressure \(q_{\mathrm{hold}}\) that
  holds the volume of the extra times fixed. Unit: GeV\(^4\). Value:
  \(-\lambda_{\mathrm{F}} + q_{\mathrm{hold}}\); \(-32HV_*/(5\kappa)\)
  at the static block. Status: ASSUMED (the fluid), DERIVED (the value).
\item
  \(q_{\mathrm{hold}}\): The pressure along the extra times that the
  fold fluid adds to its tension to hold the volume of the extra times
  fixed (four-dimensional normalisation). Unit: GeV\(^4\). Value:
  \(-\Delta_{\mathrm{r}} = -0.793\,\rho_{\mathrm{c}}^{(0)}\) today.
  Status: DERIVED (the value), ASSUMED (that the fold fluid supplies
  it).
\item
  \(\lambda_{\mathrm{F}}\): The tension of the fold: the part of the
  fold fluid whose pressures along ordinary space and along the extra
  times are minus its energy density (a vacuum energy of the
  seven-dimensional fold), per unit volume of ordinary space of the fold
  (four-dimensional normalisation). Unit: GeV\(^4\). Value:
  \(32HV_*/(5\kappa)\). Status: DERIVED.
\item
  \(H\): The author's constant (an inverse length) that fixes the period
  \(\pi/(6H)\) of \(x_0\), the Lovelock couplings
  \(\alpha_{(p)} = w_{(p)}H^{2-2p}\) and the cosmological constant
  \(\Lambda = -7H^2\) (the author's original value \(-10H^2\) was
  corrected on 2026-10-25); in UM-3 a free parameter. Unit: GeV
  (1/length). Status: AUTHOR.
\item
  \(V_*\): The physical volume of the extra times \(x_5, x_6, x_7\) on
  the fold in the late universe,
  \(V_* = V_{\mathrm{t}}(\mathrm{e}^{-a_4}b_4)^3\) (with
  \(V_{\mathrm{t}} = \int\mathrm{d}x_5\,\mathrm{d}x_6\,\mathrm{d}x_7\)
  the coordinate volume of the extra times) at a time when
  \(\mathrm{e}^{-a_4}b_4\) is constant; it converts the
  eight-dimensional normalisation into the four-dimensional one. Unit:
  length\(^3\) (GeV\(^{-3}\)). Status: ASSUMED.
\item
  \(\kappa\): The eight-dimensional gravitational coupling,
  \(\kappa = 8\pi G_8\). Unit: GeV\(^{-6}\). Status: AUTHOR.
\item
  \(\Delta_{\mathrm{r}}\): The mismatch of the radion condition: the sum
  over the matter of the fold of its pressure along the extra times and
  \((\rho - 3\mathcal{P})/2\) (zero for radiation). Unit: GeV\(^4\).
  Value: \(0.793\,\rho_{\mathrm{c}}^{(0)}\) today. Status: DERIVED.
\item
  \(\mathsf{A}\): The expansion rate of ordinary space in units of
  \(H\), \(\mathsf{A} = a_4'/H\). Unit: dimensionless. Value: 0 at the
  static block; \(H_4^{(0)}/H\) today (\(H_4^{(0)}\) the Hubble
  constant); 1 at the author's vacuum, which is not a solution with the
  corrected weights.
\item
  \(\mathsf{A}_1\): The time derivative of \(\mathsf{A}\) in units of
  \(H\), \(\mathsf{A}_1 = \mathsf{A}'/H = a_4''/H^2\). Unit:
  dimensionless.
\item
  \(q_{\mathrm{SM}}\): The pressure of the Standard Model along the
  extra times at the fold: the mean on-shell Lagrangian density of the
  Standard-Model fields (the extra-time component of the energy-momentum
  tensor of every term of the fold is its Lagrangian, because the fields
  of the fold do not depend on the extra times). Unit: GeV\(^4\).
  Status: DERIVED.
\item
  \(q_\chi\): The pressure of the bound states of dirac16complex along
  the extra times, integrated across the fold (their on-shell
  Lagrangian, the self-interaction energy per volume; zero for a dilute
  gas). Unit: GeV\(^4\). Status: DERIVED.
\item
  \(q_{\mathrm{P}}\): The pressure of the portals along the extra times:
  their mean on-shell Lagrangian density. Unit: GeV\(^4\). Status:
  DERIVED.
\item
  \(\rho\): The total energy density of the matter of the fold (all
  species); \(\Delta\rho\) its perturbation. Unit: GeV\(^4\). Status:
  STANDARD.
\item
  \(\mathcal{P}\): The total pressure of the matter of the fold;
  \(\Delta\mathcal{P}\) its perturbation. Unit: GeV\(^4\). Status:
  STANDARD.
\end{itemize}

The equation of the volume of the extra times (the radion) shows what
\(q_{\mathrm{hold}}\) does:

\[
H^2\big(\mathsf{V}_1 + 3\mathsf{V}(\mathsf{A} + \mathsf{V})\big) = \frac{8\pi G_{\mathrm{N}}}{5}\big(\rho - 3\mathcal{P} + 2q_{\mathrm{tot}}\big) \tag{FE-72}
\] \textbf{Symbols in (FE-72).}

\begin{itemize}
\tightlist
\item
  \(H\): The author's constant (an inverse length) that fixes the period
  \(\pi/(6H)\) of \(x_0\), the Lovelock couplings
  \(\alpha_{(p)} = w_{(p)}H^{2-2p}\) and the cosmological constant
  \(\Lambda = -7H^2\) (the author's original value \(-10H^2\) was
  corrected on 2026-10-25); in UM-3 a free parameter. Unit: GeV
  (1/length). Status: AUTHOR.
\item
  \(\mathsf{V}_1\): The time derivative of \(\mathsf{V}\) in units of
  \(H\), \(\mathsf{V}_1 = \mathsf{V}'/H\). Unit: dimensionless.
\item
  \(\mathsf{V}\): The rate of change of the scale factor
  \(\mathrm{e}^{-a_4}b_4\) of the extra times on the fold in units of
  \(H\), \(\mathsf{V} = \mathsf{U} - \mathsf{A}\); zero with frozen
  extra times. Unit: dimensionless.
\item
  \(\mathsf{A}\): The expansion rate of ordinary space in units of
  \(H\), \(\mathsf{A} = a_4'/H\). Unit: dimensionless. Value: 0 at the
  static block; \(H_4^{(0)}/H\) today (\(H_4^{(0)}\) the Hubble
  constant); 1 at the author's vacuum, which is not a solution with the
  corrected weights.
\item
  \(\pi\): The number pi, the ratio of the circumference of a circle to
  its diameter. Unit: dimensionless. Value: 3.141592654. Status:
  STANDARD.
\item
  \(G_{\mathrm{N}}\): Newton's constant of four-dimensional gravity (the
  gravity of our world, on the fold),
  \(G_{\mathrm{N}} = 1/M_{\mathrm{Pl}}^2\) with the Planck mass
  \(M_{\mathrm{Pl}} = 1.221\times10^{19}\) GeV. Unit: GeV\(^{-2}\).
  Value: \(6.708\times10^{-39}\) GeV\(^{-2}\). Status: STANDARD.
\item
  \(\rho\): The total energy density of the matter of the fold (all
  species); \(\Delta\rho\) its perturbation. Unit: GeV\(^4\). Status:
  STANDARD.
\item
  \(\mathcal{P}\): The total pressure of the matter of the fold;
  \(\Delta\mathcal{P}\) its perturbation. Unit: GeV\(^4\). Status:
  STANDARD.
\item
  \(q_{\mathrm{tot}}\): The total pressure along the extra times of the
  matter of the fold (the Standard Model, the bound states of
  dirac16complex, the portals and the fold fluid without its tension).
  Unit: GeV\(^4\). Status: STANDARD.
\end{itemize}

The rate \(\mathsf{V}\) of the scale factor of the extra times stays
zero only if \(\rho - 3\mathcal{P} + 2q_{\mathrm{tot}} = 0\), and
\(q_{\mathrm{hold}}\) makes this so. It acts as a potential that holds
the size of the extra times, a role that the ``block stabiliser'' of the
first universe model of this repository played for its own metric; what
produces it is OPEN. Status: DERIVED (the requirement and its value,
\texttt{um3\_series.py}, SR4 to SR6), ASSUMED (that the fold fluid
supplies it), OPEN (its origin).

\hypertarget{the-five-points-verified}{%
\section{The five points, verified}\label{the-five-points-verified}}

\textbf{{[}1{]} The stabiliser couples only to gravity.} By its
definition it had no term in the action of UM-3 other than its own: it
entered the field equations only through its energy-momentum tensor in
FE-1, and it exerted no force on dirac16complex or on the Standard Model
except through the curvature that it caused. On the static bulk this was
consistent: the stabiliser was conserved by itself (ST4). In an
expanding universe the bulk source that the original equations required
released energy (SB-15) that flowed into the fold (SB-14) and was taken
up there (SB-16). What the equations fixed was the total bulk source and
its flow; whether this meant an exchange of energy between the
stabiliser and the matter of the fold, or a second medium of the bulk,
or a thin layer of the stabiliser itself, they did not say (ASSUMED, an
interpretation). \textbf{Verdict, original weights:} exact on the static
bulk; in an expanding universe a medium that couples only to gravity and
keeps \(p = 5\rho\) was not enough to supply the required source (SB-13
against SB-14). \textbf{Verdict, corrected weights:} there is no
stabiliser, and the bulk of an expanding fold with frozen extra times is
a vacuum without any flow (FE-33). Status: DERIVED (ST4, ST6, ST9).

\textbf{{[}2{]} It is absent in the primordial era, whose bulk is the
author's vacuum (FE-31).} With the original weights the author's vacuum,
\(a_4 = Hx_4\) with constant \(b_4\) (\(\mathsf{A} = 1\),
\(\mathsf{U} = 0\)), solved the field equations of the bulk with nothing
on the right (ST2), and with the stabiliser present it would not have
been a solution: the stabiliser was absent there. But a stabiliser that
couples only to gravity and keeps \(p = 5\rho\) cannot appear later out
of nothing (SB-13). Points {[}1{]} and {[}2{]} together required a
mechanism that created it between the primordial era and the late
universe: an exchange of energy with the fold (which breaks {[}1{]}), a
conversion from or into dirac16complex in the bulk (the bulk source is
the sum of both), or a change of its equation of state (which a single
medium with \(p = 5\rho\) does not have); UM-3 contained none.
\textbf{Verdict, original weights:} {[}2{]} true, with the creation of
the stabiliser OPEN. \textbf{Verdict, corrected weights:} the author's
vacuum is not a solution (FE-31); the primordial era is an inflation of
the fold with frozen extra times, whose bulk is a vacuum; no era of UM-3
has a stabiliser, so the question of its creation no longer arises.
Status: DERIVED (ST2, ST6, ST9; \texttt{um3\_primordial.py}, PR0).

\textbf{{[}3{]} Its existence and value are DERIVED from the requirement
of a static bulk.} With the original weights the value was DERIVED: SB-9
fixed every entry (ST1). The existence of a bulk source was DERIVED more
generally than for an exactly static bulk: without any bulk source the
bulk was at every instant the author's vacuum or its time reverse
(SB-17, ST3, ST8), so a slowly expanding bulk needed one. What was
ASSUMED is that the source is a separate medium, with dirac16complex
contributing nothing; its microscopic nature was OPEN, and among
homogeneous scalar fields only a ghost reproduced it (SB-12, ST5).
\textbf{Verdict:} value DERIVED; the existence of a bulk source DERIVED;
its identification with a separate medium ASSUMED. With the corrected
weights the same derivation gives the value zero (IN-7, ST9): the
existence of the stabiliser was a consequence of the original weights.

\textbf{{[}4{]} The static bulk} is defined in SB-5 and illustrated by
five examples and a comparison of static and empty spacetimes in the
part ``The static bulk''.

\textbf{{[}5{]} Warp and warp factor} are defined in SB-2 and
illustrated by four examples in the part ``Warp and warp factor''.

\hypertarget{the-importance-of-the-stabiliser-for-um-3}{%
\section{The importance of the stabiliser for
UM-3}\label{the-importance-of-the-stabiliser-for-um-3}}

\textbf{With the original weights, no late universe like ours without
it.} A bulk without any source was at every instant the author's vacuum
or its time reverse (SB-17): ordinary space had to grow or shrink at the
full rate \(H\). The only alternative to the stabiliser would have been
to make the author's vacuum itself the late universe, with \(H\) equal
to today's expansion rate \(H_4^{(0)}\). But in the author's vacuum the
extra times shrink, as \(\mathrm{e}^{-a_4}b_4 = \mathrm{e}^{-Hx_4}b_4\),
and every coupling of the fold is divided by the volume of the extra
times, so every squared coupling of the Standard Model grows as
\(\mathrm{e}^{3Hx_4}\). The natural nuclear reactor of Oklo (Gabon),
which ran about 1.8 billion years ago, bounds the change of the
fine-structure constant \(\alpha_{\mathrm{em}}\) since then to about one
part in ten million
(\(-0.9\times10^{-7} < \Delta\alpha_{\mathrm{em}}/\alpha_{\mathrm{em}} < 1.2\times10^{-7}\)).
The author's vacuum would change it by

\[
\Delta\ln\alpha_{\mathrm{em}} = 3H_4^{(0)}\,\Delta x_4 = 0.372 \tag{SB-18}
\] \textbf{Symbols in (SB-18).}

\begin{itemize}
\tightlist
\item
  \(\Delta\): A change, difference or small variation of the quantity
  that follows (for example an interval of the time, a perturbation of a
  density, a difference of two squared masses, or an addition to an
  effective number of species). Status: STANDARD.
\item
  \(\alpha_{\mathrm{em}}\): The fine-structure constant. Unit:
  dimensionless. Value: 1/137.036. Status: STANDARD.
\item
  \(H_4\): The expansion rate (Hubble rate) of the fold, \(H_4 = a_4'\),
  as a function of the time or of the redshift; \(H_4^{(0)}\) its value
  today, the Hubble constant. Unit: GeV (1/time). Value:
  \(H_4^{(0)} = 100\,h_{100}\) km/s/Mpc with \(h_{100} = 0.674\) (Planck
  2018): 67.4 km/s/Mpc \(= 1.44\times10^{-42}\) GeV. Status: STANDARD.
\item
  \(x_4\): The time; a prime on a function of \(x_4\) alone is
  \(\mathrm{d}/\mathrm{d}x_4\). Unit: length (= time, \(c = 1\)).
  Status: AUTHOR.
\end{itemize}

(\(H_4^{(0)} = 67.4\) km/s/Mpc, \(\Delta x_4 = 1.8\) Gyr; \(\Delta\)
marks a change): 3.1 million times the bound. With the stabiliser and
frozen extra times the couplings stayed constant. Status: DERIVED (ST7),
STANDARD (the bound of Oklo).

\textbf{With the original weights, the zero point of the expansion.}
With frozen extra times the energy constraint was a Friedmann equation
of the bulk,
\(3a_4'^2 = 16\kappa(\rho_{\mathrm{bulk}} - \rho_{\mathrm{s}})\)
(SB-17): the energy density of the stabiliser was the energy density of
the bulk at which ordinary space stands still, and the expansion
measured the excess of the bulk energy density over it. Integrated over
the hidden direction and the extra times, per unit volume of ordinary
space on the fold, the stabiliser carried the energy

\[
\rho_{4,\mathrm{s}} = \frac{\rho_{\mathrm{s}}V_*}{3H} = -\frac{HV_*}{48\kappa} \tag{SB-19}
\] \textbf{Symbols in (SB-19).}

\begin{itemize}
\tightlist
\item
  \(\rho_{4,\mathrm{s}}\): The energy of the stabiliser per unit volume
  of ordinary space of the fold with the author's original weights: its
  energy density integrated over the hidden direction and the extra
  times. Unit: GeV\(^4\). Value: \(-HV_*/(48\kappa)\) with the author's
  original weights; 0 with the corrected weights. Status: DERIVED.
\item
  \(\rho_{\mathrm{s}}\): The energy density of the stabiliser of the
  bulk with the author's original weights; with frozen extra times it
  was the energy density of the bulk at which ordinary space stands
  still. Unit: GeV\(^8\). Value: \(\kappa\rho_{\mathrm{s}} = -H^2/16\)
  with the author's original weights; 0 with the corrected weights.
  Status: DERIVED.
\item
  \(V_*\): The physical volume of the extra times \(x_5, x_6, x_7\) on
  the fold in the late universe,
  \(V_* = V_{\mathrm{t}}(\mathrm{e}^{-a_4}b_4)^3\) (with
  \(V_{\mathrm{t}} = \int\mathrm{d}x_5\,\mathrm{d}x_6\,\mathrm{d}x_7\)
  the coordinate volume of the extra times) at a time when
  \(\mathrm{e}^{-a_4}b_4\) is constant; it converts the
  eight-dimensional normalisation into the four-dimensional one. Unit:
  length\(^3\) (GeV\(^{-3}\)). Status: ASSUMED.
\item
  \(H\): The author's constant (an inverse length) that fixes the period
  \(\pi/(6H)\) of \(x_0\), the Lovelock couplings
  \(\alpha_{(p)} = w_{(p)}H^{2-2p}\) and the cosmological constant
  \(\Lambda = -7H^2\) (the author's original value \(-10H^2\) was
  corrected on 2026-10-25); in UM-3 a free parameter. Unit: GeV
  (1/length). Status: AUTHOR.
\item
  \(\kappa\): The eight-dimensional gravitational coupling,
  \(\kappa = 8\pi G_8\). Unit: GeV\(^{-6}\). Status: AUTHOR.
\end{itemize}

(the volume factor \(\mathrm{e}^{-6H\lvert\sigma\rvert}\) of the warped
directions, integrated over both sides of the fold, gives \(1/(3H)\),
and \(V_*\) is the physical volume of the extra times on the fold in the
late universe). It was not a vacuum energy that makes the fold expand:
the static block as a whole did not expand at all. Status: DERIVED (ST7,
ST8).

\textbf{With the original weights, the most exotic ingredient.} Its
violation of the null energy condition (SB-11) put it among the sources
that come with instabilities or with signals faster than light;
Einstein's static universe, the example most like it, is unstable. If
the stabiliser had been a ghost field, its quanta would carry negative
energy, and even through gravity alone a ghost lets the vacuum decay
into pairs of ghosts and pairs of ordinary particles, at a rate that is
finite only if the momenta of the ghosts are cut off (Cline, Jeon and
Moore, 2004). Whether the static block with the stabiliser was stable
was never analysed. Status: STANDARD (the general results).

\textbf{With the corrected weights, its absence changes UM-3.} The
weights for which gravitons travel at the speed of light make the static
bulk a vacuum, and the stabiliser disappears with them. Four
consequences:

\begin{enumerate}
\def\labelenumi{\arabic{enumi}.}
\tightlist
\item
  Gravitational waves travel at the speed of light at every wavelength,
  and the gravitational waves of the neutron-star merger GW170817, which
  arrived together with its light, are explained for every value of
  \(H\) (FE-69).
\item
  The fold obeys the Friedmann equations of four-dimensional general
  relativity exactly, with Newton's constant
  \(8\pi G_{\mathrm{N}} = \kappa_7/V_*\), \(\kappa_7 = 15\kappa H/16\)
  (FE-33): the Gauss-Bonnet and cubic terms give the fold an
  Einstein-Hilbert term of its own.
\item
  The bulk needs no exotic medium in any era: no violation of the null
  energy condition, no ghost, and no question of how the stabiliser
  would be created after the primordial era.
\item
  Two new tasks appear. The fold fluid must hold the volume of the extra
  times fixed with the pressure \(q_{\mathrm{hold}}\) (FE-65, FE-72),
  whose origin is OPEN; and the static block is a degenerate vacuum (the
  double and triple roots of FE-70), around which the bulk equations
  have no linear part, so the response of the bulk to sources that
  depend on \(\sigma\) is strongly coupled and not computed (OPEN).
\end{enumerate}

\textbf{Summary.} The stabiliser was the price of a slowly expanding
late universe with constant couplings in the author's metric with the
author's original weights and cosmological constant. Its value was fixed
exactly and introduced no new parameter. It was not a property of the
author's metric but of those weights: the weights that make gravitons
travel at the speed of light make the static bulk empty, and the
stabiliser vanishes with them.

\hypertarget{provenance}{%
\section{Provenance}\label{provenance}}

\textbf{What comes from the author (AUTHOR).} The coordinates
\(X = \{x_0, \dots, x_7\}\) and their roles
(\texttt{Dirac\_Einstein-Lovelock\_4+4\_Nash.nb}, section DEFCHART); the
metric with \(b_4[x_4]^2\) and its two functions \(a_4(x_4)\) and
\(b_4(x_4)\) (the author's order of 2026-10-24); the Einstein-Lovelock
action with its original weights \(w_{(1)} = 1\), \(w_{(2)} = 1/32\),
\(w_{(3)} = 1/1152\) and the cosmological constant \(\Lambda = -10H^2\)
(\texttt{Dirac\_Einstein-Lovelock\_4+4\_Nash.nb}, In{[}134{]}, where the
author writes the pure number 10); the author's vacuum \(a_4 = Hx_4\)
with constant \(b_4\); the field dirac16complex; the order of
2026-10-25, 05:09 PDT, to document the stabiliser with proper
attribution and to verify the five points; and the order of 2026-10-25,
07:55 PDT: ``Gravitons must travel at c.~If these incorrect claims come
from my choice of the Lovelock couplings and my choice of the
cosmological constant \(\Lambda\) THEN my choices are wrong, and you
MUST correct them.''

\textbf{What UM-3 adds.} The proper distance \(\sigma\) (DERIVED from
the metric); the static block as the first approximation of the late
universe (ASSUMED); the stabiliser of the original weights (found and
named by Claude on 2026-10-24; its value DERIVED, the medium ASSUMED);
the corrected weights \(w_{(2)} = 1/20\), \(w_{(3)} = 1/360\) and
\(\Lambda = -7H^2\) (DERIVED at the order of 07:55); the fold fluid with
its tension and its holding pressure (ASSUMED, with DERIVED values).

\textbf{The lineage.} The first universe model of this repository (the
folder \texttt{Revision/universe\_model}, version 1.1, present at the
first commit of the repository, \texttt{0fdf375} of 2026-10-19) assumed
a ``block stabiliser'': a potential for the size
\(\mathrm{e}^{-a_4}b_4^{1/2}\) of the block of the hidden direction and
the extra times of its metric (whose extra-time components carried
\(b_4\) instead of \(b_4^2\)), with an origin left open; that version
had itself changed the third weight to \(1/768\). The second universe
model (\texttt{universe\_model\_2}, version 3 of 2026-10-18) assumed a
block stabiliser whose energy density falls as the inverse eighth power
of the block scale,
\(\kappa\rho_{\mathrm{stab}} = (26 + 4\sqrt{51})H^2 = 54.566H^2\) at its
static point. The stabiliser of UM-3 was a new object, for the author's
new metric with \(b_4^2\), with a negative energy density and
\(p = 5\rho\).

\textbf{Who found what, and when.} On 2026-10-24 Claude computed, in the
gravity module of UM-3 (\texttt{um3\_geometry.py}, check G4, commit
\texttt{4448938} of 2026-10-24 17:54 UTC-7), that a static bulk of the
author's metric needs the energy density \(-H^2/(16\kappa)\) and that
the bulk with frozen extra times needs the energy density
\((3\mathsf{A}^2 - 1)H^2/(16\kappa)\) and a flow of energy along the
hidden direction; the pressures \(-5H^2/(16\kappa)\) were added when the
medium was introduced and named ``the stabiliser of the bulk'', as the
term IN-7 of the action of UM-3, in the commit \texttt{170d8d3} of
2026-10-24 19:05 UTC-7 (the field equations of UM-3, where its existence
was labelled ASSUMED and its value DERIVED). On 2026-10-25, at the
author's order of 05:09, Claude wrote the first version of this document
with its LaTeX source and its PDF and the checks ST1 to ST8 (commit
\texttt{cd5709d} of 2026-10-25 06:55 UTC-7), and an independent referee
read it; its corrections (the energy conditions of a negative vacuum
energy, the use of ``warp'', the constraints of a sourceless bulk, the
rates of the extra times, the derivation of SB-9, the junction
conditions, the lineage and the references) are in this version. After
the author's order of 07:55 Claude derived the weights for which
gravitons travel at the speed of light (\texttt{um3\_weights.py}), found
that with them the static bulk is a vacuum and the stabiliser vanishes
(check ST9), and rewrote this document as the record of the stabiliser
and of its removal. All of this work was done in Claude Code sessions
configured for the model claude-opus-5-5 (session
https://claude.ai/code/session\_01AyG7BZ5rkLxBARRjLHHquB).

\textbf{How it was computed.} The left-hand side of the field equations
was computed exactly with the computer algebra system sympy for the
general diagonal warped metric
\(\mathrm{d}\sigma^2 + \mathrm{e}^{2\mathfrak{a}}(\mathrm{d}x_1^2 + \mathrm{d}x_2^2 + \mathrm{d}x_3^2) - \mathrm{e}^{2\mathfrak{n}}\mathrm{d}x_4^2 - \mathrm{e}^{2\mathfrak{b}}(\mathrm{d}x_5^2 + \mathrm{d}x_6^2 + \mathrm{d}x_7^2)\),
with functions \(\mathfrak{a}\), \(\mathfrak{n}\), \(\mathfrak{b}\) of
\(\sigma\) and \(x_4\): the curvature in an orthonormal frame, and the
Lovelock tensors of orders 1, 2 and 3 contracted with the generalised
Kronecker delta. It was then specialised to the author's metric,
\(\mathfrak{a} = a_4 - H\lvert\sigma\rvert\),
\(\mathfrak{b} = \ln b_4 - a_4 - H\lvert\sigma\rvert\),
\(\mathfrak{n} = 0\). The source of the fold was computed from the jump
of the first \(\sigma\)-derivatives of the metric across the fold (the
junction conditions FD-11), as the difference of a junction potential on
the two sides, and independently by smoothing the kink. The program
\texttt{universe\_model\_3/code/um3\_stabiliser.py} of the repository
computes and checks the following:

\begin{longtable}[]{@{}
  >{\raggedright\arraybackslash}p{(\columnwidth - 4\tabcolsep) * \real{0.3333}}
  >{\raggedright\arraybackslash}p{(\columnwidth - 4\tabcolsep) * \real{0.3333}}
  >{\raggedright\arraybackslash}p{(\columnwidth - 4\tabcolsep) * \real{0.3333}}@{}}
\toprule\noalign{}
\begin{minipage}[b]{\linewidth}\raggedright
check
\end{minipage} & \begin{minipage}[b]{\linewidth}\raggedright
what it verifies
\end{minipage} & \begin{minipage}[b]{\linewidth}\raggedright
result
\end{minipage} \\
\midrule\noalign{}
\endhead
\bottomrule\noalign{}
\endlastfoot
ST1 & original weights: the static bulk requires the source SB-9 on both
sides, without flow; the order-by-order values of SB-8 and its table;
the source of the fold SB-10 order by order & PASS \\
ST2 & original weights: the author's vacuum needs no source; its fold
source \((-10, -6, -10)H\); the expansion around it & PASS \\
ST3 & original weights: without a source the bulk is at every instant
the author's vacuum or its time reverse (the two constraints and their
resultant, SB-17) & PASS \\
ST4 & the stabiliser is conserved by itself on the static bulk & PASS \\
ST5 & from the values of ST1: the weak and the null energy condition
fail (SB-11); among homogeneous scalar fields only a ghost reproduces it
(SB-12) & PASS \\
ST6 & SB-13; the frozen branch SB-14 with its \(\mathsf{A}_1\) terms and
its fold source; the energy budget SB-15; the tension and pressureless
part of the fold source; the balance of the fold SB-16 & PASS \\
ST7 & the drift SB-18 in the author's vacuum, and the share SB-19 &
PASS \\
ST8 & the two constraints contain no second time derivative; their form
in \((\mathsf{A}, \mathsf{V})\); the Friedmann equation of the bulk &
PASS \\
ST9 & corrected weights: the static bulk is a vacuum (IN-7), the fold
source is a pure tension (FE-32), the frozen branch and static ordinary
space are vacua without flow, the constraints FE-66, the author's vacuum
is no longer a solution & PASS \\
\end{longtable}

The other checks used here are in the programs \texttt{um3\_weights.py}
(W1 to W8: the gravitational waves FE-68 and FE-69 for any weights, the
unique weights, FE-70, the vacua and the fold of the corrected weights),
\texttt{um3\_geometry.py} (G2 to G4: the curvature and the bulk
equations of the author's metric), \texttt{um3\_fold.py} (J4 to J7: the
junction conditions and the sources of the fold),
\texttt{um3\_series.py} (SR4 to SR6: the volume of the extra times) and
\texttt{um3\_primordial.py} (PR0). \texttt{um3\_stabiliser.py} writes
\texttt{outputs/stabiliser.json} and
\texttt{reports/stabiliser\_checks.json}, reads the general source of
the fold from the output of \texttt{um3\_fold.py}, and runs in about 45
seconds (\texttt{python\ universe\_model\_3/code/um3\_stabiliser.py}
from the root of the repository).

\textbf{Sources outside UM-3 (STANDARD).} A. Einstein, ``Kosmologische
Betrachtungen zur allgemeinen Relativitätstheorie'', Sitzungsberichte
der Königlich Preußischen Akademie der Wissenschaften (1917) 142 (the
static universe, SB-6); A. S. Eddington, ``On the instability of
Einstein's spherical world'', Monthly Notices of the Royal Astronomical
Society 90 (1930) 668; K. Schwarzschild, ``Über das Gravitationsfeld
eines Massenpunktes nach der Einsteinschen Theorie'', Sitzungsberichte
der Königlich Preußischen Akademie der Wissenschaften (1916) 189 (the
static field outside a mass); D. Lovelock, ``The Einstein tensor and its
generalizations'', Journal of Mathematical Physics 12 (1971) 498; W.
Israel, ``Singular hypersurfaces and thin shells in general
relativity'', Il Nuovo Cimento B 44 (1966) 1 (the junction conditions);
S. C. Davis, ``Generalised Israel junction conditions for a Gauss-Bonnet
brane world'', Physical Review D 67 (2003) 024030; E. Gravanis and S.
Willison, ``Israel conditions for the Gauss-Bonnet theory and the
Friedmann equation on the brane universe'', Physics Letters B 562 (2003)
118; R. V. Pound and G. A. Rebka, ``Apparent weight of photons'',
Physical Review Letters 4 (1960) 337; L. Randall and R. Sundrum, ``A
large mass hierarchy from a small extra dimension'', Physical Review
Letters 83 (1999) 3370, and ``An alternative to compactification'',
Physical Review Letters 83 (1999) 4690 (FD-8); P. Hořava and E. Witten,
``Heterotic and type I string dynamics from eleven dimensions'', Nuclear
Physics B 460 (1996) 506, and ``Eleven-dimensional supergravity on a
manifold with boundary'', Nuclear Physics B 475 (1996) 94 (example 9 of
the definition of the fold); H. Whitney, ``On singularities of mappings
of Euclidean spaces. I. Mappings of the plane into the plane'', Annals
of Mathematics 62 (1955) 374 (the fold map of example 5); S. W. Hawking
and G. F. R. Ellis, ``The Large Scale Structure of Space-Time''
(Cambridge University Press, 1973) (the energy conditions); S. Dubovsky,
T. Grégoire, A. Nicolis and R. Rattazzi, ``Null energy condition and
superluminal propagation'', Journal of High Energy Physics 03 (2006)
025; J. M. Cline, S. Jeon and G. D. Moore, ``The phantom menaced:
constraints on low-energy effective ghosts'', Physical Review D 70
(2004) 043543; the bound of Oklo on the drift of the fine-structure
constant over 1.8 Gyr (A. I. Shlyakhter, ``Direct test of the constancy
of fundamental nuclear constants'', Nature 264 (1976) 340; T. Damour and
F. Dyson, ``The Oklo bound on the time variation of the fine-structure
constant revisited'', Nuclear Physics B 480 (1996) 37); the Hubble
constant \(H_4^{(0)} = 67.4\) km/s/Mpc (Planck Collaboration, ``Planck
2018 results. VI. Cosmological parameters'', Astronomy and Astrophysics
641 (2020) A6); the speed of gravitational waves (B. P. Abbott et al.,
``Gravitational waves and gamma-rays from a binary neutron star merger:
GW170817 and GRB 170817A'', Astrophysical Journal Letters 848 (2017)
L13).

\end{document}
